\documentclass[11pt,a4paper,twoside]{book}

%% ============================================================
%% PACKAGES
%% ============================================================
\usepackage[utf8]{inputenc}
\usepackage[T1]{fontenc}
\usepackage[french]{babel}
\usepackage{lmodern}
\usepackage[a4paper,top=2.5cm,bottom=2.5cm,left=2.5cm,right=2.2cm]{geometry}
\usepackage{amsmath,amssymb,amsthm}
\usepackage{mathtools}
\usepackage{array,booktabs,longtable}
\usepackage{enumitem}
\usepackage[dvipsnames,svgnames]{xcolor}
\usepackage[most]{tcolorbox}
\usepackage{pgfplots}
\pgfplotsset{compat=1.18}
\usepackage{tikz}
\usetikzlibrary{arrows.meta,calc,positioning,shapes}
\usepackage{fancyhdr}
\usepackage[hidelinks]{hyperref}
\usepackage{stmaryrd}
\usepackage{mathrsfs}
\usepackage{caption}
\setlength{\headheight}{14pt}

%% ============================================================
%% COULEURS
%% ============================================================
\definecolor{caplpblue}{HTML}{1E3A5F}
\definecolor{caplpbluebg}{HTML}{EFF6FF}
\definecolor{caplppurple}{HTML}{5B21B6}
\definecolor{caplppurplebg}{HTML}{F5F3FF}
\definecolor{caplpgreen}{HTML}{15803D}
\definecolor{caplpgreenbg}{HTML}{F0FDF4}
\definecolor{caplpred}{HTML}{B91C1C}
\definecolor{caplpredbg}{HTML}{FEF2F2}
\definecolor{caplpgray}{HTML}{374151}
\definecolor{caplpgraybg}{HTML}{F9FAFB}
\definecolor{caplporange}{HTML}{C2410C}
\definecolor{caplporangebg}{HTML}{FFF7ED}

%% ============================================================
%% ENVIRONNEMENTS TCOLORBOX
%% ============================================================
\newtcolorbox{definitionbox}[1][]{
  breakable, enhanced,
  colback=caplpbluebg, colframe=caplpblue,
  fonttitle=\bfseries\sffamily,
  title={Définition\ifx\relax#1\relax\else\ \textemdash\ #1\fi}
}

\newtcolorbox{theoremebox}[1][]{
  breakable, enhanced,
  colback=caplppurplebg, colframe=caplppurple,
  fonttitle=\bfseries\sffamily,
  title={Théorème\ifx\relax#1\relax\else\ \textemdash\ #1\fi}
}

\newtcolorbox{propositionbox}[1][]{
  breakable, enhanced,
  colback=caplppurplebg, colframe=caplppurple!70!black,
  fonttitle=\bfseries\sffamily,
  title={Proposition\ifx\relax#1\relax\else\ \textemdash\ #1\fi}
}

\newtcolorbox{corollairebox}[1][]{
  breakable, enhanced,
  colback=caplppurplebg, colframe=caplppurple!50!black,
  fonttitle=\bfseries\sffamily,
  title={Corollaire\ifx\relax#1\relax\else\ \textemdash\ #1\fi}
}

\newtcolorbox{lemmebox}[1][]{
  breakable, enhanced,
  colback=caplppurplebg, colframe=caplppurple!40!black,
  fonttitle=\bfseries\sffamily,
  title={Lemme\ifx\relax#1\relax\else\ \textemdash\ #1\fi}
}

\newtcolorbox{preuvebox}{
  breakable, enhanced,
  colback=caplpgraybg, colframe=caplpgray!40,
  fonttitle=\bfseries\itshape\sffamily,
  title={Démonstration},
  after upper={\hfill$\square$}
}

\newtcolorbox{exemplebox}[1][]{
  breakable, enhanced,
  colback=caplpgreenbg, colframe=caplpgreen,
  fonttitle=\bfseries\sffamily,
  title={Exemple\ifx\relax#1\relax\else\ \textemdash\ #1\fi}
}

\newtcolorbox{erreurbox}[1][]{
  breakable, enhanced,
  colback=caplpredbg, colframe=caplpred,
  fonttitle=\bfseries\sffamily,
  title={Erreur fréquente\ifx\relax#1\relax\else\ \textemdash\ #1\fi}
}

\newtcolorbox{methodebox}[1][]{
  breakable, enhanced,
  colback=caplporangebg, colframe=caplporange,
  fonttitle=\bfseries\sffamily,
  title={Méthode\ifx\relax#1\relax\else\ \textemdash\ #1\fi}
}

\newtcolorbox{retenirbox}{
  breakable, enhanced,
  colback=white, colframe=caplpblue,
  coltitle=white, colbacktitle=caplpblue,
  fonttitle=\bfseries\sffamily,
  title={À retenir}, boxrule=1.5pt
}

\newtcolorbox{exercicebox}[1][]{
  breakable, enhanced,
  colback=white, colframe=caplpgray,
  fonttitle=\bfseries\sffamily,
  title={Exercice\ifx\relax#1\relax\else\ \textemdash\ #1\fi}
}

\newtcolorbox{correctionbox}[1][]{
  breakable, enhanced,
  colback=caplpgreenbg, colframe=caplpgreen!60!black,
  fonttitle=\bfseries\itshape\sffamily,
  title={Correction\ifx\relax#1\relax\else\ \textemdash\ #1\fi}
}

\newtcolorbox{remarquebox}[1][]{
  breakable, enhanced,
  colback=caplpgraybg, colframe=caplpgray!50,
  fonttitle=\bfseries\sffamily,
  title={Remarque\ifx\relax#1\relax\else\ \textemdash\ #1\fi}
}

%% ============================================================
%% EN-TÊTES ET PIEDS DE PAGE
%% ============================================================
\pagestyle{fancy}
\fancyhf{}
\fancyhead[LE]{\small\sffamily\leftmark}
\fancyhead[RO]{\small\sffamily\rightmark}
\fancyfoot[C]{\thepage}
\renewcommand{\headrulewidth}{0.4pt}

%% ============================================================
%% LISTES
%% ============================================================
\setlist[itemize]{topsep=4pt,itemsep=3pt,parsep=0pt}
\setlist[enumerate]{topsep=4pt,itemsep=3pt,parsep=0pt}

%% ============================================================
%% RACCOURCIS MATHÉMATIQUES
%% ============================================================
\newcommand{\N}{\mathbb{N}}
\newcommand{\Z}{\mathbb{Z}}
\newcommand{\Q}{\mathbb{Q}}
\newcommand{\R}{\mathbb{R}}
\newcommand{\C}{\mathbb{C}}
\newcommand{\K}{\mathbb{K}}
\newcommand{\ds}{\displaystyle}
\newcommand{\eps}{\varepsilon}
\newcommand{\bigo}{\mathcal{O}}
\DeclareMathOperator{\pgcd}{pgcd}
\DeclareMathOperator{\ppcm}{ppcm}
\DeclareMathOperator{\Card}{Card}
\DeclareMathOperator{\Im}{Im}
\DeclareMathOperator{\Ker}{Ker}

%% ============================================================
%%  DOCUMENT
%% ============================================================
\begin{document}

%% --- PAGE DE GARDE ---
\begin{titlepage}
  \centering
  \vspace*{2cm}
  {\color{caplpblue}\rule{\linewidth}{1.5pt}}\\[0.4cm]
  {\Huge\bfseries\sffamily\color{caplpblue} CAPLP Externe BAC+3}\\[0.3cm]
  {\Large\sffamily Section mathématiques -- physique chimie}\\[0.2cm]
  {\large\sffamily Session 2026}\\[0.2cm]
  {\color{caplpblue}\rule{\linewidth}{1.5pt}}\\[1.2cm]
  {\Huge\bfseries\sffamily Livret de cours}\\[0.4cm]
  {\LARGE\bfseries\sffamily\color{caplppurple} Mathématiques}\\[0.3cm]
  \begin{tcolorbox}[colback=caplpbluebg, colframe=caplpblue, width=7cm,
    halign=center, fontupper=\Large\bfseries\sffamily\color{caplpblue}]
    TOME 1
  \end{tcolorbox}
  \vspace{0.4cm}
  {\large\sffamily Programmes A.1 (partie commune) et A.2 (discipline majeure)}\\[0.3cm]
  {\normalsize\sffamily\color{caplpgray} Partie commune à tous les candidats + Discipline majeure Mathématiques}\\[2cm]
  {\large\sffamily\color{caplpgray}
    Raisonnement $\cdot$ Complexes $\cdot$ Analyse $\cdot$ Algèbre linéaire\\
    Arithmétique $\cdot$ Géométrie $\cdot$ Probabilités $\cdot$ Statistiques $\cdot$ Python\\[0.3cm]
    Fonctions 2 variables $\cdot$ DL $\cdot$ Séries $\cdot$ Diagonalisation
  }\\[2.5cm]
  {\normalsize\sffamily\color{caplpgray}
    Document de préparation au concours\\
    Niveau : Licence (BAC+3)
  }
  \vfill
  {\small\sffamily\color{caplpgray} Compilé avec \LaTeX\ --- Overleaf}
\end{titlepage}

%% --- TABLE DES MATIÈRES ---
\tableofcontents
\newpage

%% ============================================================
%% CHAPITRE 1 — RAISONNEMENT ET VOCABULAIRE ENSEMBLISTE
%% ============================================================
\chapter{Raisonnement et vocabulaire ensembliste}

\section*{Objectifs du chapitre}
\addcontentsline{toc}{section}{Objectifs}

À l'issue de ce chapitre, le candidat doit maîtriser :
\begin{itemize}
  \item les connecteurs logiques et les quantificateurs ;
  \item le vocabulaire de la théorie des ensembles et des applications ;
  \item les principales techniques de raisonnement mathématique :
        disjonction de cas, contradiction, contraposée, récurrence.
\end{itemize}

%% -----------------------------------------------------------
\section{Logique propositionnelle}
%% -----------------------------------------------------------

La logique est le socle invisible de toutes les mathématiques. Avant même de définir un nombre ou une fonction, on utilise des raisonnements — et ces raisonnements suivent des règles précises. Dans ce chapitre, on formalise ces règles pour pouvoir les manipuler rigoureusement et éviter les erreurs de raisonnement qui invalident une démonstration, même si les calculs sont corrects.

\subsection{Propositions et valeurs de vérité}

\begin{definitionbox}[Proposition]
Une \textbf{proposition} (ou assertion) est un énoncé mathématique auquel on peut attribuer
une valeur de vérité : \textbf{Vrai} (V) ou \textbf{Faux} (F).
\end{definitionbox}

\begin{exemplebox}
\begin{itemize}
  \item « $2 + 2 = 4$ » est une proposition vraie.
  \item « $\sqrt{2} \in \Q$ » est une proposition fausse.
  \item « $x > 0$ » n'est pas une proposition (valeur de vérité dépend de $x$) : c'est un
        \textbf{prédicat}.
\end{itemize}
\end{exemplebox}

\subsection{Connecteurs logiques}

Les connecteurs permettent de combiner des propositions simples pour en former de plus complexes. Le plus délicat est l'implication $P \Rightarrow Q$ : elle affirme que chaque fois que $P$ est vraie, $Q$ l'est aussi. Elle ne dit rien sur le cas où $P$ est fausse — d'où son caractère « vacueux » dans ce cas. En mathématiques, presque tous les théorèmes s'écrivent sous forme d'implication : « Si [hypothèses], alors [conclusion] ».

\begin{definitionbox}[Connecteurs logiques]
Soient $P$ et $Q$ deux propositions. On définit :
\begin{itemize}
  \item \textbf{Négation} : $\neg P$ (ou $\overline{P}$), vraie si et seulement si $P$ est fausse.
  \item \textbf{Conjonction} : $P \land Q$ (« $P$ et $Q$ »), vraie si et seulement si $P$ et $Q$
        sont toutes deux vraies.
  \item \textbf{Disjonction} : $P \lor Q$ (« $P$ ou $Q$ »), vraie si et seulement si au moins
        l'une de $P$, $Q$ est vraie.
  \item \textbf{Implication} : $P \Rightarrow Q$, fausse si et seulement si $P$ est vraie et $Q$
        est fausse.
  \item \textbf{Équivalence} : $P \Leftrightarrow Q$, vraie si et seulement si $P$ et $Q$ ont la
        même valeur de vérité.
\end{itemize}
\end{definitionbox}

\begin{center}
\begin{tabular}{cc|ccccc}
\toprule
$P$ & $Q$ & $\neg P$ & $P \land Q$ & $P \lor Q$ & $P \Rightarrow Q$ & $P \Leftrightarrow Q$ \\
\midrule
V & V & F & V & V & V & V \\
V & F & F & F & V & F & F \\
F & V & V & F & V & V & F \\
F & F & V & F & F & V & V \\
\bottomrule
\end{tabular}
\end{center}

\begin{remarquebox}[Implication]
$P \Rightarrow Q$ se lit « $P$ implique $Q$ », ou « si $P$ alors $Q$ ». Elle est équivalente à
$\neg P \lor Q$. \textbf{Une implication dont la prémisse est fausse est toujours vraie}
(vacuité).
\end{remarquebox}

\begin{propositionbox}[Propriétés des connecteurs]
Soient $P$, $Q$, $R$ des propositions. On a :
\begin{itemize}
  \item \textbf{Contraposée} : $(P \Rightarrow Q) \Leftrightarrow (\neg Q \Rightarrow \neg P)$.
  \item \textbf{Lois de De Morgan} :
        $\neg(P \land Q) \Leftrightarrow \neg P \lor \neg Q$ \quad et \quad
        $\neg(P \lor Q) \Leftrightarrow \neg P \land \neg Q$.
  \item \textbf{Transitivité} : $(P \Rightarrow Q) \land (Q \Rightarrow R) \Rightarrow (P \Rightarrow R)$.
\end{itemize}
\end{propositionbox}

\begin{preuvebox}
On vérifie par table de vérité. Pour la contraposée : $P \Rightarrow Q$ est fausse ssi
$P = \text{V}$ et $Q = \text{F}$, ce qui est équivalent à $\neg Q = \text{V}$ et
$\neg P = \text{F}$, i.e.\ $\neg Q \Rightarrow \neg P$ fausse.
\end{preuvebox}

\begin{erreurbox}[Confusion implication et équivalence]
$P \Rightarrow Q$ n'est pas symétrique : de $Q$ on ne peut pas conclure $P$.\\
\textbf{Exemple} : « $n$ pair $\Rightarrow$ $n^2$ pair » est vraie. La réciproque « $n^2$ pair
$\Rightarrow$ $n$ pair » est également vraie ici, mais ce n'est pas automatique.\\
\textbf{Contre-exemple typique} : « $n^2 > 0 \Rightarrow n \neq 0$ » est vraie, mais
« $n \neq 0 \Rightarrow n^2 > 0$ » est fausse dans $\C$.
\end{erreurbox}

\subsection{Quantificateurs}

Les quantificateurs permettent d'exprimer des propriétés portant sur un ou plusieurs éléments d'un ensemble. Ce sont des outils indispensables pour énoncer des définitions et des théorèmes avec précision. La négation d'une proposition quantifiée est une source classique d'erreurs : il faut bien comprendre que nier « pour tout » donne « il existe », et réciproquement.

\begin{definitionbox}[Quantificateurs]
Soit $P(x)$ un prédicat portant sur des éléments d'un ensemble $E$.
\begin{itemize}
  \item \textbf{Quantificateur universel} : $\forall x \in E,\; P(x)$ signifie « pour tout
        $x$ dans $E$, $P(x)$ est vraie ».
  \item \textbf{Quantificateur existentiel} : $\exists x \in E,\; P(x)$ signifie « il existe
        au moins un $x$ dans $E$ tel que $P(x)$ est vraie ».
  \item \textbf{Existentiel unique} : $\exists!\, x \in E,\; P(x)$ signifie « il existe un
        unique $x$ dans $E$ tel que $P(x)$ est vraie ».
\end{itemize}
\end{definitionbox}

\begin{propositionbox}[Négation des quantificateurs]
\[
  \neg\bigl(\forall x \in E,\; P(x)\bigr) \;\Leftrightarrow\; \exists x \in E,\; \neg P(x).
\]
\[
  \neg\bigl(\exists x \in E,\; P(x)\bigr) \;\Leftrightarrow\; \forall x \in E,\; \neg P(x).
\]
\end{propositionbox}

\begin{exemplebox}[Négation d'une proposition quantifiée]
Soit la proposition $P$ : « Toute fonction continue sur $[0,1]$ est bornée. »\\
Formellement : $\forall f \in \mathcal{C}([0,1]),\; f \text{ est bornée}$.\\
Sa négation est : $\exists f \in \mathcal{C}([0,1]),\; f \text{ n'est pas bornée}$.\\
(La proposition $P$ est vraie ; sa négation est donc fausse.)
\end{exemplebox}

\begin{erreurbox}[Ordre des quantificateurs]
L'ordre des quantificateurs est crucial et non commutatif en général :
\[
  \forall x\; \exists y,\; P(x,y) \;\not\Leftrightarrow\; \exists y\; \forall x,\; P(x,y).
\]
\textbf{Exemple} : $\forall x \in \R,\; \exists y \in \R,\; y > x$ est vraie (prendre $y = x+1$).\\
Mais $\exists y \in \R,\; \forall x \in \R,\; y > x$ est fausse ($\R$ n'est pas majoré).
\end{erreurbox}

%% -----------------------------------------------------------
\section{Théorie des ensembles}
%% -----------------------------------------------------------

La théorie des ensembles fournit le langage universel des mathématiques modernes. Tout objet mathématique — nombre, fonction, espace vectoriel — est formellement un ensemble. Il est essentiel de distinguer soigneusement l'appartenance ($\in$) de l'inclusion ($\subseteq$) : $\{1\} \in \{\{1\},2\}$ mais $\{1\} \subseteq \{1,2\}$. Cette distinction est une source fréquente de confusion en début de formation.

\subsection{Ensembles et appartenance}

\begin{definitionbox}[Ensemble, élément, appartenance]
Un \textbf{ensemble} est une collection d'objets appelés \textbf{éléments}.
On note $x \in E$ pour dire que $x$ est un élément de $E$, et $x \notin E$ dans le cas contraire.
L'\textbf{ensemble vide} $\emptyset$ ne contient aucun élément.
\end{definitionbox}

\begin{definitionbox}[Inclusion, égalité]
Soient $A$ et $B$ deux ensembles.
\begin{itemize}
  \item $A \subseteq B$ (« $A$ est inclus dans $B$ ») si $\forall x \in A,\; x \in B$.
  \item $A = B$ si $A \subseteq B$ et $B \subseteq A$.
  \item $A \subsetneq B$ (inclusion stricte) si $A \subseteq B$ et $A \neq B$.
\end{itemize}
\end{definitionbox}

\begin{remarquebox}
Pour montrer l'égalité de deux ensembles, on montre la double inclusion :
$A \subseteq B$ et $B \subseteq A$. C'est la méthode standard.
\end{remarquebox}

\subsection{Opérations ensemblistes}

\begin{definitionbox}[Opérations sur les ensembles]
Soient $A$ et $B$ des parties d'un ensemble $E$.
\begin{itemize}
  \item \textbf{Réunion} : $A \cup B = \{x \in E \mid x \in A \text{ ou } x \in B\}$.
  \item \textbf{Intersection} : $A \cap B = \{x \in E \mid x \in A \text{ et } x \in B\}$.
  \item \textbf{Complémentaire} : $\overline{A} = A^c = \{x \in E \mid x \notin A\}$.
  \item \textbf{Différence} : $A \setminus B = A \cap B^c = \{x \in A \mid x \notin B\}$.
  \item \textbf{Différence symétrique} : $A \triangle B = (A \setminus B) \cup (B \setminus A)$.
  \item \textbf{Produit cartésien} : $A \times B = \{(a,b) \mid a \in A,\; b \in B\}$.
\end{itemize}
\end{definitionbox}

\begin{propositionbox}[Lois de De Morgan ensemblistes]
Pour toute famille finie $(A_i)_{i \in I}$ de parties de $E$ :
\[
  \overline{\bigcup_{i \in I} A_i} = \bigcap_{i \in I} \overline{A_i},
  \qquad
  \overline{\bigcap_{i \in I} A_i} = \bigcup_{i \in I} \overline{A_i}.
\]
\end{propositionbox}

\begin{preuvebox}
Montrons la première égalité. Soit $x \in E$.
\begin{align*}
  x \in \overline{\bigcup_{i \in I} A_i}
  &\Leftrightarrow x \notin \bigcup_{i \in I} A_i \\
  &\Leftrightarrow \forall i \in I,\; x \notin A_i \\
  &\Leftrightarrow \forall i \in I,\; x \in \overline{A_i} \\
  &\Leftrightarrow x \in \bigcap_{i \in I} \overline{A_i}.
\end{align*}
La deuxième égalité se démontre de façon analogue.
\end{preuvebox}

\begin{center}
\begin{tikzpicture}[scale=1.1]
  % Rectangle univers
  \draw[thick, rounded corners=4pt] (-3.2,-1.8) rectangle (3.2,1.8);
  \node[font=\small\itshape] at (2.8,1.5) {$E$};
  % Deux cercles
  \begin{scope}
    \clip (-0.6,-1.3) circle (1.3);
    \fill[caplpbluebg] (0.6,-1.3) circle (1.3);
  \end{scope}
  \draw[thick, caplpblue] (-0.6,0.3) circle (1.3);
  \draw[thick, caplpred]  ( 0.6,0.3) circle (1.3);
  % Étiquettes
  \node[font=\bfseries\small, caplpblue] at (-1.5,0.6) {$A$};
  \node[font=\bfseries\small, caplpred]  at ( 1.5,0.6) {$B$};
  \node[font=\small, caplpgreen] at (0,0.3) {$A\cap B$};
  % Légende dessous
  \node[font=\small] at (0,-1.5) {$A\cup B$ = toute la zone colorée};
  \node[font=\small] at (0,-1.9) {$A\cap B$ = zone centrale (hachuré)};
\end{tikzpicture}
\quad
\begin{tikzpicture}[scale=1.1]
  % Complémentaire
  \fill[caplpgraybg, rounded corners=4pt] (-2.2,-1.4) rectangle (2.2,1.4);
  \fill[white] (0,0) circle (1.0);
  \draw[thick, rounded corners=4pt] (-2.2,-1.4) rectangle (2.2,1.4);
  \draw[thick, caplpblue] (0,0) circle (1.0);
  \node[font=\bfseries\small, caplpblue] at (0,0) {$A$};
  \node[font=\small, caplpgray] at (1.7,1.1) {$\overline{A}$};
  \node[font=\small\itshape] at (1.9,-1.1) {$E$};
  \node[font=\small] at (0,-1.7) {$\overline{A}$ = tout ce qui n'est pas dans $A$};
\end{tikzpicture}
\captionof*{figure}{\small\itshape Diagrammes de Venn : intersection $A\cap B$, union $A\cup B$, complémentaire $\overline{A}$.}
\end{center}

\begin{definitionbox}[Ensemble des parties]
L'\textbf{ensemble des parties} de $E$, noté $\mathcal{P}(E)$ (ou $2^E$), est l'ensemble de
tous les sous-ensembles de $E$ :
\[
  \mathcal{P}(E) = \{A \mid A \subseteq E\}.
\]
Si $E$ est fini et $\Card(E) = n$, alors $\Card(\mathcal{P}(E)) = 2^n$.
\end{definitionbox}

%% -----------------------------------------------------------
\section{Applications}
%% -----------------------------------------------------------

Une application est la formalisation mathématique de la notion de « règle de calcul » qui à chaque entrée associe exactement une sortie. C'est la notion centrale de toutes les mathématiques : fonction, opérateur, transformation, homomorphisme — ce sont tous des applications. L'enjeu de cette section est de comprendre précisément quand une application est « réversible » (bijective), ce qui sera fondamental pour définir des fonctions réciproques, des inverses de matrices, etc.

\begin{definitionbox}[Application]
Soient $E$ et $F$ deux ensembles non vides. Une \textbf{application} $f : E \to F$ est une
relation qui associe à tout élément $x \in E$ un \textbf{unique} élément $f(x) \in F$,
appelé \textbf{image} de $x$ par $f$.

$E$ est appelé \textbf{ensemble de départ} (ou domaine), $F$ l'\textbf{ensemble d'arrivée}
(ou codomaine).
\end{definitionbox}

\begin{definitionbox}[Image directe, image réciproque]
Soit $f : E \to F$ une application.
\begin{itemize}
  \item Pour $A \subseteq E$, l'\textbf{image directe} de $A$ est :
        $f(A) = \{f(x) \mid x \in A\} \subseteq F$.
  \item L'\textbf{image} de $f$ est $\Im(f) = f(E)$.
  \item Pour $B \subseteq F$, l'\textbf{image réciproque} de $B$ est :
        $f^{-1}(B) = \{x \in E \mid f(x) \in B\} \subseteq E$.
\end{itemize}
\end{definitionbox}

\begin{propositionbox}[Propriétés des images]
Soit $f : E \to F$, $A, A' \subseteq E$, $B, B' \subseteq F$. Alors :
\begin{enumerate}
  \item $f(A \cup A') = f(A) \cup f(A')$.
  \item $f(A \cap A') \subseteq f(A) \cap f(A')$ \quad (inclusion en général stricte).
  \item $f^{-1}(B \cup B') = f^{-1}(B) \cup f^{-1}(B')$.
  \item $f^{-1}(B \cap B') = f^{-1}(B) \cap f^{-1}(B')$.
  \item $f^{-1}(B^c) = \bigl(f^{-1}(B)\bigr)^c$.
\end{enumerate}
\end{propositionbox}

\begin{erreurbox}[Image d'une intersection]
On a seulement $f(A \cap A') \subseteq f(A) \cap f(A')$ en général.\\
\textbf{Contre-exemple} : $f : \R \to \R$, $f(x) = x^2$, $A = [1,2]$, $A' = [-2,-1]$.\\
$A \cap A' = \emptyset$, donc $f(A \cap A') = \emptyset$.\\
Mais $f(A) = [1,4]$ et $f(A') = [1,4]$, donc $f(A) \cap f(A') = [1,4] \neq \emptyset$.
\end{erreurbox}

\subsection{Injectivité, surjectivité, bijectivité}

\begin{definitionbox}[Injection, surjection, bijection]
Soit $f : E \to F$.
\begin{itemize}
  \item $f$ est \textbf{injective} si $\forall x, x' \in E,\; f(x) = f(x') \Rightarrow x = x'$.\\
        Équivalent : $\forall x \neq x' \in E,\; f(x) \neq f(x')$.
  \item $f$ est \textbf{surjective} si $\forall y \in F,\; \exists x \in E,\; f(x) = y$.\\
        Équivalent : $\Im(f) = F$.
  \item $f$ est \textbf{bijective} si elle est à la fois injective et surjective.\\
        Équivalent : $\forall y \in F,\; \exists!\, x \in E,\; f(x) = y$.
\end{itemize}
\end{definitionbox}

\begin{theoremebox}[Bijection réciproque]
Si $f : E \to F$ est bijective, il existe une unique application $f^{-1} : F \to E$,
appelée \textbf{bijection réciproque} de $f$, telle que :
\[
  f^{-1} \circ f = \mathrm{id}_E \quad \text{et} \quad f \circ f^{-1} = \mathrm{id}_F.
\]
\end{theoremebox}

\begin{propositionbox}[Composition et bijectivité]
Soient $f : E \to F$ et $g : F \to G$.
\begin{enumerate}
  \item Si $f$ et $g$ sont injectives, alors $g \circ f$ est injective.
  \item Si $f$ et $g$ sont surjectives, alors $g \circ f$ est surjective.
  \item Si $f$ et $g$ sont bijectives, alors $g \circ f$ est bijective et
        $(g \circ f)^{-1} = f^{-1} \circ g^{-1}$.
\end{enumerate}
\end{propositionbox}

\begin{preuvebox}
\textbf{Point 1.} Soient $x, x' \in E$ tels que $(g \circ f)(x) = (g \circ f)(x')$,
i.e.\ $g(f(x)) = g(f(x'))$. Comme $g$ est injective, $f(x) = f(x')$.
Comme $f$ est injective, $x = x'$. Donc $g \circ f$ est injective.\\
\textbf{Points 2 et 3 :} analogues.
\end{preuvebox}

\begin{methodebox}[Montrer qu'une application est bijective]
\begin{enumerate}
  \item \textbf{Méthode 1 (définition)} : montrer injectivité + surjectivité séparément.
  \item \textbf{Méthode 2 (inverse)} : exhiber $g : F \to E$ telle que $g \circ f = \mathrm{id}_E$
        et $f \circ g = \mathrm{id}_F$.
  \item \textbf{Méthode 3 (ensembles finis)} : si $\Card(E) = \Card(F) < +\infty$, il suffit
        de montrer l'injectivité ou la surjectivité.
\end{enumerate}
\end{methodebox}

%% -----------------------------------------------------------
\section{Techniques de raisonnement}
%% -----------------------------------------------------------

Savoir calculer ne suffit pas : il faut aussi savoir \emph{démontrer}. Une démonstration est un enchaînement logique irréfutable d'étapes qui conduit de l'hypothèse à la conclusion. Il existe plusieurs stratégies, et le choix de l'une plutôt que de l'autre dépend de la structure de ce qu'on veut prouver. La maîtrise de ces techniques est attendue au CAPLP : on ne peut pas se contenter de « vérifier sur des exemples » ou d'arguer par l'intuition.

\subsection{Raisonnement direct}

La méthode la plus naturelle : on suppose $P$ vraie et on enchaîne des implications logiques jusqu'à conclure $Q$. C'est la démarche par défaut — si elle fonctionne, c'est la plus élégante. Elle s'applique bien quand les hypothèses contiennent suffisamment d'information pour « construire » la conclusion.

\subsection{Raisonnement par contraposée}

La contraposée est utile quand la négation de la conclusion ($\neg Q$) est plus riche en information que la prémisse $P$. On suppose alors $\neg Q$ et on démontre $\neg P$ — ce qui est logiquement équivalent à $P \Rightarrow Q$.

\begin{propositionbox}
Pour démontrer $P \Rightarrow Q$, il est équivalent de démontrer $\neg Q \Rightarrow \neg P$.
\end{propositionbox}

\begin{exemplebox}[Contraposée]
\textbf{Théorème} : Si $n^2$ est impair, alors $n$ est impair.\\
\textbf{Contraposée} : Si $n$ est pair, alors $n^2$ est pair.\\
\textbf{Preuve de la contraposée} : Si $n$ est pair, $n = 2k$ pour un certain $k \in \Z$,
donc $n^2 = 4k^2 = 2(2k^2)$ est pair.
\end{exemplebox}

\subsection{Raisonnement par l'absurde}

Le raisonnement par l'absurde est particulièrement puissant pour démontrer qu'un objet \emph{n'existe pas} ou qu'une propriété est \emph{impossible}. L'idée est simple : on suppose le contraire de ce qu'on veut prouver, et on montre que cette hypothèse mène à une contradiction — ce qui est impossible, donc la supposition est fausse, donc ce qu'on voulait prouver est vrai. La démonstration de l'irrationalité de $\sqrt{2}$, présentée ci-dessous, est un modèle du genre.

\begin{methodebox}[Raisonnement par l'absurde]
Pour démontrer $P$, on suppose $\neg P$ et on en déduit une contradiction
(une proposition fausse ou une proposition et sa négation simultanément vraies).
\end{methodebox}

\begin{exemplebox}[Irrationalité de $\sqrt{2}$]
\textbf{Théorème} : $\sqrt{2} \notin \Q$.

\textbf{Preuve par l'absurde} : Supposons $\sqrt{2} \in \Q$.
Il existe $p \in \Z$ et $q \in \N^*$ avec $\pgcd(p,q) = 1$ tels que $\sqrt{2} = p/q$.
Alors $2q^2 = p^2$, donc $p^2$ est pair, donc $p$ est pair (voir contraposée).
Écrivons $p = 2k$. Alors $2q^2 = 4k^2$, i.e.\ $q^2 = 2k^2$, donc $q^2$ est pair,
donc $q$ est pair. Ainsi $2 \mid p$ et $2 \mid q$, ce qui contredit $\pgcd(p,q)=1$.
Contradiction. Donc $\sqrt{2} \notin \Q$.
\end{exemplebox}

\subsection{Raisonnement par disjonction de cas}

\begin{methodebox}[Disjonction de cas]
Pour démontrer $P$, on partitionne l'univers en cas $C_1, C_2, \ldots, C_n$
(exhaustifs et mutuellement exclusifs) et on démontre $P$ dans chacun des cas.
\end{methodebox}

\begin{exemplebox}[Disjonction de cas]
\textbf{Théorème} : Pour tout $n \in \Z$, $n(n+1)$ est pair.

\textbf{Preuve} : Soit $n \in \Z$.
\begin{itemize}
  \item \textbf{Cas 1} : $n$ est pair. Alors $n = 2k$, donc $n(n+1) = 2k(n+1)$ est pair.
  \item \textbf{Cas 2} : $n$ est impair. Alors $n+1$ est pair, soit $n+1 = 2k$,
        donc $n(n+1) = 2kn$ est pair.
\end{itemize}
Dans les deux cas, $n(n+1)$ est pair.
\end{exemplebox}

\subsection{Raisonnement par récurrence}

La récurrence est la technique de prédilection pour démontrer des propriétés portant sur tous les entiers naturels. Elle repose sur un principe fondamental : si une propriété est vraie pour $n=0$ (initialisation), et si chaque fois qu'elle est vraie pour $n$ elle l'est aussi pour $n+1$ (hérédité), alors elle est vraie pour tout entier. C'est une forme de « domino » : on fait tomber le premier, et chaque domino en fait tomber le suivant. Il est absolument indispensable de rédiger les trois étapes (initialisation, hypothèse de récurrence, hérédité) de façon explicite et complète.

\begin{theoremebox}[Principe de récurrence simple]
Soit $P(n)$ un prédicat portant sur $\N$ (ou sur $\{n_0, n_0+1, \ldots\}$).
Si :
\begin{enumerate}[label=(\roman*)]
  \item \textbf{Initialisation} : $P(n_0)$ est vraie,
  \item \textbf{Hérédité} : $\forall n \geq n_0,\; P(n) \Rightarrow P(n+1)$,
\end{enumerate}
alors $P(n)$ est vraie pour tout entier $n \geq n_0$.
\end{theoremebox}

\begin{theoremebox}[Récurrence forte]
Soit $P(n)$ un prédicat sur $\N$. Si :
\begin{enumerate}[label=(\roman*)]
  \item $P(n_0)$ est vraie,
  \item $\forall n \geq n_0,\; \bigl(P(n_0) \land P(n_0+1) \land \cdots \land P(n)\bigr)
        \Rightarrow P(n+1)$,
\end{enumerate}
alors $P(n)$ est vraie pour tout entier $n \geq n_0$.
\end{theoremebox}

\begin{exemplebox}[Somme des entiers par récurrence]
\textbf{Théorème} : Pour tout $n \in \N^*$,
\[
  \sum_{k=1}^{n} k = \frac{n(n+1)}{2}.
\]
\textbf{Preuve par récurrence} :
\begin{itemize}
  \item \textbf{Initialisation} ($n=1$) : $\sum_{k=1}^{1} k = 1 = \frac{1 \times 2}{2}$. \checkmark
  \item \textbf{Hérédité} : Supposons la formule vraie au rang $n \geq 1$. Alors :
  \[
    \sum_{k=1}^{n+1} k = \sum_{k=1}^{n} k + (n+1)
    = \frac{n(n+1)}{2} + (n+1)
    = (n+1)\left(\frac{n}{2} + 1\right)
    = \frac{(n+1)(n+2)}{2}.
  \]
  La formule est vraie au rang $n+1$.
\end{itemize}
Par le principe de récurrence, la formule est établie pour tout $n \geq 1$.
\end{exemplebox}

\begin{methodebox}[Rédiger une récurrence]
\begin{enumerate}
  \item Énoncer clairement la proposition $P(n)$ que l'on veut démontrer.
  \item \textbf{Initialisation} : vérifier $P(n_0)$ (calcul explicite).
  \item \textbf{Hérédité} : écrire « Soit $n \geq n_0$. Supposons $P(n)$ vraie. »,
        puis démontrer $P(n+1)$ en utilisant $P(n)$.
  \item Conclure par « Par le principe de récurrence, $P(n)$ est vraie pour tout $n \geq n_0$. »
\end{enumerate}
\end{methodebox}

\begin{erreurbox}[Erreurs classiques en récurrence]
\begin{itemize}
  \item \textbf{Oublier l'initialisation} : l'hérédité seule ne suffit pas.
  \item \textbf{Utiliser ce qu'on veut démontrer} : en hérédité, seul $P(n)$ (ou
        $P(n_0), \ldots, P(n)$ en forte) est supposé ; $P(n+1)$ doit être déduit.
  \item \textbf{Raisonner sur $n$ au lieu de $n+1$} : l'étape d'hérédité doit
        conclure sur le rang $n+1$.
\end{itemize}
\end{erreurbox}

%% -----------------------------------------------------------
\section{Exercices types}
%% -----------------------------------------------------------

\begin{exercicebox}[Négation et quantificateurs]
Écrire la négation des propositions suivantes, puis dire si elles sont vraies ou fausses.
\begin{enumerate}
  \item $P_1$ : « Pour tout réel $x$, $x^2 \geq 0$. »
  \item $P_2$ : « Il existe un entier $n$ tel que $n^2 = 2$. »
  \item $P_3$ : « Pour tout $\eps > 0$, il existe $n \in \N$ tel que $\frac{1}{n} < \eps$. »
  \item $P_4$ : « Pour tout $x \in \R$, pour tout $y \in \R$, $x + y > 0$. »
\end{enumerate}
\end{exercicebox}

\begin{correctionbox}
\begin{enumerate}
  \item $\neg P_1$ : « Il existe un réel $x$ tel que $x^2 < 0$. »\\
        $P_1$ est \textbf{vraie} (carré d'un réel toujours positif ou nul) ;
        $\neg P_1$ est fausse.

  \item $\neg P_2$ : « Pour tout entier $n$, $n^2 \neq 2$. »\\
        $P_2$ est \textbf{fausse} ($\sqrt{2} \notin \Z$) ; $\neg P_2$ est vraie.

  \item $\neg P_3$ : « Il existe $\eps > 0$ tel que pour tout $n \in \N$,
        $\frac{1}{n} \geq \eps$. »\\
        $P_3$ est \textbf{vraie} (propriété archimédienne de $\R$) ;
        $\neg P_3$ est fausse.

  \item $\neg P_4$ : « Il existe $x \in \R$ et $y \in \R$ tels que $x + y \leq 0$. »\\
        $P_4$ est \textbf{fausse} (prendre $x = -1, y = -1$) ; $\neg P_4$ est vraie.
\end{enumerate}
\end{correctionbox}

\begin{exercicebox}[Applications — injection, surjection]
Étudier l'injectivité et la surjectivité des applications suivantes. En cas de bijectivité,
déterminer la bijection réciproque.
\begin{enumerate}
  \item $f : \R \to \R$, $f(x) = 2x - 3$.
  \item $g : \R \to \R$, $g(x) = x^2$.
  \item $h : \R_+ \to \R_+$, $h(x) = x^2$.
  \item $\varphi : \Z \to \Z$, $\varphi(n) = 2n$.
\end{enumerate}
\end{exercicebox}

\begin{correctionbox}
\begin{enumerate}
  \item \textbf{Injection} : $f(x) = f(x') \Rightarrow 2x-3 = 2x'-3 \Rightarrow x = x'$. \\
        \textbf{Surjection} : Soit $y \in \R$. Posons $x = \frac{y+3}{2} \in \R$. Alors
        $f(x) = 2\cdot\frac{y+3}{2} - 3 = y$. \\
        Donc $f$ est \textbf{bijective}. Réciproque : $f^{-1}(y) = \frac{y+3}{2}$.

  \item \textbf{Non injective} : $g(1) = g(-1) = 1$ avec $1 \neq -1$. \\
        \textbf{Non surjective} : $-1 \in \R$ n'a pas d'antécédent ($x^2 \geq 0$).

  \item \textbf{Injection} : $h(x) = h(x') \Rightarrow x^2 = x'^2$. Sur $\R_+$, $x = x'$.\\
        \textbf{Surjection} : Soit $y \in \R_+$. Posons $x = \sqrt{y} \in \R_+$.
        Alors $h(x) = y$. \\
        Donc $h$ est \textbf{bijective}. Réciproque : $h^{-1}(y) = \sqrt{y}$.

  \item \textbf{Injective} : $\varphi(n) = \varphi(m) \Rightarrow 2n = 2m \Rightarrow n = m$.\\
        \textbf{Non surjective} : $1 \in \Z$ n'a pas d'antécédent dans $\Z$ ($2n = 1$
        n'a pas de solution entière).
\end{enumerate}
\end{correctionbox}

\begin{exercicebox}[Récurrence — type concours]
\begin{enumerate}
  \item Démontrer que pour tout entier $n \geq 1$ :
  \[
    \sum_{k=0}^{n} k^2 = \frac{n(n+1)(2n+1)}{6}.
  \]
  \item Démontrer que pour tout entier $n \geq 1$, $7^n - 1$ est divisible par $6$.
  \item (Récurrence forte) Soit la suite définie par $u_0 = 0$, $u_1 = 1$ et
        $u_{n+2} = u_{n+1} + u_n$ pour tout $n \geq 0$ (suite de Fibonacci).
        Démontrer que $\pgcd(u_n, u_{n+1}) = 1$ pour tout $n \geq 0$.
\end{enumerate}
\end{exercicebox}

\begin{correctionbox}
\textbf{Question 1.} Notons $P(n)$ : $\ds\sum_{k=0}^{n} k^2 = \frac{n(n+1)(2n+1)}{6}$.

\textit{Initialisation} ($n=1$) : $\sum_{k=0}^{1} k^2 = 0 + 1 = 1 = \frac{1 \cdot 2 \cdot 3}{6}$. \checkmark

\textit{Hérédité} : Supposons $P(n)$ vraie. Alors :
\begin{align*}
  \sum_{k=0}^{n+1} k^2
  &= \sum_{k=0}^{n} k^2 + (n+1)^2
  = \frac{n(n+1)(2n+1)}{6} + (n+1)^2 \\
  &= \frac{(n+1)}{6}\bigl[n(2n+1) + 6(n+1)\bigr]
  = \frac{(n+1)(2n^2+7n+6)}{6}
  = \frac{(n+1)(n+2)(2n+3)}{6}.
\end{align*}
C'est bien $P(n+1)$. Par récurrence, $P(n)$ est vraie pour tout $n \geq 1$.

\medskip
\textbf{Question 2.} Notons $P(n)$ : $6 \mid 7^n - 1$.

\textit{Initialisation} ($n=1$) : $7^1 - 1 = 6 = 6 \times 1$. \checkmark

\textit{Hérédité} : Supposons $6 \mid 7^n - 1$, i.e.\ $7^n - 1 = 6k$ pour un certain
$k \in \Z$. Alors :
\[
  7^{n+1} - 1 = 7 \cdot 7^n - 1 = 7(7^n - 1) + 7 - 1 = 7 \cdot 6k + 6 = 6(7k + 1).
\]
Donc $6 \mid 7^{n+1} - 1$. Par récurrence, $6 \mid 7^n - 1$ pour tout $n \geq 1$.

\medskip
\textbf{Question 3.} Notons $P(n)$ : $\pgcd(u_n, u_{n+1}) = 1$.

\textit{Initialisation} ($n=0$) : $\pgcd(u_0, u_1) = \pgcd(0,1) = 1$. \checkmark

\textit{Hérédité} : Supposons $\pgcd(u_n, u_{n+1}) = 1$. Alors :
\[
  \pgcd(u_{n+1}, u_{n+2}) = \pgcd(u_{n+1}, u_{n+1} + u_n).
\]
Or, pour tous entiers $a, b$ : $\pgcd(a, a+b) = \pgcd(a, b)$. Donc :
\[
  \pgcd(u_{n+1}, u_{n+2}) = \pgcd(u_{n+1}, u_n) = \pgcd(u_n, u_{n+1}) = 1.
\]
Par récurrence, $\pgcd(u_n, u_{n+1}) = 1$ pour tout $n \geq 0$.
\end{correctionbox}

%% -----------------------------------------------------------
\section{Bilan — À retenir}
%% -----------------------------------------------------------

\begin{retenirbox}
\textbf{Connecteurs et quantificateurs}
\begin{itemize}
  \item $P \Rightarrow Q \Leftrightarrow \neg Q \Rightarrow \neg P$ (contraposée).
  \item $\neg(\forall x,\; P(x)) \Leftrightarrow \exists x,\; \neg P(x)$.
  \item $\neg(\exists x,\; P(x)) \Leftrightarrow \forall x,\; \neg P(x)$.
  \item L'ordre des quantificateurs est essentiel.
\end{itemize}

\textbf{Ensembles}
\begin{itemize}
  \item $A = B \Leftrightarrow A \subseteq B \text{ et } B \subseteq A$.
  \item Lois de De Morgan : $\overline{A \cup B} = \overline{A} \cap \overline{B}$.
  \item $f(A \cap B) \subseteq f(A) \cap f(B)$ (pas d'égalité en général).
  \item $f^{-1}$ commute avec $\cup$, $\cap$ et le complémentaire.
\end{itemize}

\textbf{Applications}
\begin{itemize}
  \item Injective : un antécédent au plus.
  \item Surjective : au moins un antécédent.
  \item Bijective : exactement un antécédent. La réciproque existe et est unique.
\end{itemize}

\textbf{Récurrence}
\begin{itemize}
  \item Toujours : initialisation + hérédité + conclusion.
  \item Récurrence forte si $P(n+1)$ nécessite $P(0), \ldots, P(n)$.
\end{itemize}
\end{retenirbox}


%% === CHAPITRE 2 ===
\chapter{Nombres complexes}

\section*{Objectifs}
\begin{itemize}[leftmargin=1.5em]
  \item Maîtriser les différentes formes d'un nombre complexe (algébrique, trigonométrique, exponentielle).
  \item Calculer module, argument, conjugué ; utiliser l'exponentielle complexe.
  \item Retrouver les formules trigonométriques classiques via les complexes.
  \item Résoudre les équations du second degré à coefficients réels dans $\C$.
  \item Calculer les racines $n$-ièmes d'un complexe.
\end{itemize}

\section{Le corps des nombres complexes}

Pendant longtemps, l'équation $x^2 = -1$ n'avait pas de solution : aucun nombre réel ne vérifie cette propriété. Pour y remédier, les mathématiciens ont étendu $\R$ en ajoutant un nouvel élément $i$ tel que $i^2 = -1$. Cette extension est $\C$, le corps des nombres complexes. Loin d'être une curiosité abstraite, $\C$ est fondamental en physique (signaux, oscillations, électricité), en analyse (intégration, équations différentielles) et en géométrie. Dans $\C$, tout polynôme de degré $n$ admet exactement $n$ racines (théorème de d'Alembert-Gauss) — propriété que $\R$ n'a pas.

\begin{definitionbox}[Corps $\C$]
L'ensemble des \emph{nombres complexes} est $\C = \{a + ib \mid a, b \in \R\}$
où $i$ est un élément vérifiant $i^2 = -1$.
Pour $z = a + ib$ :
\begin{itemize}
  \item $\mathrm{Re}(z) = a$ est la \emph{partie réelle} ;
  \item $\mathrm{Im}(z) = b$ est la \emph{partie imaginaire} ;
  \item $\bar{z} = a - ib$ est le \emph{conjugué} de $z$ ;
  \item $|z| = \sqrt{a^2 + b^2}$ est le \emph{module} de $z$.
\end{itemize}
$\C$ est muni de l'addition et de la multiplication usuelles (avec $i^2=-1$) qui en font un corps commutatif contenant $\R$.
\end{definitionbox}

\begin{theoremebox}[Propriétés du conjugué et du module]
Pour tous $z, w \in \C$ :
\begin{align*}
  \overline{z+w} &= \bar{z}+\bar{w}, &\overline{zw} &= \bar{z}\,\bar{w}, &\overline{\bar{z}} &= z, \\
  |zw| &= |z|\,|w|, &\left|\tfrac{z}{w}\right| &= \tfrac{|z|}{|w|}\;(w\neq0), &|z|^2 &= z\bar{z}.
\end{align*}
De plus : $\mathrm{Re}(z) = \dfrac{z+\bar{z}}{2}$ et $\mathrm{Im}(z) = \dfrac{z-\bar{z}}{2i}$.
\end{theoremebox}

\begin{definitionbox}[Inverse d'un complexe non nul]
Pour $z = a+ib \neq 0$ :
\[
  z^{-1} = \frac{\bar{z}}{|z|^2} = \frac{a - ib}{a^2+b^2}.
\]
\end{definitionbox}

\section{Représentation géométrique et forme trigonométrique}

La grande force des nombres complexes est leur double nature : algébrique (forme $a+ib$) et géométrique (point du plan). Cette dualité permet de traduire des propriétés algébriques en propriétés géométriques et inversement. En particulier, la multiplication par $e^{i\theta}$ correspond à une rotation d'angle $\theta$ dans le plan — ce qui donne une interprétation géométrique immédiate des transformations complexes et des formules trigonométriques.

\begin{definitionbox}[Plan complexe, affixe]
On représente $z = a+ib$ par le point $M(a,b)$ du plan muni d'un repère orthonormé
$(O,\vec{e}_1,\vec{e}_2)$, appelé \emph{plan complexe}. $z$ est l'\emph{affixe} de $M$.
\end{definitionbox}

\begin{center}
\begin{tikzpicture}[scale=1.2,>=Stealth]
  \draw[->] (-0.5,0) -- (3,0) node[right,font=\small] {$\mathrm{Re}$};
  \draw[->] (0,-0.5) -- (0,2.5) node[above,font=\small] {$\mathrm{Im}$};
  % Point z
  \coordinate (Z) at (2,1.6);
  \fill[caplpblue] (Z) circle (2pt) node[above right,font=\small] {$M(a,b)$};
  % Module
  \draw[thick, caplpblue] (0,0) -- (Z) node[midway, above left, font=\small] {$|z|$};
  % Projections
  \draw[dashed, gray] (Z) -- (2,0) node[below,font=\scriptsize] {$a$};
  \draw[dashed, gray] (Z) -- (0,1.6) node[left,font=\scriptsize] {$b$};
  % Angle
  \draw[->, caplpgreen] (0.6,0) arc[start angle=0, end angle=38.66, radius=0.6]
    node[right,font=\scriptsize, caplpgreen] {$\theta$};
  % Cercle unitaire (partiel)
  \draw[dotted, gray] (0,0) circle (1);
  \node[font=\scriptsize, gray] at (0.75,-0.15) {$1$};
  \node[font=\small] at (1,-1.2) {$z = a+ib = |z|e^{i\theta}$, \quad $a = |z|\cos\theta$, $b = |z|\sin\theta$};
\end{tikzpicture}
\captionof*{figure}{\small\itshape Représentation de $z=a+ib$ dans le plan complexe.}
\end{center}

\begin{definitionbox}[Argument d'un complexe non nul]
Soit $z \neq 0$. Tout réel $\theta$ vérifiant
\[
  z = |z|(\cos\theta + i\sin\theta)
\]
est un \emph{argument} de $z$, noté $\arg(z)$. L'argument est défini modulo $2\pi$.
L'\emph{argument principal} est l'unique $\theta \in ]-\pi,\pi]$.
\end{definitionbox}

\begin{theoremebox}[Propriétés de l'argument]
Pour $z, w \in \C^*$ :
\[
  \arg(zw) \equiv \arg(z)+\arg(w) \pmod{2\pi}, \quad
  \arg\!\left(\tfrac{z}{w}\right) \equiv \arg(z)-\arg(w) \pmod{2\pi}, \quad
  \arg(\bar{z}) \equiv -\arg(z) \pmod{2\pi}.
\]
\end{theoremebox}

\section{Exponentielle complexe et formules trigonométriques}

La formule d'Euler $e^{i\theta} = \cos\theta + i\sin\theta$ est l'un des résultats les plus remarquables des mathématiques. Elle relie en une seule expression l'exponentielle, la trigonométrie et les nombres complexes. En particulier, pour $\theta = \pi$, on obtient $e^{i\pi} + 1 = 0$, qui réunit les cinq constantes fondamentales des mathématiques. Au-delà de l'élégance, cette formule est un outil de calcul redoutable : les formules trigonométriques les plus pénibles à démontrer classiquement deviennent immédiates par manipulation des exponentielles complexes.

\begin{definitionbox}[Exponentielle complexe]
Pour tout $\theta \in \R$, on définit :
\[
  e^{i\theta} = \cos\theta + i\sin\theta \quad \textbf{(formule d'Euler)}.
\]
Plus généralement, pour $z = a + ib$ :
\[
  e^z = e^a(\cos b + i\sin b).
\]
\end{definitionbox}

\begin{theoremebox}[Formule de Moivre]
Pour tout $\theta \in \R$ et $n \in \Z$ :
\[
  \left(e^{i\theta}\right)^n = e^{in\theta}, \quad \text{soit} \quad
  (\cos\theta + i\sin\theta)^n = \cos(n\theta) + i\sin(n\theta).
\]
\end{theoremebox}

\begin{methodebox}[Formules trigonométriques via les complexes]
On dispose de deux méthodes :
\begin{enumerate}
  \item \textbf{Linéarisation} : exprimer $\cos^n\theta$ ou $\sin^n\theta$ en fonctions de $\cos(k\theta)$ ou $\sin(k\theta)$. On utilise :
    \[ \cos\theta = \frac{e^{i\theta}+e^{-i\theta}}{2}, \qquad \sin\theta = \frac{e^{i\theta}-e^{-i\theta}}{2i}. \]
  \item \textbf{Développement de $e^{in\theta}$} : développer $(\cos\theta+i\sin\theta)^n$ par le binôme, puis identifier parties réelle et imaginaire pour obtenir $\cos(n\theta)$ et $\sin(n\theta)$.
\end{enumerate}
\end{methodebox}

\begin{exemplebox}[Formules d'addition et linéarisation]
\textbf{Formule d'addition.} $e^{i(\alpha+\beta)} = e^{i\alpha}e^{i\beta}$, soit :
\[ \cos(\alpha+\beta)+i\sin(\alpha+\beta) = (\cos\alpha+i\sin\alpha)(\cos\beta+i\sin\beta). \]
En identifiant : $\cos(\alpha+\beta) = \cos\alpha\cos\beta - \sin\alpha\sin\beta$,
$\sin(\alpha+\beta) = \sin\alpha\cos\beta + \cos\alpha\sin\beta$.

\textbf{Linéarisation de $\cos^2\theta$.}
$\cos^2\theta = \left(\dfrac{e^{i\theta}+e^{-i\theta}}{2}\right)^2
= \dfrac{e^{2i\theta}+2+e^{-2i\theta}}{4} = \dfrac{1+\cos(2\theta)}{2}$.
\end{exemplebox}

\section{Racines $n$-ièmes d'un complexe}

Dans $\R$, certains nombres n'ont pas de racine carrée (les négatifs) et $1$ n'a que deux racines cubiques ($1$ et $-1$). Dans $\C$, la situation est parfaitement symétrique : tout nombre complexe non nul admet exactement $n$ racines $n$-ièmes distinctes, régulièrement espacées sur un cercle de rayon $r^{1/n}$. Ce résultat est une conséquence directe du théorème de d'Alembert-Gauss.

\begin{theoremebox}[Racines $n$-ièmes]
Soit $z_0 = r e^{i\varphi} \in \C^*$ et $n \geq 1$ entier.
Les \emph{racines $n$-ièmes} de $z_0$ sont les $n$ complexes :
\[
  z_k = r^{1/n}\, e^{i\frac{\varphi + 2k\pi}{n}}, \quad k = 0, 1, \ldots, n-1.
\]
Elles forment un polygone régulier à $n$ sommets inscrit dans le cercle de rayon $r^{1/n}$.
\end{theoremebox}

\begin{definitionbox}[Racines $n$-ièmes de l'unité]
Les racines $n$-ièmes de $1$ sont $\omega_k = e^{2ik\pi/n}$, $k=0,\ldots,n-1$.
On pose $\omega = e^{2i\pi/n}$ (racine primitive) ; les $n$ racines sont $1, \omega, \omega^2, \ldots, \omega^{n-1}$.

Propriété remarquable : $\displaystyle\sum_{k=0}^{n-1}\omega^k = 0$.
\end{definitionbox}

\begin{center}
\begin{tikzpicture}[scale=1.4,>=Stealth]
  \draw[->] (-1.4,0) -- (1.4,0) node[right,font=\scriptsize] {$\mathrm{Re}$};
  \draw[->] (0,-1.4) -- (0,1.4) node[above,font=\scriptsize] {$\mathrm{Im}$};
  \draw[dotted] (0,0) circle(1);
  % Racines 6ièmes de l'unité
  \foreach \k in {0,...,5}{
    \pgfmathsetmacro\angle{60*\k}
    \pgfmathsetmacro\cx{cos(\angle)}
    \pgfmathsetmacro\cy{sin(\angle)}
    \fill[caplpblue] (\cx,\cy) circle(2pt);
    \draw[thin, caplpblue!40] (0,0) -- (\cx,\cy);
  }
  \node[font=\scriptsize] at (1.15,0.1) {$1$};
  \node[font=\scriptsize] at (0.6,1.1) {$\omega$};
  \node[font=\scriptsize] at (-0.6,1.1) {$\omega^2$};
  \node[font=\scriptsize] at (-1.3,0.1) {$\omega^3$};
  \node[font=\small,below] at (0,-1.5) {Racines $6$-ièmes de l'unité ($n=6$)};
\end{tikzpicture}
\captionof*{figure}{\small\itshape Les $n$ racines $n$-ièmes de l'unité forment un polygone régulier inscrit dans le cercle unité.}
\end{center}

\section{Équation du second degré à coefficients réels}

\begin{theoremebox}[Équation $az^2 + bz + c = 0$, $a,b,c \in \R$]
Soit $\Delta = b^2 - 4ac \in \R$.
\begin{itemize}
  \item $\Delta > 0$ : deux racines réelles distinctes $z = \dfrac{-b \pm \sqrt{\Delta}}{2a}$.
  \item $\Delta = 0$ : une racine réelle double $z = -\dfrac{b}{2a}$.
  \item $\Delta < 0$ : deux racines complexes conjuguées $z = \dfrac{-b \pm i\sqrt{|\Delta|}}{2a}$.
\end{itemize}
Dans tous les cas : somme des racines $= -\dfrac{b}{a}$, produit $= \dfrac{c}{a}$.
\end{theoremebox}

\begin{erreurbox}[Pièges fréquents]
\begin{itemize}
  \item \textbf{Erreur 1.} Écrire $\arg(z_1 z_2) = \arg(z_1)\arg(z_2)$ (produit au lieu de somme).
  \item \textbf{Erreur 2.} Oublier que $\arg$ est défini modulo $2\pi$ : ne pas écrire $\arg(z) = \theta$ sans préciser l'intervalle.
  \item \textbf{Erreur 3.} Pour $\Delta < 0$ avec $a,b,c$ réels, les racines sont toujours \emph{conjuguées l'une de l'autre} — ne pas les chercher indépendamment.
\end{itemize}
\end{erreurbox}

\section{Exercices types}

\begin{exercicebox}[Formes et calculs]
Soit $z_1 = 1+i$ et $z_2 = \sqrt{3}-i$.
\begin{enumerate}
  \item Mettre $z_1$ et $z_2$ sous forme exponentielle $re^{i\theta}$.
  \item Calculer $z_1 z_2$, $\dfrac{z_1}{z_2}$, et $z_1^6$ sous forme algébrique.
  \item Calculer $\arg(z_1^6)$ et $|z_1^6|$.
\end{enumerate}
\end{exercicebox}

\begin{correctionbox}[Correction]
\textbf{1.} $|z_1|=\sqrt{2}$, $\arg(z_1)=\frac{\pi}{4}$, donc $z_1=\sqrt{2}\,e^{i\pi/4}$.
$|z_2|=2$, $\arg(z_2)=-\frac{\pi}{6}$, donc $z_2=2e^{-i\pi/6}$.

\textbf{2.} $z_1 z_2 = 2\sqrt{2}\,e^{i(\pi/4-\pi/6)} = 2\sqrt{2}\,e^{i\pi/12}$.
En algébrique : $z_1 z_2 = (\sqrt{3}+1) + i(\sqrt{3}-1)$.
$\dfrac{z_1}{z_2} = \dfrac{\sqrt{2}}{2}e^{i5\pi/12}$.
$z_1^6 = (\sqrt{2})^6 e^{i6\pi/4} = 8\,e^{i3\pi/2} = 8(0-i) = -8i$.

\textbf{3.} $|z_1^6|=8$, $\arg(z_1^6) = -\frac{\pi}{2}$ (argument principal).
\end{correctionbox}

\begin{exercicebox}[Linéarisation]
Exprimer $\cos^3\theta$ en fonction de $\cos(3\theta)$ et $\cos\theta$.
\end{exercicebox}

\begin{correctionbox}[Correction]
$\cos\theta = \dfrac{e^{i\theta}+e^{-i\theta}}{2}$, donc :
\[
\cos^3\theta = \frac{(e^{i\theta}+e^{-i\theta})^3}{8}
= \frac{e^{3i\theta}+3e^{i\theta}+3e^{-i\theta}+e^{-3i\theta}}{8}
= \frac{2\cos(3\theta)+6\cos\theta}{8} = \frac{\cos(3\theta)}{4}+\frac{3\cos\theta}{4}.
\]
\end{correctionbox}

\begin{exercicebox}[Racines et équation — type concours]
\begin{enumerate}
  \item Trouver les racines cubiques de $z_0 = 8i$.
  \item Résoudre dans $\C$ : $z^2 - 2z + 5 = 0$.
  \item En déduire une factorisation de $P(z)=z^2-2z+5$ sur $\C$.
\end{enumerate}
\end{exercicebox}

\begin{correctionbox}[Correction]
\textbf{1.} $8i = 8e^{i\pi/2}$. Racines cubiques : $z_k = 2e^{i(\pi/6+2k\pi/3)}$, $k=0,1,2$.
$z_0 = 2e^{i\pi/6} = \sqrt{3}+i$,\; $z_1 = 2e^{i5\pi/6} = -\sqrt{3}+i$,\; $z_2 = 2e^{-i\pi/2} = -2i$.

\textbf{2.} $\Delta = 4-20 = -16$. Racines : $z = \dfrac{2 \pm 4i}{2} = 1 \pm 2i$.

\textbf{3.} $P(z) = (z-(1+2i))(z-(1-2i))$.
\end{correctionbox}

\begin{retenirbox}
\textbf{Nombres complexes — À retenir}
\begin{itemize}
  \item $z = a+ib = re^{i\theta}$ avec $r=|z|$, $\theta=\arg(z)$.
  \item $|z|^2 = z\bar{z}$ ; $z^{-1} = \bar{z}/|z|^2$.
  \item Euler : $e^{i\theta} = \cos\theta + i\sin\theta$.
  \item $\cos\theta = \frac{e^{i\theta}+e^{-i\theta}}{2}$, $\sin\theta = \frac{e^{i\theta}-e^{-i\theta}}{2i}$.
  \item Moivre : $(e^{i\theta})^n = e^{in\theta}$.
  \item Racines $n$-ièmes de $re^{i\varphi}$ : $r^{1/n}e^{i(\varphi+2k\pi)/n}$, $k=0,\ldots,n-1$.
  \item $\Delta < 0$ avec $a,b,c \in \R$ : racines complexes conjuguées.
\end{itemize}
\end{retenirbox}

%% === CHAPITRE 3 ===
\chapter{Fonctions d'une variable réelle}

\section*{Objectifs}
\begin{itemize}[leftmargin=1.5em]
  \item Maîtriser les fonctions de référence et leurs propriétés.
  \item Calculer des limites (définition, théorèmes opératoires, règle de l'Hôpital).
  \item Établir continuité, prolongement, TVI.
  \item Définir la dérivée, l'interpréter géométriquement, appliquer les règles de dérivation.
  \item Étudier parité, monotonie, convexité et tangente d'une courbe.
\end{itemize}

\section{Fonctions de référence}

Toute étude de fonction repose sur la maîtrise d'un répertoire de fonctions élémentaires dont on connaît parfaitement les propriétés. Les fonctions polynomiales, l'exponentielle, le logarithme et les fonctions trigonométriques constituent une boîte à outils complète : elles permettent de modéliser la plupart des phénomènes continus rencontrés en sciences et en ingénierie. La maîtrise de leurs domaines, dérivées et comportements asymptotiques est un prérequis indispensable pour aborder l'analyse.

\begin{definitionbox}[Fonctions usuelles]
Les fonctions de référence à maîtriser :
\begin{center}
\begin{tabular}{llll}
\toprule
Fonction & Domaine & Dérivée & Propriétés \\
\midrule
$x^n$ ($n\in\N$) & $\R$ & $nx^{n-1}$ & paire si $n$ pair, impaire si $n$ impair \\
$x^\alpha$ ($\alpha\in\R$) & $]0,+\infty[$ & $\alpha x^{\alpha-1}$ & croissante si $\alpha>0$ \\
$e^x$ & $\R$ & $e^x$ & toujours $>0$, convexe \\
$\ln x$ & $]0,+\infty[$ & $1/x$ & strictement croissante \\
$\sin x$ & $\R$ & $\cos x$ & $2\pi$-périodique, impaire \\
$\cos x$ & $\R$ & $-\sin x$ & $2\pi$-périodique, paire \\
$\tan x$ & $\R\setminus\{\pi/2+k\pi\}$ & $1+\tan^2 x = 1/\cos^2 x$ & $\pi$-périodique, impaire \\
$\arcsin x$ & $[-1,1]$ & $1/\sqrt{1-x^2}$ & image $[-\pi/2,\pi/2]$ \\
$\arccos x$ & $[-1,1]$ & $-1/\sqrt{1-x^2}$ & image $[0,\pi]$ \\
$\arctan x$ & $\R$ & $1/(1+x^2)$ & image $]-\pi/2,\pi/2[$ \\
\bottomrule
\end{tabular}
\end{center}
\end{definitionbox}

\section{Limites}

La notion de limite formalise rigoureusement l'intuition de « se rapprocher sans nécessairement atteindre ». La définition $\varepsilon$-$\delta$ de Cauchy-Weierstrass traduit cette idée : pour toute précision $\varepsilon$ arbitrairement petite, on peut trouver un voisinage de $a$ sur lequel $f$ reste à distance $\varepsilon$ de sa valeur limite. Ce formalisme est indispensable pour établir les théorèmes opératoires et justifier rigoureusement les calculs de limites rencontrés en pratique.

\begin{definitionbox}[Limite finie en un point]
Soit $f$ définie au voisinage de $a \in \R$ (sauf peut-être en $a$).
On dit que $f(x) \to L$ quand $x \to a$ si :
\[
  \forall \varepsilon > 0,\;\exists \delta > 0 : 0 < |x-a| < \delta \implies |f(x)-L| < \varepsilon.
\]
\end{definitionbox}

\begin{theoremebox}[Théorèmes opératoires sur les limites]
Si $f(x) \to L$ et $g(x) \to M$ quand $x \to a$ :
\begin{itemize}
  \item $f(x)+g(x) \to L+M$ ; $f(x)g(x) \to LM$ ; $f(x)/g(x) \to L/M$ si $M \neq 0$.
  \item \textbf{Théorème des gendarmes.} Si $g(x) \leq f(x) \leq h(x)$ et $g(x),h(x) \to L$, alors $f(x) \to L$.
  \item \textbf{Limite par composition.} Si $f(x) \to L$ et $g$ continue en $L$, alors $g(f(x)) \to g(L)$.
\end{itemize}
\end{theoremebox}

\begin{theoremebox}[Règle de l'Hôpital]
Si $f$ et $g$ sont dérivables au voisinage de $a$, $f(a)=g(a)=0$ (ou $\pm\infty$), $g'(x)\neq 0$ au voisinage de $a$, et si $\lim_{x\to a}\dfrac{f'(x)}{g'(x)} = L$, alors :
\[
  \lim_{x\to a}\frac{f(x)}{g(x)} = L.
\]
\end{theoremebox}

\section{Continuité}

La continuité formalise l'idée intuitive de « tracer la courbe sans lever le crayon » : une fonction est continue si son graphe ne présente ni saut ni trou. Le théorème des valeurs intermédiaires en est la conséquence géométrique la plus frappante : une courbe continue qui passe d'une valeur à une autre doit nécessairement traverser toutes les valeurs intermédiaires. Ce résultat, apparemment évident, repose en profondeur sur la complétude de $\mathbb{R}$ et fournit un outil fondamental pour prouver l'existence de solutions d'équations.

\begin{definitionbox}[Continuité en un point]
$f$ est \emph{continue en $a$} si $\lim_{x\to a} f(x) = f(a)$.
$f$ est \emph{continue sur $I$} si elle est continue en tout point de $I$.
\end{definitionbox}

\begin{theoremebox}[Théorème des valeurs intermédiaires (TVI)]
Soit $f : [a,b] \to \R$ continue. Pour tout $k$ compris entre $f(a)$ et $f(b)$, il existe $c \in [a,b]$ tel que $f(c) = k$.
\end{theoremebox}

\begin{preuvebox}
On traite le cas $f(a) < k < f(b)$. Posons $A = \{x \in [a,b] \mid f(x) \leq k\}$.
$A$ est non vide ($a \in A$) et majoré par $b$, donc $c = \sup A$ existe.
Par continuité de $f$ en $c$, $f(c) = k$ (sinon on obtiendrait une contradiction sur la définition de sup). \qed
\end{preuvebox}

\begin{theoremebox}[Corollaire — existence d'une racine]
Si $f$ est continue sur $[a,b]$ et $f(a)f(b) < 0$, alors $f$ admet au moins une racine dans $]a,b[$.
\end{theoremebox}

\begin{definitionbox}[Prolongement par continuité]
Si $f$ est définie sur $I \setminus \{a\}$ et $\lim_{x\to a} f(x) = L$, on peut prolonger $f$ en une fonction continue sur $I$ en posant $f(a) = L$.
\end{definitionbox}

\section{Dérivabilité}

La dérivée est le concept central de l'analyse différentielle : elle quantifie le taux de variation instantané d'une fonction et fournit la meilleure approximation affine (linéaire) de la fonction au voisinage d'un point. Géométriquement, $f'(a)$ est le coefficient directeur de la tangente à la courbe en $(a, f(a))$ ; analytiquement, elle permet de contrôler la monotonie et de localiser les extrema. Le théorème des accroissements finis, qui en est la conséquence la plus utile, relie la variation globale de $f$ sur $[a,b]$ à une valeur locale de la dérivée.

\begin{definitionbox}[Dérivée en un point]
$f$ est \emph{dérivable en $a$} si la limite suivante existe et est finie :
\[
  f'(a) = \lim_{h\to 0}\frac{f(a+h)-f(a)}{h} = \lim_{x\to a}\frac{f(x)-f(a)}{x-a}.
\]
$f'(a)$ est la \emph{dérivée} de $f$ en $a$. Géométriquement, c'est le coefficient directeur de la tangente à la courbe en $(a,f(a))$.
\end{definitionbox}

\begin{theoremebox}[Règles de dérivation]
\begin{align*}
  (f+g)' &= f'+g', & (\lambda f)' &= \lambda f', \\
  (fg)' &= f'g+fg', & \left(\frac{f}{g}\right)' &= \frac{f'g-fg'}{g^2}\;(g\neq0), \\
  (g\circ f)'(x) &= g'(f(x))\cdot f'(x) & &\text{(règle de la chaîne)}.
\end{align*}
\end{theoremebox}

\begin{theoremebox}[Théorème des accroissements finis (TAF)]
Si $f$ est continue sur $[a,b]$ et dérivable sur $]a,b[$, il existe $c \in ]a,b[$ tel que :
\[
  f(b)-f(a) = f'(c)(b-a).
\]
\end{theoremebox}

\begin{preuvebox}
Appliquer le théorème de Rolle à $g(x) = f(x) - f(a) - \dfrac{f(b)-f(a)}{b-a}(x-a)$,
qui vérifie $g(a)=g(b)=0$. \qed
\end{preuvebox}

\section{Parité, monotonie et convexité}

L'étude qualitative d'une fonction s'appuie sur ses propriétés globales : la parité réduit le domaine d'étude par symétrie, la périodicité ramène l'étude à une période, la monotonie décrit l'évolution générale et la convexité précise la courbure. Ces propriétés structurent l'analyse d'une fonction bien avant tout calcul numérique, et permettent de dresser un tableau de variations complet à partir du signe de $f'$ et $f''$.

\begin{definitionbox}[Parité et périodicité]
\begin{itemize}
  \item $f$ est \emph{paire} si $f(-x) = f(x)$ pour tout $x$ du domaine (courbe symétrique par rapport à l'axe des ordonnées).
  \item $f$ est \emph{impaire} si $f(-x) = -f(x)$ (courbe symétrique par rapport à l'origine).
  \item $f$ est \emph{$T$-périodique} si $f(x+T) = f(x)$ pour tout $x$.
\end{itemize}
\end{definitionbox}

\begin{theoremebox}[Monotonie et signe de la dérivée]
Soit $f$ dérivable sur $I$.
\begin{itemize}
  \item $f$ croissante sur $I$ $\iff$ $f'(x) \geq 0$ sur $I$.
  \item $f$ strictement croissante sur $I$ $\impliedby$ $f'(x) > 0$ sur $I$ (sauf en un nombre fini de points).
  \item $f$ constante sur $I$ $\iff$ $f'(x) = 0$ sur $I$.
\end{itemize}
\end{theoremebox}

\begin{definitionbox}[Convexité]
$f$ est \emph{convexe} sur $I$ si pour tous $x,y \in I$ et $t \in [0,1]$ :
\[
  f(tx+(1-t)y) \leq t f(x)+(1-t)f(y).
\]
Géométriquement : la corde est au-dessus de la courbe.
$f$ est \emph{concave} si $-f$ est convexe.
\end{definitionbox}

\begin{theoremebox}[Convexité et dérivée seconde]
Soit $f$ deux fois dérivable sur $I$.
\begin{itemize}
  \item $f$ convexe sur $I$ $\iff$ $f''(x) \geq 0$ pour tout $x \in I$.
  \item $f$ concave sur $I$ $\iff$ $f''(x) \leq 0$ pour tout $x \in I$.
  \item Un \emph{point d'inflexion} est un point où $f''$ change de signe.
\end{itemize}
\end{theoremebox}

\begin{center}
\begin{tikzpicture}
\begin{axis}[
  width=10cm, height=5.5cm,
  xlabel={$x$}, ylabel={$y$},
  title={\small Convexité, tangente et point d'inflexion},
  domain=-2.5:2.5, samples=100,
  xmin=-2.5, xmax=2.5, ymin=-2, ymax=5,
  grid=major, grid style={dotted, gray!40},
  tick label style={font=\scriptsize},
  label style={font=\small},
  legend style={font=\scriptsize, at={(0.98,0.98)}, anchor=north east},
]
% Courbe f(x) = x^3 - x
\addplot[thick, black] {x^3 - x};
\addlegendentry{$f(x)=x^3-x$}
% Tangente en x=0 (inflexion) : f'(0)=-1, f(0)=0 => y=-x
\addplot[dashed, caplpred, domain=-1.5:1.5] {-x};
\addlegendentry{Tangente en $0$}
% Point d'inflexion
\addplot[mark=*, mark size=3pt, caplpred] coordinates {(0,0)};
% Zone convexe (x>0) et concave (x<0)
\node[font=\scriptsize, caplpgreen] at (axis cs:1.8,3.5) {convexe ($f''>0$)};
\node[font=\scriptsize, caplpblue] at (axis cs:-1.8,3.5) {concave ($f''<0$)};
\end{axis}
\end{tikzpicture}
\captionof*{figure}{\small\itshape $f(x)=x^3-x$ : point d'inflexion en $x=0$ où $f''$ change de signe ; la tangente traverse la courbe.}
\end{center}

\section{Tangente et équation de la tangente}

La tangente en un point est la meilleure approximation affine de la courbe au voisinage de ce point : parmi toutes les droites passant par $(a, f(a))$, c'est celle qui s'écarte le moins vite de la courbe quand $x \to a$. Cette idée est au cœur du calcul différentiel et préfigure directement la notion de différentielle pour les fonctions de plusieurs variables.

\begin{definitionbox}[Équation de la tangente]
Si $f$ est dérivable en $a$, la tangente à la courbe de $f$ au point $(a, f(a))$ a pour équation :
\[
  y = f(a) + f'(a)(x-a).
\]
\end{definitionbox}

\section{Exercices types}

\begin{exercicebox}[Limite et règle de l'Hôpital]
Calculer les limites suivantes :
\begin{enumerate}
  \item $\displaystyle\lim_{x\to 0}\frac{\sin x - x}{x^3}$
  \item $\displaystyle\lim_{x\to +\infty} x\ln\!\left(1+\frac{1}{x}\right)$
  \item $\displaystyle\lim_{x\to 0^+} x^\alpha \ln x$ pour $\alpha > 0$
\end{enumerate}
\end{exercicebox}

\begin{correctionbox}[Correction]
\textbf{1.} Forme $0/0$. On applique l'Hôpital deux fois :
$\dfrac{\sin x - x}{x^3} \xrightarrow{H} \dfrac{\cos x - 1}{3x^2} \xrightarrow{H} \dfrac{-\sin x}{6x} \to -\dfrac{1}{6}$.

\textbf{2.} Posons $u = 1/x \to 0^+$ : $\dfrac{\ln(1+u)}{u} \to 1$. Donc la limite est $1$.

\textbf{3.} $x^\alpha \ln x = \dfrac{\ln x}{x^{-\alpha}}$. Hôpital : $\dfrac{1/x}{-\alpha x^{-\alpha-1}} = -\dfrac{x^\alpha}{\alpha} \to 0$.
\end{correctionbox}

\begin{exercicebox}[Étude complète]
Étudier la fonction $f(x) = x e^{-x}$ : domaine, limites, variations, convexité, tangente en $x=1$.
\end{exercicebox}

\begin{correctionbox}[Correction]
\textbf{Domaine :} $\R$. \textbf{Limites :} $f(x)\to 0$ en $\pm\infty$ ($xe^{-x}\to 0$ en $+\infty$ par croissances comparées ; $\to -\infty$ en $-\infty$).

\textbf{Dérivée :} $f'(x) = e^{-x} - xe^{-x} = (1-x)e^{-x}$.
$f'(x) > 0 \iff x < 1$ : $f$ croissante sur $]-\infty,1[$, décroissante sur $]1,+\infty[$. Maximum en $x=1$ : $f(1)=e^{-1}$.

\textbf{Dérivée seconde :} $f''(x) = -(1-x)e^{-x} + (-1)e^{-x} = (x-2)e^{-x}$.
$f''(x) > 0 \iff x > 2$ : convexe sur $]2,+\infty[$, concave sur $]-\infty,2[$. Point d'inflexion en $x=2$, $f(2)=2e^{-2}$.

\textbf{Tangente en $x=1$ :} $f(1)=e^{-1}$, $f'(1)=0$. Équation : $y = e^{-1}$.
\end{correctionbox}

\begin{exercicebox}[TVI — type concours]
Montrer que l'équation $x^3 + x - 1 = 0$ admet une unique solution réelle dans $[0,1]$, puis l'encadrer à $10^{-1}$ près.
\end{exercicebox}

\begin{correctionbox}[Correction]
Posons $f(x)=x^3+x-1$. $f$ est continue sur $[0,1]$.
$f(0)=-1 < 0$ et $f(1)=1 > 0$ : par le TVI, $\exists c \in ]0,1[$ tel que $f(c)=0$.

\textbf{Unicité.} $f'(x)=3x^2+1 > 0$ sur $\R$ : $f$ strictement croissante, donc la racine est unique.

\textbf{Encadrement.} $f(0{,}6)=0{,}216+0{,}6-1=-0{,}184 < 0$ et $f(0{,}7)=0{,}343+0{,}7-1=0{,}043 > 0$.
Donc $c \in ]0{,}6 ; 0{,}7[$, soit $c \approx 0{,}68$ à $10^{-1}$ près.
\end{correctionbox}

\begin{retenirbox}
\textbf{Fonctions d'une variable réelle — À retenir}
\begin{itemize}
  \item Dérivée = pente de la tangente. Tangente : $y=f(a)+f'(a)(x-a)$.
  \item $f'>0 \Rightarrow$ croissante ; $f''> 0 \Rightarrow$ convexe (corde au-dessus).
  \item TAF : $f(b)-f(a)=f'(c)(b-a)$, $c \in ]a,b[$.
  \item TVI : $f$ continue, $f(a)f(b)<0 \Rightarrow$ racine dans $]a,b[$.
  \item Hôpital : forme $0/0$ ou $\infty/\infty$ : $\lim f/g = \lim f'/g'$.
  \item $xe^{-x}\to 0$ ($x\to+\infty$) ; $x^\alpha\ln x\to 0$ ($x\to 0^+$, $\alpha>0$) : croissances comparées.
\end{itemize}
\end{retenirbox}

%% === CHAPITRE 4 ===
\chapter{Courbes paramétrées}

\section*{Objectifs}
\begin{itemize}[leftmargin=1.5em]
  \item Définir et étudier une courbe paramétrée plane.
  \item Déterminer les éléments de symétrie, les points singuliers, les tangentes.
  \item Tracer les courbes polynomiales, rationnelles et trigonométriques.
  \item Appliquer le schéma global d'étude.
\end{itemize}

\section{Définitions fondamentales}

Les courbes paramétrées permettent de représenter des trajectoires qui ne sont pas graphes de fonctions : un cercle, une cycloïde ou une spirale ne peuvent pas s'écrire sous la forme $y = f(x)$ sur tout leur support. En introduisant un paramètre $t$ (souvent interprété comme le temps), on décrit le mouvement point par point, ce qui ouvre la voie à l'étude des tangentes, des vitesses et des singularités indépendamment de toute contrainte de fonctionnalité.

\begin{definitionbox}[Courbe paramétrée]
Une \emph{courbe paramétrée} est la donnée d'un arc $\mathcal{C}$ défini par :
\[
  \mathcal{C} : \begin{cases} x = f(t) \\ y = g(t) \end{cases}, \quad t \in I \subseteq \R,
\]
où $f$ et $g$ sont des fonctions de classe $\mathcal{C}^1$ sur $I$.
Le point $M(t) = (f(t),g(t))$ est le point courant.
Le vecteur $\overrightarrow{M'(t)} = (f'(t), g'(t))$ est le \emph{vecteur dérivé}.
\end{definitionbox}

\begin{definitionbox}[Point régulier, point singulier]
\begin{itemize}
  \item $M(t)$ est \emph{régulier} si $\overrightarrow{M'(t)} \neq \vec{0}$, i.e.\ $f'(t)^2+g'(t)^2 \neq 0$.
  \item $M(t)$ est \emph{singulier} (ou stationnaire) si $f'(t) = g'(t) = 0$.
\end{itemize}
En un point régulier, la tangente à $\mathcal{C}$ en $M(t)$ est dirigée par $\overrightarrow{M'(t)}$.
\end{definitionbox}

\begin{definitionbox}[Tangente en un point singulier]
Si $M(t_0)$ est singulier, on cherche la limite de la direction de $\overrightarrow{M'(t)}$ quand $t \to t_0$.
Plus précisément, on développe $f(t_0+h)$ et $g(t_0+h)$ en $h$ :
\[
  f(t_0+h) \sim a_p h^p, \quad g(t_0+h) \sim b_q h^q \quad (h \to 0),
\]
et la tangente est dirigée par $(a_p, 0)$ si $p < q$, ou $(0, b_q)$ si $q < p$, ou $(a_p, b_p)$ si $p = q$.
\end{definitionbox}

\section{Symétries d'une courbe paramétrée}

Exploiter les symétries d'une courbe paramétrée permet de réduire considérablement l'intervalle d'étude effectif. Lorsque l'on peut établir qu'un arc sur $[t_0, t_1]$ se déduit d'un arc sur un sous-intervalle par une symétrie axiale ou centrale, il suffit d'étudier ce sous-intervalle et de compléter le tracé par réflexion. Cette économie de calcul est systématiquement recherchée dans les épreuves de concours.

\begin{methodebox}[Détection des symétries]
\begin{enumerate}
  \item \textbf{Symétrie par rapport à l'axe $Ox$} : si $\exists \varphi$ tel que $f(\varphi(t)) = f(t)$ et $g(\varphi(t)) = -g(t)$.
  \item \textbf{Symétrie par rapport à l'axe $Oy$} : si $f(\varphi(t)) = -f(t)$ et $g(\varphi(t)) = g(t)$.
  \item \textbf{Symétrie par rapport à l'origine} : si $f(\varphi(t)) = -f(t)$ et $g(\varphi(t)) = -g(t)$.
  \item \textbf{Périodicité} : si $(f(t+T), g(t+T)) = (f(t), g(t))$, étudier sur une période.
\end{enumerate}
\end{methodebox}

\section{Schéma global d'étude}

L'étude d'une courbe paramétrée suit un protocole rigoureux et systématique qui garantit de n'omettre aucune caractéristique géométrique importante : domaine, réduction par symétries, tableau de variations simultané de $x(t)$ et $y(t)$, traitement des points singuliers et des branches infinies. La rigueur de ce schéma est attendue dans toute réponse de concours.

\begin{methodebox}[Schéma d'étude d'une courbe paramétrée]
\begin{enumerate}
  \item \textbf{Domaine} : intervalle de définition de $f$ et $g$.
  \item \textbf{Réduction} : utiliser périodicité et symétries pour se ramener à un intervalle minimal.
  \item \textbf{Tableau de variations} : variations de $x(t)=f(t)$ et $y(t)=g(t)$ ; signe de $f'$ et $g'$.
  \item \textbf{Points singuliers} : chercher $f'(t)=g'(t)=0$ et étudier la tangente par développement limité.
  \item \textbf{Branches infinies} : étudier le comportement quand $t \to \pm\infty$ ou aux bords du domaine.
  \item \textbf{Tracé} : calculer quelques points remarquables et tracer la courbe.
\end{enumerate}
\end{methodebox}

\section{Exemples de courbes classiques}

Le répertoire des courbes paramétrées classiques — cycloïde, courbes de Lissajous, astroïde, cardioïde, spirale d'Archimède — constitue un patrimoine géométrique riche, illustrant la diversité des formes que peut engendrer une simple dépendance paramétrique. Chacune de ces courbes possède une signification mécanique ou physique précise, et leur étude concrète permet de fixer les méthodes.

\begin{exemplebox}[Cycloïde (courbe polynomiale trigonométrique)]
$\mathcal{C} : x(t) = t - \sin t$, $y(t) = 1 - \cos t$, $t \in \R$.

$x'(t) = 1-\cos t \geq 0$, $y'(t) = \sin t$.
Points singuliers : $x'(t)=0 \iff t = 2k\pi$. En $t=0$ : $x(0)=0$, $y(0)=0$.
Développement : $x(h) \sim h^2/2$, $y(h) \sim h^2/2$, tangente verticale … en réalité horizontale car $y/x \to 1$.
La cycloïde est la trajectoire d'un point d'un cercle de rayon 1 roulant sur l'axe $Ox$.
\end{exemplebox}

\begin{exemplebox}[Courbe de Lissajous (trigonométrique)]
$\mathcal{C} : x(t) = \cos(2t)$, $y(t) = \sin(3t)$, $t \in [0, 2\pi]$.
Période : $\mathrm{ppcm}(\pi, 2\pi/3) = 2\pi$, on étudie sur $[0, 2\pi]$.
$|x| \leq 1$, $|y| \leq 1$ : courbe bornée dans $[-1,1]^2$.
Points singuliers : $x'(t) = -2\sin(2t) = 0 \iff t = k\pi/2$ ; $y'(t) = 3\cos(3t) = 0 \iff t = \pi/6+k\pi/3$.
Aux points réguliers : tangente dirigée par $(-2\sin(2t), 3\cos(3t))$.
\end{exemplebox}

\begin{center}
\begin{tikzpicture}
\begin{axis}[
  width=7cm, height=5.5cm,
  xlabel={$x$}, ylabel={$y$},
  title={\small Courbe de Lissajous : $x=\cos(2t)$, $y=\sin(3t)$},
  xmin=-1.2, xmax=1.2, ymin=-1.2, ymax=1.2,
  axis equal,
  grid=major, grid style={dotted, gray!30},
  tick label style={font=\scriptsize},
  label style={font=\small},
]
\addplot[thick, caplpblue, domain=0:360, samples=200]
  ({cos(2*x)}, {sin(3*x)});
\end{axis}
\end{tikzpicture}
\quad
\begin{tikzpicture}
\begin{axis}[
  width=7cm, height=5.5cm,
  xlabel={$x$}, ylabel={$y$},
  title={\small Astroïde : $x=\cos^3 t$, $y=\sin^3 t$},
  xmin=-1.2, xmax=1.2, ymin=-1.2, ymax=1.2,
  axis equal,
  grid=major, grid style={dotted, gray!30},
  tick label style={font=\scriptsize},
  label style={font=\small},
]
\addplot[thick, caplpred, domain=0:360, samples=200]
  ({(cos(x))^3}, {(sin(x))^3});
\end{axis}
\end{tikzpicture}
\captionof*{figure}{\small\itshape Courbe de Lissajous (gauche) et astroïde (droite) — deux courbes paramétrées trigonométriques classiques.}
\end{center}

\begin{exercicebox}[Étude d'une courbe paramétrée]
Étudier et tracer la courbe $\mathcal{C}$ définie par $x(t) = t^2 - 1$, $y(t) = t^3 - t$, $t \in \R$.
\end{exercicebox}

\begin{correctionbox}[Correction]
\textbf{Symétrie.} $x(-t)=x(t)$, $y(-t)=-y(t)$ : symétrie par rapport à $Ox$. On étudie sur $[0,+\infty[$.

\textbf{Variations.} $x'(t)=2t\geq 0$ (croissant), $y'(t)=3t^2-1$.
$y'(t)=0 \iff t=1/\sqrt{3}$.

\textbf{Tableau :}
\begin{center}
\begin{tabular}{c|ccccc}
$t$ & $0$ & & $1/\sqrt{3}$ & & $+\infty$ \\
\hline
$x'$ & $0$ & $+$ & & $+$ & \\
$x$ & $-1$ & $\nearrow$ & & $\nearrow$ & $+\infty$ \\
\hline
$y'$ & $-1$ & $-$ & $0$ & $+$ & \\
$y$ & $0$ & $\searrow$ & $-2/(3\sqrt{3})$ & $\nearrow$ & $+\infty$ \\
\end{tabular}
\end{center}

\textbf{Point singulier.} $x'(0)=0$, $y'(0)=-1\neq 0$ : point régulier, tangente verticale en $M(0)=(-1,0)$.

\textbf{Point double.} $M(t_1)=M(t_2)$ avec $t_1\neq t_2$ : $t_1^2=t_2^2$ et $t_1^3-t_1=t_2^3-t_2$.
$t_2=-t_1$, donc par symétrie : point double en $x=-1+t^2$... En $t=1$ et $t=-1$ : $M(1)=(0,0)=M(-1)$. Point double en l'origine.
\end{correctionbox}

\begin{retenirbox}
\textbf{Courbes paramétrées — À retenir}
\begin{itemize}
  \item Vecteur dérivé $\vec{M'(t)}=(f'(t),g'(t))$ dirige la tangente (si $\neq\vec{0}$).
  \item Point singulier : $f'=g'=0$ — étudier via DL.
  \item Exploiter les symétries pour réduire l'intervalle d'étude.
  \item Chercher les points doubles : $M(t_1)=M(t_2)$, $t_1\neq t_2$.
  \item Branches infinies : étudier $y/x$ ou $y-ax-b$ quand $t\to$ bord.
\end{itemize}
\end{retenirbox}

%% === CHAPITRE 5 ===
\chapter{Calcul intégral}

\section*{Objectifs}
\begin{itemize}[leftmargin=1.5em]
  \item Définir l'intégrale de Riemann et connaître ses propriétés.
  \item Calculer des primitives des fonctions usuelles.
  \item Appliquer les techniques : intégration par parties (IPP), changement de variable.
  \item Calculer des aires et des valeurs moyennes.
  \item Étudier les intégrales impropres (convergence).
\end{itemize}

\section{Intégrale de Riemann}

L'intégrale de Riemann fournit une construction rigoureuse de l'aire sous une courbe à partir de sommes de rectangles dont les bases tendent vers zéro. Cette construction par sommes de Riemann est fondamentale car elle justifie les calculs numériques d'intégrales et donne un sens précis à la notion d'aire algébrique, tenant compte du signe de la fonction. La convergence de ces sommes vers une limite finie repose sur la régularité (continuité ou continuité par morceaux) de la fonction intégrée.

\begin{definitionbox}[Intégrale de Riemann]
Soit $f : [a,b] \to \R$ continue (ou par morceaux). L'\emph{intégrale de Riemann} de $f$ sur $[a,b]$ est :
\[
  \int_a^b f(x)\,\mathrm{d}x = \lim_{n\to+\infty} \sum_{k=0}^{n-1} f\!\left(a+k\,\frac{b-a}{n}\right)\frac{b-a}{n}.
\]
C'est l'aire algébrique entre la courbe $y=f(x)$ et l'axe $Ox$ sur $[a,b]$.
\end{definitionbox}

\begin{theoremebox}[Propriétés de l'intégrale]
Pour $f, g$ continues sur $[a,b]$ et $\lambda \in \R$ :
\begin{align*}
  &\int_a^b (\lambda f+g)\,\mathrm{d}x = \lambda\int_a^b f\,\mathrm{d}x + \int_a^b g\,\mathrm{d}x & &\text{(linéarité)} \\
  &\int_a^b f\,\mathrm{d}x = \int_a^c f\,\mathrm{d}x + \int_c^b f\,\mathrm{d}x & &\text{(relation de Chasles)} \\
  &f \geq 0 \text{ sur } [a,b] \implies \int_a^b f\,\mathrm{d}x \geq 0 & &\text{(positivité)} \\
  &\left|\int_a^b f\,\mathrm{d}x\right| \leq \int_a^b |f|\,\mathrm{d}x & &\text{(inégalité triangulaire)}
\end{align*}
Par convention : $\int_b^a f\,\mathrm{d}x = -\int_a^b f\,\mathrm{d}x$.
\end{theoremebox}

\section{Théorème fondamental de l'analyse}

Le théorème fondamental de l'analyse est le résultat le plus important du calcul intégral : il établit que la dérivation et l'intégration sont des opérations inverses l'une de l'autre. Ce lien profond entre les deux grandes branches de l'analyse — le calcul différentiel et le calcul intégral — permet de calculer des intégrales en cherchant des primitives, transformant un problème de passage à la limite en un simple calcul algébrique.

\begin{theoremebox}[Théorème fondamental]
Soit $f$ continue sur $[a,b]$. La fonction $F : x \mapsto \displaystyle\int_a^x f(t)\,\mathrm{d}t$ est l'unique primitive de $f$ sur $[a,b]$ s'annulant en $a$.

De plus, si $G$ est une primitive quelconque de $f$ sur $[a,b]$ :
\[
  \int_a^b f(x)\,\mathrm{d}x = G(b) - G(a) = \big[G(x)\big]_a^b.
\]
\end{theoremebox}

\begin{preuvebox}
$F'(x) = \lim_{h\to0}\dfrac{1}{h}\displaystyle\int_x^{x+h}f(t)\,\mathrm{d}t$.
Par continuité de $f$, pour $h$ petit, $\displaystyle\int_x^{x+h}f(t)\,\mathrm{d}t \approx f(x)\cdot h$, d'où $F'(x)=f(x)$. \qed
\end{preuvebox}

\section{Primitives usuelles}

La stratégie générale pour calculer une intégrale est de reconnaître la forme de la fonction intégrée dans la table des primitives usuelles, ou de la ramener à cette table par une technique d'intégration. La maîtrise de cette liste est donc un préalable indispensable : elle constitue le socle à partir duquel toutes les techniques d'intégration s'appliquent.

\begin{retenirbox}
Primitives à connaître par cœur ($C \in \R$ constante d'intégration) :
\[
\begin{array}{ll@{\qquad}ll}
x^n \;(n\neq-1) & \dfrac{x^{n+1}}{n+1}+C &
e^x & e^x+C \\[8pt]
\dfrac{1}{x} & \ln|x|+C &
\ln x & x\ln x - x + C \\[8pt]
\cos x & \sin x+C &
\sin x & -\cos x+C \\[8pt]
\dfrac{1}{\cos^2 x} & \tan x+C &
\dfrac{1}{1+x^2} & \arctan x+C \\[8pt]
\dfrac{1}{\sqrt{1-x^2}} & \arcsin x+C &
a^x \;(a>0) & \dfrac{a^x}{\ln a}+C
\end{array}
\]
\end{retenirbox}

\section{Techniques d'intégration}

Lorsque la primitive d'une fonction n'est pas immédiatement lisible dans la table, deux outils fondamentaux permettent de transformer l'intégrale : l'intégration par parties, qui joue le rôle de la règle de Leibniz pour les primitives, et le changement de variable, qui transpose le problème dans un nouveau paramètre plus maniable. Le choix judicieux de $u$ et $v'$ dans l'IPP, ou du changement $x = \varphi(t)$, est au cœur de la résolution de la majorité des intégrales rencontrées en concours.

\begin{theoremebox}[Intégration par parties (IPP)]
Soient $u, v$ de classe $\mathcal{C}^1$ sur $[a,b]$. Alors :
\[
  \int_a^b u(x)v'(x)\,\mathrm{d}x = \big[u(x)v(x)\big]_a^b - \int_a^b u'(x)v(x)\,\mathrm{d}x.
\]
\end{theoremebox}

\begin{methodebox}[Choix de $u$ et $v'$ en IPP]
Règle mnémotechnique LIATE (par ordre de priorité pour $u$) :
Logarithme · Inverse · Algébrique (polynôme) · Trigonométrique · Exponentielle.
Choisir $v'$ de sorte que $v$ soit simple à calculer.
\end{methodebox}

\begin{theoremebox}[Changement de variable]
Si $\varphi : [\alpha,\beta] \to [a,b]$ est de classe $\mathcal{C}^1$, $\varphi(\alpha)=a$, $\varphi(\beta)=b$ :
\[
  \int_a^b f(x)\,\mathrm{d}x = \int_\alpha^\beta f(\varphi(t))\,\varphi'(t)\,\mathrm{d}t.
\]
\end{theoremebox}

\begin{exemplebox}[IPP et changement de variable]
\textbf{IPP :} $\displaystyle I = \int_0^1 x e^x\,\mathrm{d}x$.
$u=x$, $v'=e^x$ : $I = [xe^x]_0^1 - \int_0^1 e^x\,\mathrm{d}x = e - [e^x]_0^1 = e-(e-1) = 1$.

\textbf{Changement de variable :} $\displaystyle J = \int_0^1 \frac{x}{\sqrt{1-x^2}}\,\mathrm{d}x$.
Posons $t = 1-x^2$, $\mathrm{d}t = -2x\,\mathrm{d}x$. $x:0\to1$ donne $t:1\to0$.
$J = \int_1^0 \dfrac{-\mathrm{d}t}{2\sqrt{t}} = \dfrac{1}{2}\int_0^1 t^{-1/2}\,\mathrm{d}t = \big[t^{1/2}\big]_0^1 = 1$.
\end{exemplebox}

\section{Aires et valeur moyenne}

L'intégrale possède deux interprétations géométrique et physique immédiatement opérationnelles. D'une part, elle mesure l'aire algébrique entre des courbes, ce qui est utile en probabilités (calcul de probabilités pour les lois continues) et en mécanique (travail d'une force). D'autre part, la valeur moyenne d'une fonction sur un intervalle généralise la moyenne arithmétique à une infinité de valeurs : c'est la « hauteur » du rectangle de même aire que l'intégrale.

\begin{definitionbox}[Aire entre deux courbes]
L'aire du domaine compris entre les courbes $y=f(x)$ et $y=g(x)$ sur $[a,b]$ est :
\[
  \mathcal{A} = \int_a^b |f(x)-g(x)|\,\mathrm{d}x.
\]
\end{definitionbox}

\begin{definitionbox}[Valeur moyenne]
La \emph{valeur moyenne} de $f$ continue sur $[a,b]$ est :
\[
  \langle f \rangle = \frac{1}{b-a}\int_a^b f(x)\,\mathrm{d}x.
\]
\end{definitionbox}

\begin{center}
\begin{tikzpicture}
\begin{axis}[
  width=10cm, height=5.5cm,
  xlabel={$x$}, ylabel={$y$},
  title={\small Aire entre $f(x)=x^2$ et $g(x)=x$ sur $[0,1]$},
  domain=0:1, samples=80,
  xmin=-0.1, xmax=1.2, ymin=-0.1, ymax=1.2,
  grid=major, grid style={dotted, gray!30},
  tick label style={font=\scriptsize},
  label style={font=\small},
]
\addplot[name path=F, thick, caplpblue] {x};
\addplot[name path=G, thick, caplpred] {x^2};
\addplot[caplpgreenbg, opacity=0.6] fill between[of=F and G, soft clip={domain=0:1}];
\addlegendimage{area legend, fill=caplpgreenbg}
\addlegendentry{Aire $= \int_0^1(x-x^2)\,dx = 1/6$}
\addlegendentry{$g(x)=x$}
\addlegendentry{$f(x)=x^2$}
\end{axis}
\end{tikzpicture}
\captionof*{figure}{\small\itshape Aire entre $y=x$ et $y=x^2$ sur $[0,1]$ : $\int_0^1(x-x^2)\,\mathrm{d}x = \left[\frac{x^2}{2}-\frac{x^3}{3}\right]_0^1 = \frac{1}{6}$.}
\end{center}

\section{Intégrales impropres}

Les intégrales impropres étendent la définition riemannienne aux cas où l'intervalle d'intégration est non borné ou la fonction présente une singularité à la borne. La convergence n'est alors plus automatique et doit être établie par passage à la limite. Les intégrales de référence — série de Bertrand $\int_1^{+\infty} x^{-\alpha}\,\mathrm{d}x$ et leur analogue en $0$ — servent de pierre de touche pour les critères de comparaison.

\begin{definitionbox}[Intégrale impropre]
\begin{itemize}
  \item \textbf{Borne infinie :} $\displaystyle\int_a^{+\infty} f(x)\,\mathrm{d}x = \lim_{A\to+\infty}\int_a^A f(x)\,\mathrm{d}x$ (si la limite existe).
  \item \textbf{Singularité en $b$ :} $\displaystyle\int_a^b f(x)\,\mathrm{d}x = \lim_{\varepsilon\to0^+}\int_a^{b-\varepsilon} f(x)\,\mathrm{d}x$.
\end{itemize}
\end{definitionbox}

\begin{theoremebox}[Intégrales de référence]
\begin{itemize}
  \item $\displaystyle\int_1^{+\infty}\frac{\mathrm{d}x}{x^\alpha}$ converge $\iff \alpha > 1$.
  \item $\displaystyle\int_0^1 \frac{\mathrm{d}x}{x^\alpha}$ converge $\iff \alpha < 1$.
  \item $\displaystyle\int_0^{+\infty} e^{-x}\,\mathrm{d}x = 1$.
\end{itemize}
\end{theoremebox}

\begin{exercicebox}[Calculs d'intégrales]
\begin{enumerate}
  \item $\displaystyle\int_1^e \ln x\,\mathrm{d}x$
  \item $\displaystyle\int_0^{\pi/2}\sin^2 x\,\mathrm{d}x$
  \item $\displaystyle\int_0^{+\infty} x e^{-x}\,\mathrm{d}x$
\end{enumerate}
\end{exercicebox}

\begin{correctionbox}[Correction]
\textbf{1.} IPP avec $u=\ln x$, $v'=1$ :
$\displaystyle\int_1^e\ln x\,\mathrm{d}x = [x\ln x]_1^e - \int_1^e 1\,\mathrm{d}x = e - (e-1) = 1$.

\textbf{2.} Linéarisation : $\sin^2 x = \dfrac{1-\cos(2x)}{2}$.
$\displaystyle\int_0^{\pi/2}\sin^2 x\,\mathrm{d}x = \left[\frac{x}{2}-\frac{\sin(2x)}{4}\right]_0^{\pi/2} = \frac{\pi}{4}$.

\textbf{3.} IPP : $[{-xe^{-x}}]_0^{+\infty}+\int_0^{+\infty}e^{-x}\,\mathrm{d}x = 0 + 1 = 1$.
(La limite $xe^{-x}\to0$ en $+\infty$ par croissances comparées.)
\end{correctionbox}

\begin{retenirbox}
\textbf{Calcul intégral — À retenir}
\begin{itemize}
  \item TFA : $\left(\int_a^x f\right)' = f(x)$ ; $\int_a^b f = [F]_a^b$.
  \item IPP : $\int uv' = [uv] - \int u'v$.
  \item Changement de variable : $\mathrm{d}x \to \varphi'(t)\,\mathrm{d}t$, adapter les bornes.
  \item Aire entre courbes : $\int_a^b|f-g|$.
  \item Valeur moyenne : $\frac{1}{b-a}\int_a^b f$.
  \item Intégrale de Riemann $\int_1^{+\infty}x^{-\alpha}$ : CV $\iff \alpha>1$.
\end{itemize}
\end{retenirbox}

%% === CHAPITRE 6 ===
\chapter{Équations différentielles}

\section*{Objectifs}
\begin{itemize}[leftmargin=1.5em]
  \item Résoudre les équations différentielles du premier ordre linéaires à coefficients constants (homogène + variation de la constante).
  \item Résoudre les équations différentielles du premier ordre séparables.
  \item Résoudre les équations différentielles du second ordre à coefficients constants.
  \item Trouver une solution particulière par identification (méthode des coefficients indéterminés).
\end{itemize}

\section{Équations différentielles du premier ordre}

Les équations différentielles sont le langage naturel de la modélisation des phénomènes d'évolution continue : croissance d'une population, refroidissement d'un corps, charge d'un condensateur. Une ED du premier ordre relie la dérivée de l'inconnue à la fonction elle-même ; sa résolution détermine l'état futur d'un système à partir de son état initial. La superposition de la solution homogène et d'une solution particulière constitue la structure fondamentale de l'ensemble des solutions.

\begin{definitionbox}[ED1 linéaire à coefficients constants]
Une \emph{équation différentielle du premier ordre linéaire à coefficients constants} est de la forme :
\[
  y' + ay = b(x), \quad a \in \R,\; b : I \to \R \text{ continue}.
\]
L'\emph{équation homogène associée} est $y' + ay = 0$.
\end{definitionbox}

\begin{theoremebox}[Résolution de l'homogène]
Les solutions de $y' + ay = 0$ sont :
\[
  y_h(x) = C e^{-ax}, \quad C \in \R.
\]
\end{theoremebox}

\begin{preuvebox}
$y' = -ay \iff \frac{y'}{y} = -a \iff (\ln|y|)' = -a \iff \ln|y| = -ax + \mathrm{cte}$, d'où $y = Ce^{-ax}$. \qed
\end{preuvebox}

\begin{theoremebox}[Structure de l'ensemble des solutions]
La solution générale de $y' + ay = b(x)$ est :
\[
  y(x) = y_h(x) + y_p(x) = Ce^{-ax} + y_p(x),
\]
où $y_p$ est une \emph{solution particulière} quelconque de l'équation complète.
\end{theoremebox}

\begin{methodebox}[Variation de la constante]
Pour trouver $y_p$, on pose $y_p(x) = C(x)e^{-ax}$ avec $C(x)$ inconnue.
En substituant : $C'(x)e^{-ax} = b(x)$, donc $C'(x) = b(x)e^{ax}$, soit $C(x) = \int b(x)e^{ax}\,\mathrm{d}x$.
\end{methodebox}

\begin{definitionbox}[ED1 séparable]
Une équation $y' = f(x)g(y)$ est dite \emph{séparable} si le second membre est un produit d'une fonction de $x$ et d'une fonction de $y$.
\end{definitionbox}

\begin{methodebox}[Résolution d'une ED1 séparable]
Si $g(y) \neq 0$ :
\[
  \frac{y'}{g(y)} = f(x) \implies \int \frac{\mathrm{d}y}{g(y)} = \int f(x)\,\mathrm{d}x + C.
\]
On intègre les deux membres puis on exprime $y$ en fonction de $x$ (si possible).
\end{methodebox}

\begin{exemplebox}[ED1 séparable — refroidissement de Newton]
$\dfrac{\mathrm{d}T}{\mathrm{d}t} = -k(T - T_\infty)$, $k > 0$, $T_\infty$ température ambiante.
Séparation : $\dfrac{\mathrm{d}T}{T-T_\infty} = -k\,\mathrm{d}t$, d'où $\ln|T-T_\infty| = -kt + C$.
Solution : $T(t) = T_\infty + (T_0 - T_\infty)e^{-kt}$.
\end{exemplebox}

\section{Équations différentielles du second ordre à coefficients constants}

L'équation caractéristique $ar^2 + br + c = 0$ est la clé de résolution des ED2 à coefficients constants : l'ansatz $y = e^{rx}$ transforme l'équation différentielle en une équation algébrique, dont les racines déterminent la forme de la solution homogène. Le signe du discriminant $\Delta$ distingue trois comportements qualitativement différents — amortissement surcritique, critique ou oscillant — qui correspondent respectivement aux régimes surramorti, apériodique et oscillatoire en physique.

\begin{definitionbox}[ED2 linéaire à coefficients constants]
\[
  ay'' + by' + cy = f(x), \quad a,b,c \in \R,\; a \neq 0.
\]
L'\emph{équation caractéristique} associée est : $ar^2 + br + c = 0$, de discriminant $\Delta = b^2 - 4ac$.
\end{definitionbox}

\begin{theoremebox}[Solutions de l'homogène $ay''+by'+cy=0$]
\begin{itemize}
  \item $\Delta > 0$ : deux racines réelles $r_1, r_2$. $y_h = C_1 e^{r_1 x} + C_2 e^{r_2 x}$.
  \item $\Delta = 0$ : racine double $r_0 = -b/(2a)$. $y_h = (C_1 + C_2 x)e^{r_0 x}$.
  \item $\Delta < 0$ : racines complexes conjuguées $\alpha \pm i\beta$. $y_h = e^{\alpha x}(C_1\cos\beta x + C_2\sin\beta x)$.
\end{itemize}
\end{theoremebox}

\begin{methodebox}[Solution particulière par identification (second membre classique)]
\begin{center}
\begin{tabular}{ll}
\toprule
Second membre $f(x)$ & Forme de $y_p$ \\
\midrule
$P_n(x)$ (polynôme degré $n$), $0$ non racine & $Q_n(x)$ (polynôme degré $n$) \\
$P_n(x)$, $0$ racine de mult.\ $m$ & $x^m Q_n(x)$ \\
$e^{\alpha x}P_n(x)$, $\alpha$ non racine & $e^{\alpha x}Q_n(x)$ \\
$e^{\alpha x}P_n(x)$, $\alpha$ racine de mult.\ $m$ & $x^m e^{\alpha x}Q_n(x)$ \\
$e^{\alpha x}(A\cos\beta x + B\sin\beta x)$ & $e^{\alpha x}(P\cos\beta x + Q\sin\beta x)$ \\
\bottomrule
\end{tabular}
\end{center}
\end{methodebox}

\begin{center}
\begin{tikzpicture}
\begin{axis}[
  width=10cm, height=5cm,
  xlabel={$t$}, ylabel={$y(t)$},
  title={\small Solutions de $y''+y=0$ (oscillateur harmonique)},
  domain=0:4*pi, samples=150,
  xmin=0, xmax=4*pi,
  grid=major, grid style={dotted,gray!30},
  tick label style={font=\scriptsize},
  label style={font=\small},
  legend style={font=\scriptsize, at={(0.98,0.98)}, anchor=north east},
  xtick={0, 3.14159, 6.28318, 9.42478, 12.56637},
  xticklabels={$0$,$\pi$,$2\pi$,$3\pi$,$4\pi$},
]
\addplot[thick, caplpblue] {cos(deg(x))};
\addlegendentry{$\cos t$}
\addplot[thick, caplpred, dashed] {sin(deg(x))};
\addlegendentry{$\sin t$}
\addplot[thick, caplpgreen, dash dot] {cos(deg(x)) + sin(deg(x))};
\addlegendentry{$\cos t + \sin t$}
\end{axis}
\end{tikzpicture}
\captionof*{figure}{\small\itshape Solutions de l'oscillateur harmonique $y''+y=0$ : $\Delta < 0$, racines $\pm i$, solutions $\cos t$ et $\sin t$.}
\end{center}

\begin{exercicebox}[ED1 complète et ED2]
\begin{enumerate}
  \item Résoudre $y' - 2y = e^x$ avec $y(0) = 0$.
  \item Résoudre $y'' - 3y' + 2y = 4e^{3x}$.
\end{enumerate}
\end{exercicebox}

\begin{correctionbox}[Correction]
\textbf{1.} Homogène : $y_h = Ce^{2x}$. Variation de la constante : $C'(x)e^{2x} = e^x$, $C'(x)=e^{-x}$, $C(x)=-e^{-x}$.
$y_p = -e^{-x}\cdot e^{2x} = -e^x$. Solution générale : $y = Ce^{2x} - e^x$.
C.I. $y(0)=0$ : $C-1=0$, $C=1$. Solution : $y=e^{2x}-e^x$.

\textbf{2.} Éq. caractéristique : $r^2-3r+2=0$, racines $r_1=1$, $r_2=2$.
$y_h = C_1 e^x + C_2 e^{2x}$.
$\alpha=3$ n'est pas racine : $y_p = Ae^{3x}$. Substitution : $9A-9A+2A = 4$, $2A=4$, $A=2$.
$y_p = 2e^{3x}$. Solution générale : $y = C_1 e^x + C_2 e^{2x} + 2e^{3x}$.
\end{correctionbox}

\begin{retenirbox}
\textbf{Équations différentielles — À retenir}
\begin{itemize}
  \item ED1 : sol. générale $=$ homogène $+$ particulière.
  \item Homogène $y'+ay=0$ : $y=Ce^{-ax}$.
  \item Séparable $y'=f(x)g(y)$ : séparer, intégrer.
  \item ED2 : équation caractéristique $ar^2+br+c=0$, trois cas ($\Delta>0$, $=0$, $<0$).
  \item Sol. particulière par identification (tableau des formes).
\end{itemize}
\end{retenirbox}

%% === CHAPITRE 7 ===
\chapter{Suites numériques réelles}

\section*{Objectifs}
\begin{itemize}[leftmargin=1.5em]
  \item Définir et étudier les suites explicites et récurrentes.
  \item Maîtriser les théorèmes de convergence (monotonie bornée, gendarmes, suites adjacentes).
  \item Analyser graphiquement une suite récurrente (toile d'araignée).
  \item Connaître les suites arithmétiques, géométriques et leurs propriétés.
  \item Extraire des sous-suites et utiliser les suites extraites.
\end{itemize}

\section{Convergence d'une suite}

La convergence d'une suite formalise rigoureusement l'idée qu'une suite « tend vers » une valeur limite : la définition $\varepsilon$-$N$ de Cauchy-Weierstrass exige qu'à partir d'un certain rang $N$, tous les termes restent arbitrairement proches de la limite $\ell$. Ce formalisme permet de démontrer des théorèmes opératoires et des critères de convergence (suites monotones bornées, suites adjacentes) qui sont les outils de travail de l'analyse séquentielle.

\begin{definitionbox}[Limite d'une suite]
La suite $(u_n)$ \emph{converge vers $\ell \in \R$} si :
\[
  \forall \varepsilon > 0,\;\exists N \in \N : n \geq N \implies |u_n - \ell| < \varepsilon.
\]
On note $\lim_{n\to+\infty} u_n = \ell$ ou $u_n \to \ell$.
Si $(u_n)$ ne converge pas, elle \emph{diverge}.
\end{definitionbox}

\begin{theoremebox}[Unicité de la limite]
Si $(u_n)$ converge, sa limite est unique.
\end{theoremebox}

\begin{theoremebox}[Théorèmes opératoires]
Si $u_n \to \ell$ et $v_n \to m$ :
\begin{itemize}
  \item $u_n + v_n \to \ell + m$ ; $\lambda u_n \to \lambda \ell$ ; $u_n v_n \to \ell m$ ; $u_n/v_n \to \ell/m$ si $m\neq 0$.
  \item \textbf{Théorème des gendarmes.} Si $u_n \leq w_n \leq v_n$ et $u_n, v_n \to \ell$, alors $w_n \to \ell$.
  \item \textbf{Comparaison.} Si $u_n \leq v_n$ pour $n$ assez grand et $u_n,v_n$ convergent, alors $\ell \leq m$.
\end{itemize}
\end{theoremebox}

\begin{theoremebox}[Suites monotones et bornées]
Toute suite \emph{croissante et majorée} converge.
Toute suite \emph{décroissante et minorée} converge.
\end{theoremebox}

\begin{preuvebox}
Soit $(u_n)$ croissante et majorée par $M$. L'ensemble $\{u_n \mid n\in\N\}$ est non vide et majoré, il admet donc un supremum $\ell = \sup\{u_n\}$. Pour tout $\varepsilon>0$, il existe $N$ tel que $u_N > \ell - \varepsilon$. Comme $(u_n)$ est croissante, pour $n\geq N$ : $\ell - \varepsilon < u_N \leq u_n \leq \ell$, d'où $|u_n - \ell| < \varepsilon$. \qed
\end{preuvebox}

\begin{theoremebox}[Suites adjacentes]
Deux suites $(u_n)$ et $(v_n)$ sont \emph{adjacentes} si :
\begin{enumerate}[label=\roman*)]
  \item $(u_n)$ est croissante et $(v_n)$ est décroissante,
  \item $v_n - u_n \to 0$.
\end{enumerate}
Dans ce cas, $(u_n)$ et $(v_n)$ convergent vers la même limite $\ell$, et $u_n \leq \ell \leq v_n$ pour tout $n$.
\end{theoremebox}

\section{Suites récurrentes $u_{n+1} = f(u_n)$}

L'étude d'une suite récurrente $u_{n+1} = f(u_n)$ repose sur la recherche des points fixes de $f$, qui sont les seules limites candidates : si $(u_n)$ converge vers $\ell$, alors en passant à la limite dans la relation de récurrence on obtient $\ell = f(\ell)$. La méthode graphique de la toile d'araignée — construire la suite en alternant lectures sur la courbe de $f$ et sur la droite $y = x$ — permet de visualiser la dynamique et de conjecturer la convergence avant tout calcul.

\begin{methodebox}[Étude d'une suite récurrente]
\begin{enumerate}
  \item \textbf{Points fixes.} Résoudre $f(x) = x$ : ce sont les limites candidates.
  \item \textbf{Invariance d'un intervalle.} Chercher $I$ tel que $f(I) \subseteq I$ et $u_0 \in I$.
  \item \textbf{Monotonie.} Étudier le signe de $u_{n+1}-u_n = f(u_n)-u_n$.
  \item \textbf{Convergence.} Si $(u_n)$ est monotone et bornée dans $I$, elle converge.
  \item \textbf{Identification de la limite.} Par passage à la limite dans $u_{n+1}=f(u_n)$, la limite $\ell$ vérifie $\ell = f(\ell)$.
\end{enumerate}
\end{methodebox}

\begin{center}
\begin{tikzpicture}
\begin{axis}[
  width=8cm, height=7cm,
  xlabel={$x$}, ylabel={$y$},
  title={\small Toile d'araignée : $u_{n+1} = \sqrt{u_n+2}$, $u_0=0$},
  domain=0:4, samples=100,
  xmin=0, xmax=4, ymin=0, ymax=4,
  axis equal,
  grid=major, grid style={dotted,gray!30},
  tick label style={font=\scriptsize},
  label style={font=\small},
]
% Courbe f(x) = sqrt(x+2)
\addplot[thick, caplpblue] {sqrt(x+2)};
\addlegendentry{$y=\sqrt{x+2}$}
% Diagonale y=x
\addplot[thick, dashed, gray] {x};
\addlegendentry{$y=x$}
% Toile d'araignée (itérations)
\addplot[thick, caplpred, -] coordinates {
  (0,0) (0,1.4142) (1.4142,1.4142) (1.4142,1.8478)
  (1.8478,1.8478) (1.8478,1.9494) (1.9494,1.9494)
  (1.9494,1.9747) (1.9747,1.9747) (1.9747,1.9874)
  (1.9874,1.9874) (1.9874,2.0)
};
% Point fixe
\addplot[mark=*, mark size=3pt, caplpgreen] coordinates {(2,2)};
\node[font=\scriptsize, caplpgreen, right] at (axis cs:2.05,2.1) {$\ell=2$};
\end{axis}
\end{tikzpicture}
\captionof*{figure}{\small\itshape Toile d'araignée : la suite $u_{n+1}=\sqrt{u_n+2}$ converge vers le point fixe $\ell=2$.}
\end{center}

\section{Suites classiques}

Les suites arithmétiques et géométriques sont les deux familles fondamentales, dont toutes les propriétés se calculent explicitement. Elles interviennent comme cas particuliers de suites récurrentes affines, et leurs formules de somme sont des outils de calcul incontournables en probabilités (loi géométrique) comme en finance (capitalisation). La maîtrise parfaite de ces deux familles conditionne la résolution rapide d'une grande variété d'exercices.

\begin{definitionbox}[Suite arithmétique]
$(u_n)$ est \emph{arithmétique} de raison $r$ si $u_{n+1} = u_n + r$ pour tout $n$.
Terme général : $u_n = u_0 + nr$.
Somme : $\displaystyle\sum_{k=0}^{n} u_k = (n+1)\frac{u_0+u_n}{2}$.
\end{definitionbox}

\begin{definitionbox}[Suite géométrique]
$(u_n)$ est \emph{géométrique} de raison $q \neq 0$ si $u_{n+1} = qu_n$ pour tout $n$.
Terme général : $u_n = u_0 q^n$.
Somme : $\displaystyle\sum_{k=0}^{n} u_k = u_0\,\frac{1-q^{n+1}}{1-q}$ si $q \neq 1$.
\end{definitionbox}

\begin{theoremebox}[Limites usuelles de suites]
Pour $a > 1$, $\alpha > 0$, $q \in \R$ :
\[
  \lim_{n\to+\infty} \frac{n^\alpha}{a^n} = 0, \quad
  \lim_{n\to+\infty} \frac{\ln n}{n^\alpha} = 0, \quad
  \lim_{n\to+\infty} q^n = \begin{cases} 0 & \text{si } |q|<1 \\ 1 & \text{si } q=1 \\ +\infty & \text{si } q>1 \end{cases}.
\]
\end{theoremebox}

\section{Suites extraites}

Les sous-suites (ou suites extraites) sont un outil puissant pour caractériser la convergence ou la divergence d'une suite. Le théorème de Bolzano-Weierstrass affirme que toute suite bornée admet une sous-suite convergente, résultat profond qui repose sur la complétude de $\mathbb{R}$. En pratique, montrer que deux sous-suites extraites ont des limites différentes est la méthode standard pour prouver qu'une suite diverge.

\begin{definitionbox}[Suite extraite]
Une \emph{suite extraite} (ou sous-suite) de $(u_n)$ est une suite $(u_{\varphi(n)})$ où
$\varphi : \N \to \N$ est strictement croissante.
\end{definitionbox}

\begin{theoremebox}
Si $(u_n) \to \ell$, alors toute suite extraite de $(u_n)$ tend vers $\ell$.
Réciproque utile : si $(u_{2n})$ et $(u_{2n+1})$ tendent toutes deux vers $\ell$, alors $(u_n) \to \ell$.
\end{theoremebox}

\begin{exercicebox}[Suite récurrente — type concours]
Soit $u_0 = 0$ et $u_{n+1} = \dfrac{u_n + 3}{u_n + 1}$ pour $n \geq 0$.
\begin{enumerate}
  \item Montrer que $0 \leq u_n \leq \sqrt{3}$ pour tout $n$ (raisonnement par récurrence).
  \item Montrer que $(u_n)$ est croissante.
  \item En déduire que $(u_n)$ converge et calculer sa limite.
\end{enumerate}
\end{exercicebox}

\begin{correctionbox}[Correction]
Posons $f(x) = \dfrac{x+3}{x+1}$.

\textbf{1.} Récurrence. $u_0=0\in[0,\sqrt{3}]$. Si $u_n\in[0,\sqrt{3}]$ : $f(x)\geq 0$ et $f(x)^2-3 = \dfrac{(x+3)^2-3(x+1)^2}{(x+1)^2} = \dfrac{-2x^2+6x+6\cdot x\cdots}{(x+1)^2}$... En simplifiant : $f(x)^2-3 = \dfrac{x^2-2x+6-3x^2-6x-3}{(x+1)^2}$... Calculons : $f(x)^2 = \dfrac{(x+3)^2}{(x+1)^2}$, $3 = \dfrac{3(x+1)^2}{(x+1)^2}$. $(x+3)^2-3(x+1)^2 = x^2+6x+9-3x^2-6x-3 = -2x^2+6 = 2(3-x^2) \geq 0$ pour $x\in[0,\sqrt{3}]$. Donc $f(x)\leq\sqrt{3}$. Par récurrence, $u_n\in[0,\sqrt{3}]$.

\textbf{2.} $u_{n+1}-u_n = \dfrac{u_n+3}{u_n+1}-u_n = \dfrac{3-u_n^2}{u_n+1} \geq 0$ car $u_n\leq\sqrt{3}$. Suite croissante.

\textbf{3.} Croissante et majorée par $\sqrt{3}$ : converge vers $\ell$. Passage à la limite : $\ell = \dfrac{\ell+3}{\ell+1}$, soit $\ell^2+\ell=\ell+3$, $\ell^2=3$, $\ell=\sqrt{3}$ (positif).
\end{correctionbox}

\begin{retenirbox}
\textbf{Suites — À retenir}
\begin{itemize}
  \item Monotone bornée $\Rightarrow$ convergente.
  \item Gendarmes : $u_n\leq w_n\leq v_n$, $u_n,v_n\to\ell$ $\Rightarrow$ $w_n\to\ell$.
  \item $u_{n+1}=f(u_n)$ : point fixe = limite candidate ; étudier monotonie + bornes.
  \item Arithmétique : $u_n=u_0+nr$ ; géométrique : $u_n=u_0 q^n$.
  \item Suite extraite d'une suite convergente : même limite.
\end{itemize}
\end{retenirbox}

%% === CHAPITRE 8 ===
\chapter{Algèbre linéaire — Systèmes et matrices (A.1)}

\section*{Objectifs}
\begin{itemize}[leftmargin=1.5em]
  \item Résoudre les systèmes linéaires $n \leq 3$ par la méthode du pivot de Gauss.
  \item Effectuer les opérations matricielles : addition, produit, transposition.
  \item Calculer l'inverse d'une matrice $2\times2$ et $3\times3$.
  \item Discuter un système selon un paramètre.
\end{itemize}

\section{Systèmes linéaires}

La méthode du pivot de Gauss est l'algorithme universel de résolution des systèmes linéaires : elle transforme le système initial en un système triangulaire équivalent par des opérations élémentaires sur les lignes, puis résout par substitution ascendante. Cet algorithme est à la base de tous les solveurs numériques modernes et son analyse en termes de rang de la matrice permet de discuter complètement le nombre de solutions d'un système selon ses paramètres.

\begin{definitionbox}[Système linéaire]
Un \emph{système linéaire} de $n$ équations à $p$ inconnues $(x_1,\ldots,x_p)$ est de la forme :
\[
  \begin{cases}
    a_{11}x_1 + \cdots + a_{1p}x_p = b_1 \\
    \hspace{2cm}\vdots \\
    a_{n1}x_1 + \cdots + a_{np}x_p = b_n
  \end{cases}
\]
Il est \emph{homogène} si $b_1=\cdots=b_n=0$. Sous forme matricielle : $AX = B$.
\end{definitionbox}

\begin{methodebox}[Pivot de Gauss]
Opérations élémentaires sur les lignes (ne changent pas les solutions) :
\begin{enumerate}
  \item Échanger deux lignes : $L_i \leftrightarrow L_j$.
  \item Multiplier une ligne par un scalaire non nul : $L_i \leftarrow \lambda L_i$.
  \item Ajouter un multiple d'une ligne à une autre : $L_i \leftarrow L_i + \lambda L_j$.
\end{enumerate}
On transforme la matrice augmentée $[A|B]$ en forme échelonnée réduite (triangulaire supérieure), puis on remonte par substitution.
\end{methodebox}

\begin{exemplebox}[Résolution par pivot]
$\begin{cases} x+2y+z=4 \\ 2x+y-z=1 \\ x-y+2z=3 \end{cases}$

Matrice augmentée $\to$ pivot :
$L_2 \leftarrow L_2-2L_1$, $L_3 \leftarrow L_3-L_1$ :
$\begin{pmatrix}1&2&1&|&4\\0&-3&-3&|&-7\\0&-3&1&|&-1\end{pmatrix}$

$L_3 \leftarrow L_3-L_2$ :
$\begin{pmatrix}1&2&1&|&4\\0&-3&-3&|&-7\\0&0&4&|&6\end{pmatrix}$

Remontée : $z=3/2$, $y=(-7+9/2)/(-3)=5/6$... (calcul complet donne $z=3/2$, $y=1/6$, $x=11/6$).
\end{exemplebox}

\begin{theoremebox}[Discussion d'un système]
Un système $AX=B$ est :
\begin{itemize}
  \item \emph{Compatible} (au moins une solution) si le rang de $A$ égale le rang de $[A|B]$.
  \item \emph{À solution unique} si $\mathrm{rang}(A)=\mathrm{rang}([A|B])=p$ (nombre d'inconnues).
  \item \emph{À infinité de solutions} si $\mathrm{rang}(A)=\mathrm{rang}([A|B])<p$.
  \item \emph{Incompatible} (sans solution) si $\mathrm{rang}(A)<\mathrm{rang}([A|B])$.
\end{itemize}
\end{theoremebox}

\section{Calcul matriciel}

Les matrices ne sont pas seulement des tableaux de nombres : elles représentent des applications linéaires entre espaces vectoriels, et les opérations matricielles reflètent la composition de ces applications. En particulier, le produit matriciel $AB$ correspond à la composition de l'application représentée par $B$ suivie de celle représentée par $A$, ce qui explique sa non-commutativité générale. Cette interprétation géométrique est indispensable pour comprendre la signification des valeurs propres et de la diagonalisation.

\begin{definitionbox}[Opérations sur les matrices]
Soit $A=(a_{ij}) \in \mathcal{M}_{n,p}(\R)$, $B=(b_{ij}) \in \mathcal{M}_{n,p}(\R)$, $C=(c_{jk})\in\mathcal{M}_{p,q}(\R)$.
\begin{itemize}
  \item \textbf{Somme :} $(A+B)_{ij} = a_{ij}+b_{ij}$.
  \item \textbf{Produit scalaire :} $(\lambda A)_{ij} = \lambda a_{ij}$.
  \item \textbf{Produit :} $(AC)_{ik} = \sum_{j=1}^p a_{ij}c_{jk}$ (matrice $n\times q$).
  \item \textbf{Transposée :} $A^\top \in \mathcal{M}_{p,n}(\R)$, $(A^\top)_{ji}=a_{ij}$.
\end{itemize}
Attention : $AB \neq BA$ en général.
\end{definitionbox}

\begin{theoremebox}[Inverse d'une matrice $2\times 2$]
Pour $A = \begin{pmatrix}a&b\\c&d\end{pmatrix}$ avec $\det(A)=ad-bc\neq 0$ :
\[
  A^{-1} = \frac{1}{ad-bc}\begin{pmatrix}d&-b\\-c&a\end{pmatrix}.
\]
\end{theoremebox}

\begin{methodebox}[Inverse d'une matrice $3\times 3$ (Gauss-Jordan)]
Former $[A|I_3]$ et appliquer le pivot de Gauss jusqu'à obtenir $[I_3|A^{-1}]$.
Si $A$ ne se réduit pas en $I_3$ : $A$ n'est pas inversible.
\end{methodebox}

\begin{erreurbox}
\begin{itemize}
  \item \textbf{Erreur 1.} $(AB)^{-1} = B^{-1}A^{-1}$ (et non $A^{-1}B^{-1}$).
  \item \textbf{Erreur 2.} $(A+B)^{-1} \neq A^{-1}+B^{-1}$.
  \item \textbf{Erreur 3.} $AB=0 \not\Rightarrow A=0$ ou $B=0$ (diviseurs de zéro matriciels).
\end{itemize}
\end{erreurbox}

\begin{exercicebox}[Système avec paramètre]
Discuter selon $m\in\R$ le système :
$\begin{cases}mx + y = 1\\x + my = 1\end{cases}$
\end{exercicebox}

\begin{correctionbox}[Correction]
Matrice $A=\begin{pmatrix}m&1\\1&m\end{pmatrix}$, $\det(A)=m^2-1=(m-1)(m+1)$.

\textbf{Si $m\neq\pm1$} : $A$ inversible, unique solution $x=y=\dfrac{1}{m+1}$.

\textbf{Si $m=1$} : système $x+y=1$ (deux fois) $\to$ infinité de solutions : $\{(t,1-t),\;t\in\R\}$.

\textbf{Si $m=-1$} : $-x+y=1$ et $x-y=1$ (contradictoires) $\to$ incompatible.
\end{correctionbox}

\begin{retenirbox}
\textbf{Algèbre linéaire A.1 — À retenir}
\begin{itemize}
  \item Pivot de Gauss : échelonner $[A|B]$ puis remonter.
  \item Rang et discussion : unique $\iff$ rang$=n=p$ ; $\infty$ solutions $\iff$ rang$<p$ ; incompatible $\iff$ rangs discordants.
  \item $(AB)^{-1}=B^{-1}A^{-1}$ ; $(AB)^\top=B^\top A^\top$.
  \item Inverse $2\times2$ : swap diagonale, changer signes, diviser par $\det$.
\end{itemize}
\end{retenirbox}

%% === CHAPITRE 9 ===
\chapter{Arithmétique des entiers}

\section*{Objectifs}
\begin{itemize}[leftmargin=1.5em]
  \item Maîtriser la représentation en base 10 et base 2 ; effectuer des conversions.
  \item Appliquer la division euclidienne ; définir divisibilité et multiples.
  \item Connaître les propriétés des nombres premiers.
  \item Calculer le PGCD et le PPCM par l'algorithme d'Euclide.
  \item Utiliser l'identité de Bézout et les congruences.
\end{itemize}

\section{Représentation des entiers}

La notation positionnelle est le fondement de notre système numérique : chaque chiffre d'une écriture en base $b$ représente une puissance de $b$, et l'entier est la somme pondérée de ces puissances. La base 2 (binaire), utilisée par tous les processeurs, et la base 16 (hexadécimale) sont les systèmes de référence en informatique. La maîtrise des conversions est donc à la fois un exercice mathématique fondamental et une compétence directement utile en algorithmique.

\begin{definitionbox}[Écriture en base $b$]
Tout entier $n \geq 0$ s'écrit de façon unique en base $b \geq 2$ :
\[
  n = a_k b^k + a_{k-1}b^{k-1} + \cdots + a_1 b + a_0, \quad 0 \leq a_i < b.
\]
On note $n = (a_k a_{k-1}\ldots a_0)_b$.
En base 2, les chiffres sont des \emph{bits} (0 ou 1) ; en base 16, on ajoute les symboles A–F.
\end{definitionbox}

\begin{methodebox}[Conversion décimal $\to$ base $b$ (divisions successives)]
Diviser $n$ par $b$, noter le reste, recommencer avec le quotient jusqu'à 0.
Lire les restes de bas en haut.

\textbf{Exemple.} $26$ en base 2 :
$26 = 13\times2+0$, $13=6\times2+1$, $6=3\times2+0$, $3=1\times2+1$, $1=0\times2+1$.
Résultat : $(11010)_2$.
\end{methodebox}

\section{Division euclidienne et divisibilité}

La division euclidienne est le théorème fondateur de l'arithmétique des entiers : elle garantit l'existence et l'unicité du quotient et du reste, deux entiers qui encodent toute la relation entre dividende et diviseur. La divisibilité, définie par l'annulation du reste, structure l'ensemble $\mathbb{Z}$ et engendre toute la théorie des nombres premiers, dont le théorème fondamental de l'arithmétique (décomposition unique en facteurs premiers) est le résultat central.

\begin{theoremebox}[Division euclidienne]
Pour tout $a \in \Z$ et $b \in \Z^*$, il existe un unique couple $(q,r) \in \Z\times\N$ tel que :
\[
  a = bq + r, \quad 0 \leq r < |b|.
\]
$q$ est le \emph{quotient} et $r$ le \emph{reste}.
\end{theoremebox}

\begin{definitionbox}[Divisibilité]
$b$ \emph{divise} $a$ (noté $b \mid a$) si le reste de la division de $a$ par $b$ est nul, i.e.\ $\exists k\in\Z : a=kb$.
\end{definitionbox}

\begin{definitionbox}[Nombres premiers]
Un entier $p \geq 2$ est \emph{premier} s'il n'admet pas d'autre diviseur positif que $1$ et lui-même.
Tout entier $n \geq 2$ se décompose de façon unique (à l'ordre des facteurs près) en produit de facteurs premiers (\emph{théorème fondamental de l'arithmétique}).
\end{definitionbox}

\section{PGCD, PPCM et algorithme d'Euclide}

L'algorithme d'Euclide est l'un des plus anciens algorithmes connus, et l'un des plus efficaces : il calcule le PGCD de deux entiers en $O(\log(\min(a,b)))$ divisions euclidiennes. Son prolongement — l'algorithme d'Euclide étendu — permet de construire explicitement les coefficients de Bézout, c'est-à-dire les entiers $u, v$ tels que $au + bv = \pgcd(a,b)$. Ces coefficients sont indispensables pour inverser un entier modulo $n$ et résoudre les équations de congruence.

\begin{definitionbox}[PGCD et PPCM]
\begin{itemize}
  \item Le \emph{PGCD} de $a$ et $b$ (notés $a\wedge b$ ou $\pgcd(a,b)$) est le plus grand diviseur commun.
  \item Le \emph{PPCM} de $a$ et $b$ (noté $a\vee b$ ou $\ppcm(a,b)$) est le plus petit multiple commun positif.
  \item $\pgcd(a,b)\times\ppcm(a,b) = |ab|$.
\end{itemize}
\end{definitionbox}

\begin{theoremebox}[Algorithme d'Euclide]
$\pgcd(a,b) = \pgcd(b, a \bmod b)$. On itère jusqu'au reste nul : le dernier reste non nul est le PGCD.
\end{theoremebox}

\begin{preuvebox}
Si $a = bq+r$, alors $d \mid a$ et $d \mid b \iff d \mid b$ et $d \mid r$. Donc les diviseurs communs de $a$ et $b$ sont exactement les diviseurs communs de $b$ et $r$, d'où $\pgcd(a,b)=\pgcd(b,r)$. \qed
\end{preuvebox}

\begin{theoremebox}[Identité de Bézout]
Pour $a, b \in \Z$ non tous deux nuls : $\exists (u,v)\in\Z^2$ tels que $au + bv = \pgcd(a,b)$.

\emph{Corollaire.} $a$ et $b$ sont \emph{premiers entre eux} ($\pgcd(a,b)=1$) $\iff \exists u,v : au+bv=1$.
\end{theoremebox}

\begin{definitionbox}[Congruences]
On dit que $a \equiv b \pmod{n}$ (lire « $a$ est congru à $b$ modulo $n$ ») si $n \mid (a-b)$.
Les congruences sont compatibles avec l'addition et la multiplication :
si $a\equiv a'$ et $b\equiv b' \pmod n$, alors $a+b\equiv a'+b'$ et $ab\equiv a'b' \pmod n$.
\end{definitionbox}

\begin{center}
\begin{tikzpicture}[>=Stealth, node distance=1.2cm, font=\small]
  \node[draw, rounded corners, fill=caplpbluebg] (a) {$a=252$, $b=198$};
  \node[draw, rounded corners, fill=caplpgreenbg, below of=a] (s1) {$252 = 1\times198+54$};
  \node[draw, rounded corners, fill=caplpgreenbg, below of=s1] (s2) {$198 = 3\times54+36$};
  \node[draw, rounded corners, fill=caplpgreenbg, below of=s2] (s3) {$54 = 1\times36+18$};
  \node[draw, rounded corners, fill=caplpgreenbg, below of=s3] (s4) {$36 = 2\times18+0$};
  \node[draw, rounded corners, fill=caplpredbg, below of=s4] (res) {$\pgcd(252,198)=18$};
  \draw[->] (a)--(s1); \draw[->](s1)--(s2); \draw[->](s2)--(s3); \draw[->](s3)--(s4); \draw[->](s4)--(res);
\end{tikzpicture}
\captionof*{figure}{\small\itshape Algorithme d'Euclide : $\pgcd(252,198)=18$.}
\end{center}

\begin{exercicebox}[Bézout et congruences]
\begin{enumerate}
  \item Calculer $\pgcd(105,70)$ et trouver $(u,v)\in\Z^2$ tels que $105u+70v=\pgcd(105,70)$.
  \item Résoudre la congruence $3x \equiv 1 \pmod 7$.
\end{enumerate}
\end{exercicebox}

\begin{correctionbox}[Correction]
\textbf{1.} $105=1\times70+35$, $70=2\times35+0$. $\pgcd=35$.
Remontée : $35=105-1\times70$, donc $u=1$, $v=-1$.

\textbf{2.} On cherche l'inverse de $3$ mod $7$. $\pgcd(3,7)=1$ (Bézout) : $3\times5=15=2\times7+1$, soit $3\times5\equiv1\pmod7$. Donc $x\equiv5\pmod7$.
\end{correctionbox}

\begin{retenirbox}
\textbf{Arithmétique — À retenir}
\begin{itemize}
  \item Division euclidienne : $a=bq+r$, $0\leq r<|b|$.
  \item Euclide : $\pgcd(a,b)=\pgcd(b,a\bmod b)$ jusqu'au reste nul.
  \item Bézout : $\exists u,v : au+bv=\pgcd(a,b)$.
  \item Congruences : compatibles avec $+$ et $\times$.
  \item Base 2 : divisions successives par 2, lire les restes.
\end{itemize}
\end{retenirbox}

%% === CHAPITRE 10 ===
\chapter{Calcul vectoriel et espaces euclidiens}

\section*{Objectifs}
\begin{itemize}[leftmargin=1.5em]
  \item Manipuler les vecteurs dans $\R^2$ et $\R^3$ (opérations, norme, angle).
  \item Définir et utiliser le produit scalaire.
  \item Définir et calculer le produit vectoriel dans $\R^3$.
  \item Maîtriser les transformations usuelles : translation, rotation, symétrie.
  \item Travailler avec le barycentre.
\end{itemize}

\section{Vecteurs dans $\R^2$ et $\R^3$}

Un vecteur est une translation libre : il est défini par une direction, un sens et une norme, indépendamment du point d'application. Les coordonnées dans la base canonique de $\mathbb{R}^2$ ou $\mathbb{R}^3$ permettent de ramener toutes les opérations géométriques (somme, norme, colinéarité) à des calculs algébriques simples, ce qui est le grand avantage de la géométrie vectorielle sur la géométrie synthétique classique.

\begin{definitionbox}[Vecteur, norme, angle]
Un \emph{vecteur} $\vec{u} = (x,y)$ dans $\R^2$ (ou $(x,y,z)$ dans $\R^3$) est caractérisé par ses coordonnées dans la base canonique.
\begin{itemize}
  \item \textbf{Norme} : $\|\vec{u}\| = \sqrt{x^2+y^2}$ dans $\R^2$, $\sqrt{x^2+y^2+z^2}$ dans $\R^3$.
  \item \textbf{Vecteur unitaire} : $\|\vec{u}\|=1$.
  \item \textbf{Collinéarité} (dans $\R^2$) : $\vec{u}$ et $\vec{v}$ colinéaires $\iff xv_y - yv_x = 0$.
\end{itemize}
\end{definitionbox}

\begin{definitionbox}[Barycentre]
Le \emph{barycentre} de points pondérés $(A_i, \lambda_i)$ avec $\sum\lambda_i \neq 0$ est le point $G$ tel que :
\[
  \sum_{i} \lambda_i \overrightarrow{GA_i} = \vec{0}.
\]
En coordonnées : $G = \dfrac{\sum \lambda_i A_i}{\sum \lambda_i}$.
\end{definitionbox}

\section{Produit scalaire}

Le produit scalaire est l'outil qui encode simultanément l'angle entre deux vecteurs et leur norme : la formule $\vec{u} \cdot \vec{v} = \|\vec{u}\|\,\|\vec{v}\|\cos\theta$ relie géométrie et algèbre de façon directe. En particulier, l'orthogonalité — condition fondamentale en géométrie, en physique et en traitement du signal — se traduit par l'annulation du produit scalaire. L'inégalité de Cauchy-Schwarz, qui en découle, est l'une des inégalités les plus utilisées en mathématiques.

\begin{definitionbox}[Produit scalaire dans $\R^n$]
Le \emph{produit scalaire} de $\vec{u}=(x_1,y_1)$ et $\vec{v}=(x_2,y_2)$ dans $\R^2$ est :
\[
  \vec{u}\cdot\vec{v} = x_1x_2 + y_1y_2.
\]
Dans $\R^3$ : $\vec{u}\cdot\vec{v} = x_1x_2+y_1y_2+z_1z_2$.

Formule géométrique : $\vec{u}\cdot\vec{v} = \|\vec{u}\|\,\|\vec{v}\|\cos\theta$ où $\theta$ est l'angle entre $\vec{u}$ et $\vec{v}$.
\end{definitionbox}

\begin{theoremebox}[Propriétés du produit scalaire]
\begin{itemize}
  \item Symétrie : $\vec{u}\cdot\vec{v}=\vec{v}\cdot\vec{u}$.
  \item Bilinéarité : $(\lambda\vec{u}+\vec{v})\cdot\vec{w}=\lambda(\vec{u}\cdot\vec{w})+\vec{v}\cdot\vec{w}$.
  \item Positivité : $\vec{u}\cdot\vec{u}=\|\vec{u}\|^2\geq0$, et $=0\iff\vec{u}=\vec{0}$.
  \item \textbf{Orthogonalité} : $\vec{u}\perp\vec{v}\iff\vec{u}\cdot\vec{v}=0$.
  \item \textbf{Inégalité de Cauchy-Schwarz} : $|\vec{u}\cdot\vec{v}|\leq\|\vec{u}\|\,\|\vec{v}\|$.
\end{itemize}
\end{theoremebox}

\section{Produit vectoriel dans $\R^3$}

Le produit vectoriel $\vec{u} \wedge \vec{v}$ est le vecteur de $\mathbb{R}^3$ orthogonal au plan engendré par $\vec{u}$ et $\vec{v}$, dont la norme est l'aire du parallélogramme qu'ils délimitent et dont le sens est donné par la règle de la main droite. Ce produit, propre à $\mathbb{R}^3$, est l'outil naturel pour calculer les équations de plans, les moments en mécanique et les normales à une surface.

\begin{definitionbox}[Produit vectoriel]
Pour $\vec{u}=(u_1,u_2,u_3)$ et $\vec{v}=(v_1,v_2,v_3)$ dans $\R^3$ :
\[
  \vec{u}\wedge\vec{v} = \begin{vmatrix}\vec{e}_1&\vec{e}_2&\vec{e}_3\\u_1&u_2&u_3\\v_1&v_2&v_3\end{vmatrix}
  = \begin{pmatrix}u_2v_3-u_3v_2\\u_3v_1-u_1v_3\\u_1v_2-u_2v_1\end{pmatrix}.
\]
\end{definitionbox}

\begin{theoremebox}[Propriétés du produit vectoriel]
\begin{itemize}
  \item \textbf{Antisymétrie} : $\vec{u}\wedge\vec{v} = -\vec{v}\wedge\vec{u}$.
  \item \textbf{Orthogonalité} : $(\vec{u}\wedge\vec{v})\perp\vec{u}$ et $(\vec{u}\wedge\vec{v})\perp\vec{v}$.
  \item \textbf{Norme} : $\|\vec{u}\wedge\vec{v}\| = \|\vec{u}\|\,\|\vec{v}\|\,|\sin\theta|$ = aire du parallélogramme.
  \item \textbf{Colinéarité} : $\vec{u}\wedge\vec{v}=\vec{0}\iff\vec{u}$ et $\vec{v}$ colinéaires.
\end{itemize}
\end{theoremebox}

\section{Transformations usuelles}

Les transformations affines — rotations, réflexions, translations, homothéties — sont les isométries et similitudes du plan et de l'espace. Leur représentation matricielle (pour les parties linéaires) permet de les composer par produit matriciel et de les appliquer efficacement à des ensembles de points. Ces transformations sont omniprésentes en géométrie différentielle, en infographie et en mécanique des solides.

\begin{definitionbox}[Transformations planes]
Dans un repère orthonormé de $\R^2$ :
\begin{itemize}
  \item \textbf{Translation} de vecteur $\vec{t}=(a,b)$ : $T(x,y)=(x+a,y+b)$.
  \item \textbf{Rotation} d'angle $\theta$ autour de l'origine : $R_\theta(x,y)=(x\cos\theta-y\sin\theta,\;x\sin\theta+y\cos\theta)$.
    Matrice : $\begin{pmatrix}\cos\theta&-\sin\theta\\\sin\theta&\cos\theta\end{pmatrix}$.
  \item \textbf{Symétrie axiale} par rapport à $Ox$ : $(x,y)\mapsto(x,-y)$.
  \item \textbf{Symétrie centrale} d'centre $O$ : $(x,y)\mapsto(-x,-y)$.
  \item \textbf{Homothétie} de rapport $k$ et centre $O$ : $(x,y)\mapsto(kx,ky)$.
\end{itemize}
\end{definitionbox}

\begin{center}
\begin{tikzpicture}[scale=1.2, >=Stealth]
  \draw[->] (-0.3,0)--(3.2,0) node[right,font=\scriptsize]{$x$};
  \draw[->] (0,-0.3)--(0,2.8) node[above,font=\scriptsize]{$y$};
  % Vecteur u
  \draw[->, very thick, caplpblue] (0,0)--(2,1) node[above right,font=\small]{$\vec{u}$};
  % Vecteur v
  \draw[->, very thick, caplpred] (0,0)--(0.5,2) node[above right,font=\small]{$\vec{v}$};
  % Parallélogramme
  \draw[dashed, caplpgreen] (2,1)--(2.5,3)--(0.5,2);
  % Aire
  \node[font=\scriptsize, caplpgreen] at (1.2,1.8) {Aire$=\|\vec{u}\wedge\vec{v}\|$};
  % Angle
  \draw[caplporange, ->] (0.4,0) arc[start angle=0, end angle=76, radius=0.4];
  \node[font=\scriptsize, caplporange] at (0.55,0.3) {$\theta$};
\end{tikzpicture}
\captionof*{figure}{\small\itshape La norme du produit vectoriel $\|\vec{u}\wedge\vec{v}\|$ est l'aire du parallélogramme engendré par $\vec{u}$ et $\vec{v}$.}
\end{center}

\begin{exercicebox}[Produit scalaire et vectoriel]
Soient $\vec{u}=(1,2,-1)$ et $\vec{v}=(3,0,2)$ dans $\R^3$.
\begin{enumerate}
  \item Calculer $\vec{u}\cdot\vec{v}$, $\|\vec{u}\|$, $\|\vec{v}\|$ et l'angle $\theta$ entre $\vec{u}$ et $\vec{v}$.
  \item Calculer $\vec{u}\wedge\vec{v}$ et vérifier qu'il est orthogonal à $\vec{u}$ et $\vec{v}$.
  \item Calculer l'aire du parallélogramme engendré par $\vec{u}$ et $\vec{v}$.
\end{enumerate}
\end{exercicebox}

\begin{correctionbox}[Correction]
\textbf{1.} $\vec{u}\cdot\vec{v}=3+0-2=1$. $\|\vec{u}\|=\sqrt{6}$, $\|\vec{v}\|=\sqrt{13}$.
$\cos\theta=\dfrac{1}{\sqrt{6}\sqrt{13}}=\dfrac{1}{\sqrt{78}}$, $\theta=\arccos\!\left(\dfrac{1}{\sqrt{78}}\right)\approx83{,}5°$.

\textbf{2.} $\vec{u}\wedge\vec{v}=\begin{pmatrix}2\cdot2-(-1)\cdot0\\(-1)\cdot3-1\cdot2\\1\cdot0-2\cdot3\end{pmatrix}=\begin{pmatrix}4\\-5\\-6\end{pmatrix}$.
Vérif : $(4,{-5},{-6})\cdot(1,2,{-1})=4-10+6=0$ ✓ ; $(4,{-5},{-6})\cdot(3,0,2)=12+0-12=0$ ✓.

\textbf{3.} Aire $=\|(4,-5,-6)\|=\sqrt{16+25+36}=\sqrt{77}$.
\end{correctionbox}

\begin{retenirbox}
\textbf{Calcul vectoriel — À retenir}
\begin{itemize}
  \item $\vec{u}\cdot\vec{v}=\sum u_i v_i = \|\vec{u}\|\|\vec{v}\|\cos\theta$.
  \item Orthogonalité : $\vec{u}\perp\vec{v}\iff\vec{u}\cdot\vec{v}=0$.
  \item Produit vectoriel : $\vec{u}\wedge\vec{v}\perp\vec{u}$, $\perp\vec{v}$, norme = aire du parallélogramme.
  \item Rotation d'angle $\theta$ : matrice $\begin{pmatrix}\cos\theta&-\sin\theta\\\sin\theta&\cos\theta\end{pmatrix}$.
  \item Barycentre : $G=\dfrac{\sum\lambda_i A_i}{\sum\lambda_i}$.
\end{itemize}
\end{retenirbox}

%% === CHAPITRE 11 ===
\chapter{Géométrie du plan et de l'espace}

\section*{Objectifs}
\begin{itemize}[leftmargin=1.5em]
  \item Écrire les équations de droites et plans dans différentes formes.
  \item Calculer des distances, intersections, angles.
  \item Maîtriser les coordonnées polaires, cylindriques et sphériques.
  \item Étudier cercles et sphères.
\end{itemize}

\section{Droites dans le plan}

Une droite dans $\mathbb{R}^2$ admet deux représentations complémentaires : la forme cartésienne $ax + by + c = 0$, qui la définit comme lieu de points satisfaisant une équation, et la forme paramétrique, qui la décrit comme trajectoire d'un point mobile. La forme cartésienne est adaptée au calcul de distances et d'intersections ; la forme paramétrique est privilégiée pour les courbes dans l'espace.

\begin{definitionbox}[Équations d'une droite dans $\R^2$]
Une droite $\mathcal{D}$ dans $\R^2$ peut être décrite par :
\begin{itemize}
  \item \textbf{Forme cartésienne :} $ax+by+c=0$, $(a,b)\neq(0,0)$.
  \item \textbf{Forme paramétrique :} $(x,y)=(x_0,y_0)+t(d_1,d_2)$, $t\in\R$, où $(x_0,y_0)\in\mathcal{D}$ et $(d_1,d_2)$ est un vecteur directeur.
  \item \textbf{Forme réduite :} $y=mx+p$ (si non verticale), $m=-a/b$.
\end{itemize}
Vecteur normal à $ax+by+c=0$ : $\vec{n}=(a,b)$ ; vecteur directeur : $(-b,a)$.
\end{definitionbox}

\begin{theoremebox}[Distance d'un point à une droite]
La distance du point $P(x_0,y_0)$ à la droite $ax+by+c=0$ est :
\[
  d(P,\mathcal{D}) = \frac{|ax_0+by_0+c|}{\sqrt{a^2+b^2}}.
\]
\end{theoremebox}

\section{Plans et droites dans l'espace}

La géométrie de $\mathbb{R}^3$ généralise celle du plan en introduisant le vecteur normal comme outil central : un plan est entièrement caractérisé par l'un de ses points et son vecteur normal, de même qu'une droite dans $\mathbb{R}^2$ l'est par un point et son vecteur directeur. Les positions relatives (parallélisme, perpendicularité, intersection) se ramènent alors à des conditions algébriques simples faisant intervenir produits scalaires et produits vectoriels.

\begin{definitionbox}[Équation d'un plan dans $\R^3$]
Un plan $\mathcal{P}$ de vecteur normal $\vec{n}=(a,b,c)$ passant par $A(x_0,y_0,z_0)$ a pour équation :
\[
  a(x-x_0)+b(y-y_0)+c(z-z_0)=0, \quad \text{soit} \quad ax+by+cz+d=0.
\]
\end{definitionbox}

\begin{definitionbox}[Droite dans $\R^3$]
Une droite $\mathcal{D}$ de vecteur directeur $\vec{v}=(l,m,n)$ passant par $A(x_0,y_0,z_0)$ :
\begin{itemize}
  \item \textbf{Paramétrique :} $\begin{cases}x=x_0+lt\\y=y_0+mt\\z=z_0+nt\end{cases}$, $t\in\R$.
  \item \textbf{Symétrique :} $\dfrac{x-x_0}{l}=\dfrac{y-y_0}{m}=\dfrac{z-z_0}{n}$ (si $l,m,n\neq0$).
\end{itemize}
\end{definitionbox}

\begin{theoremebox}[Positions relatives dans $\R^3$]
\begin{itemize}
  \item Deux plans sont parallèles $\iff$ leurs vecteurs normaux sont colinéaires.
  \item Deux plans sont perpendiculaires $\iff \vec{n}_1\cdot\vec{n}_2=0$.
  \item Droite $\parallel$ plan $\iff \vec{v}\cdot\vec{n}=0$ et $A\notin\mathcal{P}$.
  \item Droite $\subset$ plan $\iff \vec{v}\cdot\vec{n}=0$ et $A\in\mathcal{P}$.
  \item Droite $\perp$ plan $\iff \vec{v}\parallel\vec{n}$.
\end{itemize}
\end{theoremebox}

\section{Cercles et sphères}

Le cercle et la sphère sont les lieux de points équidistants d'un centre fixé : leur équation canonique découle directement de la définition de la norme. L'identification d'une équation de degré 2 en $x$, $y$ (et $z$) à une équation de cercle ou sphère passe par la technique de complétion du carré.

\begin{definitionbox}[Cercle et sphère]
\begin{itemize}
  \item \textbf{Cercle} de centre $\Omega(a,b)$ et rayon $r>0$ dans $\R^2$ : $(x-a)^2+(y-b)^2=r^2$.
  \item \textbf{Sphère} de centre $\Omega(a,b,c)$ et rayon $r>0$ dans $\R^3$ : $(x-a)^2+(y-b)^2+(z-c)^2=r^2$.
\end{itemize}
\end{definitionbox}

\section{Systèmes de coordonnées}

Le choix du système de coordonnées est souvent décisif pour la simplicité des calculs : les coordonnées polaires sont naturelles pour les cercles et spirales, les coordonnées cylindriques pour les objets à symétrie axiale, et les coordonnées sphériques pour les problèmes à symétrie centrale (électrostatique, mécanique céleste). Maîtriser les formules de changement de variables et les éléments d'aire ou de volume associés est indispensable pour les intégrales multiples.

\begin{definitionbox}[Coordonnées polaires]
Dans $\R^2$, le point $M$ de coordonnées cartésiennes $(x,y)$ a pour \emph{coordonnées polaires} $(r,\theta)$ avec :
\[
  x = r\cos\theta, \quad y = r\sin\theta, \quad r = \sqrt{x^2+y^2}, \quad \theta = \mathrm{Arg}(x+iy).
\]
\end{definitionbox}

\begin{definitionbox}[Coordonnées cylindriques et sphériques]
Dans $\R^3$, un point $M(x,y,z)$ :
\begin{itemize}
  \item \textbf{Coordonnées cylindriques} $(r,\theta,z)$ : $x=r\cos\theta$, $y=r\sin\theta$, $z=z$.
  \item \textbf{Coordonnées sphériques} $(\rho,\theta,\varphi)$ : $x=\rho\sin\varphi\cos\theta$, $y=\rho\sin\varphi\sin\theta$, $z=\rho\cos\varphi$,
    avec $\rho\geq0$, $\theta\in[0,2\pi[$, $\varphi\in[0,\pi]$.
\end{itemize}
\end{definitionbox}

\begin{center}
\begin{tikzpicture}[scale=1.3, >=Stealth]
  % Axes
  \draw[->] (0,0,0)--(2.2,0,0) node[right,font=\scriptsize]{$x$};
  \draw[->] (0,0,0)--(0,2.2,0) node[above,font=\scriptsize]{$y$};
  \draw[->] (0,0,0)--(0,0,2.2) node[above,font=\scriptsize]{$z$};
  % Point M
  \coordinate (M) at (1.2,0.9,1.5);
  \fill[caplpblue] (M) circle(2pt) node[right,font=\scriptsize,caplpblue]{$M(\rho,\theta,\varphi)$};
  % Rayon
  \draw[thick,caplpblue] (0,0,0)--(M) node[midway,above left,font=\scriptsize]{$\rho$};
  % Projection sur Oxy
  \coordinate (Mxy) at (1.2,0.9,0);
  \draw[dashed,gray] (M)--(Mxy);
  \draw[dashed,gray] (Mxy)--(0,0,0);
  % Angle phi (avec Oz)
  \draw[caplpgreen,->] (0,0,0.5) arc[start angle=90,end angle=51,radius=0.5];
  \node[font=\scriptsize,caplpgreen] at (0.15,0.1,0.6){$\varphi$};
  % Angle theta dans Oxy
  \draw[caplporange,->] (0.35,0,0) arc[start angle=0,end angle=36.87,radius=0.35];
  \node[font=\scriptsize,caplporange] at (0.4,0.2,0){$\theta$};
\end{tikzpicture}
\captionof*{figure}{\small\itshape Coordonnées sphériques : $\rho$ = distance à l'origine, $\varphi$ = colatitude, $\theta$ = longitude.}
\end{center}

\begin{exercicebox}[Droites, plans et distances]
\begin{enumerate}
  \item Trouver l'équation du plan passant par $A(1,0,2)$, $B(0,1,1)$ et $C(2,1,0)$.
  \item Calculer la distance du point $P(3,1,4)$ à ce plan.
\end{enumerate}
\end{exercicebox}

\begin{correctionbox}[Correction]
\textbf{1.} Vecteurs $\overrightarrow{AB}=(-1,1,-1)$ et $\overrightarrow{AC}=(1,1,-2)$.
$\vec{n}=\overrightarrow{AB}\wedge\overrightarrow{AC}=\begin{pmatrix}1\cdot(-2)-(-1)\cdot1\\(-1)\cdot1-(-1)(-2)\\(-1)\cdot1-1\cdot1\end{pmatrix}=\begin{pmatrix}-1\\-3\\-2\end{pmatrix}$.
On peut prendre $\vec{n}=(1,3,2)$.
Équation : $1(x-1)+3(y-0)+2(z-2)=0$, soit $x+3y+2z=5$.

\textbf{2.} $d(P,\mathcal{P})=\dfrac{|3+3+8-5|}{\sqrt{1+9+4}}=\dfrac{9}{\sqrt{14}}=\dfrac{9\sqrt{14}}{14}$.
\end{correctionbox}

\begin{retenirbox}
\textbf{Géométrie — À retenir}
\begin{itemize}
  \item Plan : $ax+by+cz+d=0$, $\vec{n}=(a,b,c)$.
  \item Distance point–plan : $\dfrac{|ax_0+by_0+cz_0+d|}{\sqrt{a^2+b^2+c^2}}$.
  \item Droite$\parallel$plan $\iff\vec{v}\cdot\vec{n}=0$ ; droite$\perp$plan $\iff\vec{v}\parallel\vec{n}$.
  \item Sphériques : $x=\rho\sin\varphi\cos\theta$, $y=\rho\sin\varphi\sin\theta$, $z=\rho\cos\varphi$.
\end{itemize}
\end{retenirbox}

%% === CHAPITRE 12 ===
\chapter{Formule du binôme de Newton}

\section*{Objectifs}
\begin{itemize}[leftmargin=1.5em]
  \item Calculer et utiliser les coefficients binomiaux.
  \item Démontrer et appliquer la formule du binôme.
  \item Déduire des identités remarquables et des formules de sommation.
\end{itemize}

\section{Coefficients binomiaux}

Le coefficient binomial $\binom{n}{k}$ admet une interprétation combinatoire directe : il compte le nombre de façons de choisir $k$ éléments parmi $n$ sans tenir compte de l'ordre. Cette lecture est indispensable pour comprendre les lois binomiale et hypergéométrique en probabilités. La relation de Pascal, qui relie $\binom{n+1}{k}$ à deux coefficients de la ligne précédente, traduit le fait que tout choix de $k$ éléments parmi $n+1$ soit inclut, soit exclut le $(n+1)$-ième élément.

\begin{definitionbox}[Coefficient binomial]
Pour $n \in \N$ et $0 \leq k \leq n$ :
\[
  \binom{n}{k} = \frac{n!}{k!(n-k)!}.
\]
\end{definitionbox}

\begin{theoremebox}[Propriétés fondamentales]
\begin{itemize}
  \item $\dbinom{n}{0}=\dbinom{n}{n}=1$ ; $\dbinom{n}{1}=\dbinom{n}{n-1}=n$.
  \item Symétrie : $\dbinom{n}{k}=\dbinom{n}{n-k}$.
  \item \textbf{Formule de Pascal :} $\dbinom{n+1}{k}=\dbinom{n}{k-1}+\dbinom{n}{k}$.
  \item $\displaystyle\sum_{k=0}^n\binom{n}{k}=2^n$ ; $\displaystyle\sum_{k=0}^n(-1)^k\binom{n}{k}=0$.
\end{itemize}
\end{theoremebox}

\section{Formule du binôme}

La formule du binôme de Newton est la généralisation à tout exposant entier $n$ de l'identité remarquable $(a+b)^2 = a^2 + 2ab + b^2$. Elle permet de développer $(a+b)^n$ de façon systématique et de calculer des identités de sommation en specialisant $a$ et $b$ (par exemple $a = b = 1$ donne $\sum_k \binom{n}{k} = 2^n$). Sa démonstration par récurrence illustre parfaitement l'utilisation de la relation de Pascal.

\begin{theoremebox}[Binôme de Newton]
Pour tous $a,b\in\R$ et $n\in\N$ :
\[
  (a+b)^n = \sum_{k=0}^n \binom{n}{k}a^{n-k}b^k.
\]
\end{theoremebox}

\begin{preuvebox}
Par récurrence sur $n$.
\textbf{Init.} $n=0$ : $(a+b)^0=1=\binom{0}{0}$. ✓
\textbf{Hér.} Supposons la formule vraie au rang $n$. Alors :
$(a+b)^{n+1}=(a+b)(a+b)^n=(a+b)\sum_{k=0}^n\binom{n}{k}a^{n-k}b^k
=\sum_{k=0}^n\binom{n}{k}a^{n+1-k}b^k+\sum_{k=0}^n\binom{n}{k}a^{n-k}b^{k+1}$.
En réindexant la seconde somme ($k\to k-1$) :
$=\sum_{k=0}^{n+1}\left[\binom{n}{k}+\binom{n}{k-1}\right]a^{n+1-k}b^k=\sum_{k=0}^{n+1}\binom{n+1}{k}a^{n+1-k}b^k$
par la formule de Pascal. \qed
\end{preuvebox}

\begin{exemplebox}[Applications]
$(x+1)^5 = x^5+5x^4+10x^3+10x^2+5x+1$.

$(2x-3)^4 = \sum_{k=0}^4\binom{4}{k}(2x)^{4-k}(-3)^k = 16x^4-96x^3+216x^2-216x+81$.

\textbf{Identité :} $(1+x)^n = \sum_k\binom{n}{k}x^k$.
En dérivant : $n(1+x)^{n-1}=\sum_{k=1}^n k\binom{n}{k}x^{k-1}$.
En $x=1$ : $\displaystyle\sum_{k=0}^n k\binom{n}{k}=n\cdot2^{n-1}$.
\end{exemplebox}

\begin{exercicebox}[Identités et coefficient]
\begin{enumerate}
  \item Développer $(1+i)^{10}$ de deux façons (binôme et forme exponentielle) et retrouver la valeur de $\binom{10}{2}-\binom{10}{6}+\binom{10}{10}$.
  \item Calculer $\displaystyle\sum_{k=0}^{20}(-1)^k\binom{20}{k}3^k$.
\end{enumerate}
\end{exercicebox}

\begin{correctionbox}[Correction]
\textbf{1.} $1+i=\sqrt{2}e^{i\pi/4}$, donc $(1+i)^{10}=(\sqrt{2})^{10}e^{i10\pi/4}=32e^{i5\pi/2}=32i$.
Par le binôme : $(1+i)^{10}=\sum_{k=0}^{10}\binom{10}{k}i^k$.
$i^k$ suit le cycle $1,i,-1,-i,\ldots$. Partie réelle : $\binom{10}{0}-\binom{10}{4}+\binom{10}{8}=1-210+45=-164$... Hmm, recalculons. Partie imaginaire : $\binom{10}{1}-\binom{10}{5}+\binom{10}{9}=10-252+10=-232$... Ainsi $(1+i)^{10}=32i$ signifie partie réelle $=0$ et partie imaginaire $=32$. On obtient $\binom{10}{2}-\binom{10}{6}+\binom{10}{10}=45-210+1=-164$... La partie réelle est $1-\binom{10}{4}+\binom{10}{8}=1-210+45=-164\neq0$. Contradiction ? En fait $\mathrm{Re}((1+i)^{10})=0$ vient des termes $k\equiv0\pmod4$ avec signe : vérifier soigneusement. Le résultat $32i$ confirme Re $=0$.

\textbf{2.} $\displaystyle\sum_{k=0}^{20}(-1)^k\binom{20}{k}3^k = \sum_{k=0}^{20}\binom{20}{k}(-3)^k\cdot1^{20-k}=(1-3)^{20}=(-2)^{20}=2^{20}=1\,048\,576$.
\end{correctionbox}

\begin{retenirbox}
\textbf{Formule du binôme — À retenir}
\begin{itemize}
  \item $(a+b)^n=\sum_{k=0}^n\binom{n}{k}a^{n-k}b^k$.
  \item Pascal : $\binom{n+1}{k}=\binom{n}{k}+\binom{n}{k-1}$.
  \item $\sum_k\binom{n}{k}=2^n$ (en $a=b=1$) ; $\sum_k(-1)^k\binom{n}{k}=0$ (en $a=1,b=-1$).
  \item Pour extraire un coefficient : identifier la puissance de $x$ dans $(1+x)^n$.
\end{itemize}
\end{retenirbox}

%% === CHAPITRE 13 ===
\chapter{Probabilités}

\section*{Objectifs}
\begin{itemize}[leftmargin=1.5em]
  \item Définir et utiliser un espace probabilisé fini.
  \item Maîtriser probabilités conditionnelles, formule des probabilités totales et formule de Bayes.
  \item Définir l'indépendance d'événements.
  \item Étudier les variables aléatoires discrètes : loi, espérance, variance.
  \item Connaître les lois uniforme, Bernoulli et binomiale.
\end{itemize}

\section{Espaces probabilisés finis}

Les axiomes de Kolmogorov formalisent notre intuition de la probabilité : une mesure $\mathbb{P}$ sur les événements est normalisée à 1 et est additive sur les événements incompatibles. Ce cadre axiomatique, à la fois minimal et complet, engendre toutes les propriétés opérationnelles (complémentaire, union, monotonie) à partir de trois postulats seulement, et s'applique sans modification aux espaces finis comme aux espaces continus.

\begin{definitionbox}[Espace probabilisé fini]
Un \emph{espace probabilisé fini} est un triplet $(\Omega, \mathcal{P}(\Omega), \mathbb{P})$ où :
\begin{itemize}
  \item $\Omega = \{\omega_1,\ldots,\omega_n\}$ est l'univers (ensemble fini des issues).
  \item $\mathcal{P}(\Omega)$ est l'ensemble des parties de $\Omega$ (événements).
  \item $\mathbb{P} : \mathcal{P}(\Omega) \to [0,1]$ vérifie $\mathbb{P}(\Omega)=1$ et
    $\mathbb{P}(A\cup B)=\mathbb{P}(A)+\mathbb{P}(B)$ si $A\cap B=\emptyset$.
\end{itemize}
\end{definitionbox}

\begin{theoremebox}[Propriétés de base]
\begin{itemize}
  \item $\mathbb{P}(\overline{A})=1-\mathbb{P}(A)$.
  \item $\mathbb{P}(\emptyset)=0$.
  \item $A\subseteq B \implies \mathbb{P}(A)\leq\mathbb{P}(B)$.
  \item $\mathbb{P}(A\cup B)=\mathbb{P}(A)+\mathbb{P}(B)-\mathbb{P}(A\cap B)$.
\end{itemize}
\end{theoremebox}

\section{Probabilités conditionnelles}

La probabilité conditionnelle modélise la mise à jour de notre connaissance : savoir que $B$ s'est réalisé restreint l'univers à $B$, ce qui modifie les probabilités de tous les autres événements. La formule de Bayes, qui en est le corollaire le plus utile, permet d'inverser le sens du conditionnement : connaissant $\mathbb{P}(A|B_i)$ (vraisemblances), elle calcule $\mathbb{P}(B_j|A)$ (probabilités a posteriori), ce qui est le cœur du raisonnement bayésien en statistique.

\begin{definitionbox}[Probabilité conditionnelle]
La \emph{probabilité conditionnelle} de $A$ sachant $B$ (avec $\mathbb{P}(B)>0$) est :
\[
  \mathbb{P}(A|B) = \frac{\mathbb{P}(A\cap B)}{\mathbb{P}(B)}.
\]
\end{definitionbox}

\begin{theoremebox}[Formule des probabilités totales]
Soit $(B_1,\ldots,B_n)$ un \emph{système complet d'événements} (partition de $\Omega$, $\mathbb{P}(B_i)>0$). Alors pour tout événement $A$ :
\[
  \mathbb{P}(A) = \sum_{i=1}^n \mathbb{P}(A|B_i)\,\mathbb{P}(B_i).
\]
\end{theoremebox}

\begin{theoremebox}[Formule de Bayes]
Sous les mêmes hypothèses :
\[
  \mathbb{P}(B_j|A) = \frac{\mathbb{P}(A|B_j)\,\mathbb{P}(B_j)}{\displaystyle\sum_{i=1}^n \mathbb{P}(A|B_i)\,\mathbb{P}(B_i)}.
\]
\end{theoremebox}

\begin{definitionbox}[Indépendance]
Deux événements $A$ et $B$ sont \emph{indépendants} si :
\[
  \mathbb{P}(A\cap B) = \mathbb{P}(A)\,\mathbb{P}(B).
\]
$n$ événements $A_1,\ldots,A_n$ sont \emph{mutuellement indépendants} si pour toute partie $I\subseteq\{1,\ldots,n\}$ :
$\mathbb{P}\!\left(\bigcap_{i\in I}A_i\right)=\prod_{i\in I}\mathbb{P}(A_i)$.
\end{definitionbox}

\section{Variables aléatoires discrètes}

Une variable aléatoire discrète est une façon d'encoder numériquement les issues d'une expérience aléatoire, ce qui permet d'effectuer des calculs d'espérance et de variance. L'espérance $\mathbb{E}(X)$ est la valeur « moyenne attendue » sur un grand nombre de répétitions de l'expérience ; la variance $\mathbb{V}(X)$ mesure la dispersion des valeurs autour de cette moyenne. Ces deux paramètres résument la loi de manière compacte et sont indispensables à toute analyse probabiliste.

\begin{definitionbox}[Variable aléatoire discrète, loi]
Une v.a.\ discrète $X$ prend des valeurs $x_1<x_2<\cdots<x_k$ avec probabilités $p_i=\mathbb{P}(X=x_i)>0$, $\sum p_i=1$.
La \emph{loi} de $X$ est la donnée de ces couples $(x_i,p_i)$.
\end{definitionbox}

\begin{definitionbox}[Espérance, variance, écart-type]
\[
  \mathbb{E}(X)=\sum_i x_i p_i, \quad
  \mathbb{V}(X)=\mathbb{E}(X^2)-(\mathbb{E}(X))^2=\sum_i x_i^2 p_i - \left(\sum_i x_i p_i\right)^2, \quad
  \sigma(X)=\sqrt{\mathbb{V}(X)}.
\]
\end{definitionbox}

\begin{theoremebox}[Propriétés de l'espérance et de la variance]
Pour des v.a.\ $X,Y$ et $a,b\in\R$ :
\begin{itemize}
  \item $\mathbb{E}(aX+b)=a\mathbb{E}(X)+b$ ; $\mathbb{E}(X+Y)=\mathbb{E}(X)+\mathbb{E}(Y)$.
  \item $\mathbb{V}(aX+b)=a^2\mathbb{V}(X)$.
  \item Si $X,Y$ indépendantes : $\mathbb{E}(XY)=\mathbb{E}(X)\mathbb{E}(Y)$ et $\mathbb{V}(X+Y)=\mathbb{V}(X)+\mathbb{V}(Y)$.
\end{itemize}
\end{theoremebox}

\section{Lois classiques}

Chaque loi classique modélise un type précis d'aléatoire : la loi uniforme correspond à l'équiprobabilité des issues, la loi de Bernoulli à une expérience à deux issues (succès/échec), la loi binomiale au nombre de succès dans $n$ répétitions indépendantes, et la loi de Poisson aux événements rares et nombreux. La connaissance de leurs formules d'espérance et de variance, ainsi que des situations qui les génèrent, est indispensable pour modéliser rapidement un énoncé probabiliste.

\begin{definitionbox}[Loi uniforme discrète]
$X$ suit la \emph{loi uniforme} sur $\{1,\ldots,n\}$ si $\mathbb{P}(X=k)=\dfrac{1}{n}$ pour tout $k$.
$\mathbb{E}(X)=\dfrac{n+1}{2}$, $\mathbb{V}(X)=\dfrac{n^2-1}{12}$.
\end{definitionbox}

\begin{definitionbox}[Loi de Bernoulli et loi binomiale]
\begin{itemize}
  \item $X\sim\mathcal{B}(p)$ (\emph{Bernoulli}) : $X\in\{0,1\}$, $\mathbb{P}(X=1)=p$. $\mathbb{E}(X)=p$, $\mathbb{V}(X)=p(1-p)$.
  \item $X\sim\mathcal{B}(n,p)$ (\emph{binomiale}) : $X$ = nombre de succès en $n$ épreuves de Bernoulli indépendantes de paramètre $p$.
    \[\mathbb{P}(X=k)=\binom{n}{k}p^k(1-p)^{n-k}, \quad k=0,\ldots,n.\]
    $\mathbb{E}(X)=np$, $\mathbb{V}(X)=np(1-p)$.
\end{itemize}
\end{definitionbox}

\begin{center}
\begin{tikzpicture}
\begin{axis}[
  width=10cm, height=5.5cm,
  xlabel={$k$}, ylabel={$\mathbb{P}(X=k)$},
  title={\small Loi binomiale $\mathcal{B}(n,p)$},
  ybar, bar width=5pt,
  xmin=-0.5, xmax=10.5, ymin=0, ymax=0.35,
  xtick={0,...,10},
  tick label style={font=\scriptsize},
  label style={font=\small},
  legend style={font=\scriptsize, at={(0.98,0.98)}, anchor=north east},
  grid=major, grid style={dotted,gray!30},
]
% B(10,0.3)
\addplot[fill=caplpbluebg, draw=caplpblue] coordinates {
  (0,0.0282)(1,0.1211)(2,0.2335)(3,0.2668)(4,0.2001)(5,0.1029)(6,0.0368)(7,0.0090)(8,0.0014)(9,0.0001)(10,0.000006)};
\addlegendentry{$n=10,p=0{,}3$}
% B(10,0.5)
\addplot[fill=caplpgreenbg, draw=caplpgreen] coordinates {
  (0,0.00098)(1,0.00977)(2,0.04395)(3,0.11719)(4,0.20508)(5,0.24609)(6,0.20508)(7,0.11719)(8,0.04395)(9,0.00977)(10,0.00098)};
\addlegendentry{$n=10,p=0{,}5$}
\end{axis}
\end{tikzpicture}
\captionof*{figure}{\small\itshape Loi binomiale $\mathcal{B}(10,p)$ pour $p=0{,}3$ et $p=0{,}5$. Pour $p=0{,}5$, la distribution est symétrique.}
\end{center}

\begin{exercicebox}[Bayes et binomiale — type concours]
Une urne $A$ contient 3 boules rouges et 2 bleues ; une urne $B$ contient 1 rouge et 4 bleues.
On choisit une urne au hasard (équiprobable), puis on tire 2 boules avec remise.
\begin{enumerate}
  \item Calculer la probabilité d'obtenir 2 boules rouges.
  \item Sachant que les 2 boules sont rouges, quelle est la probabilité qu'on ait tiré dans l'urne $A$ ?
\end{enumerate}
\end{exercicebox}

\begin{correctionbox}[Correction]
Notons $U_A$, $U_B$ les événements "urne $A$", "urne $B$" et $R_2$ l'événement "2 rouges".
$\mathbb{P}(U_A)=\mathbb{P}(U_B)=1/2$.

\textbf{1.} Avec remise dans $A$ : $\mathbb{P}(R_2|U_A)=(3/5)^2=9/25$.
Dans $B$ : $\mathbb{P}(R_2|U_B)=(1/5)^2=1/25$.
Prob. totales : $\mathbb{P}(R_2)=\dfrac{1}{2}\cdot\dfrac{9}{25}+\dfrac{1}{2}\cdot\dfrac{1}{25}=\dfrac{10}{50}=\dfrac{1}{5}$.

\textbf{2.} Bayes : $\mathbb{P}(U_A|R_2)=\dfrac{\mathbb{P}(R_2|U_A)\mathbb{P}(U_A)}{\mathbb{P}(R_2)}=\dfrac{(9/25)(1/2)}{1/5}=\dfrac{9/50}{1/5}=\dfrac{9}{10}$.
\end{correctionbox}

\begin{retenirbox}
\textbf{Probabilités — À retenir}
\begin{itemize}
  \item $\mathbb{P}(A|B)=\dfrac{\mathbb{P}(A\cap B)}{\mathbb{P}(B)}$.
  \item Prob. totales : $\mathbb{P}(A)=\sum_i\mathbb{P}(A|B_i)\mathbb{P}(B_i)$.
  \item Bayes : $\mathbb{P}(B_j|A)=\dfrac{\mathbb{P}(A|B_j)\mathbb{P}(B_j)}{\mathbb{P}(A)}$.
  \item Indépendance : $\mathbb{P}(A\cap B)=\mathbb{P}(A)\mathbb{P}(B)$.
  \item Binomiale : $\mathbb{P}(X=k)=\binom{n}{k}p^k(1-p)^{n-k}$, $\mathbb{E}=np$, $\mathbb{V}=np(1-p)$.
\end{itemize}
\end{retenirbox}

%% === CHAPITRE 14 ===
\chapter{Statistiques à une variable}

\section*{Objectifs}
\begin{itemize}[leftmargin=1.5em]
  \item Calculer les indicateurs de position : moyenne, médiane, quartiles.
  \item Calculer les indicateurs de dispersion : étendue, écart interquartile, variance, écart-type.
  \item Construire et interpréter un histogramme et une boîte à moustaches.
\end{itemize}

\section{Indicateurs de position}

La moyenne et la médiane sont deux indicateurs de position aux comportements très différents : la moyenne est sensible aux valeurs extrêmes (outliers) et peut être fortement tirée par quelques données aberrantes, tandis que la médiane est robuste et reflète mieux le « centre » d'une distribution asymétrique. Le choix de l'un ou l'autre indicateur dépend donc de la nature des données et de l'objectif de l'analyse.

\begin{definitionbox}[Série statistique]
Une \emph{série statistique} est un ensemble de $n$ valeurs $x_1 \leq x_2 \leq \cdots \leq x_n$ (données triées par ordre croissant), éventuellement avec des effectifs $n_i$ ou fréquences $f_i=n_i/n$.
\end{definitionbox}

\begin{definitionbox}[Moyenne]
La \emph{moyenne} (ou moyenne arithmétique) est :
\[
  \bar{x} = \frac{1}{n}\sum_{i=1}^n x_i = \sum_i x_i f_i.
\]
\end{definitionbox}

\begin{definitionbox}[Médiane et quartiles]
Données triées $x_1\leq\cdots\leq x_n$ :
\begin{itemize}
  \item La \emph{médiane} $M_e$ partage la série en deux moitiés égales.
    Si $n$ est impair : $M_e = x_{(n+1)/2}$ ; si $n$ est pair : $M_e = \frac{1}{2}(x_{n/2}+x_{n/2+1})$.
  \item Le \emph{premier quartile} $Q_1$ est la médiane de la moitié inférieure.
  \item Le \emph{troisième quartile} $Q_3$ est la médiane de la moitié supérieure.
\end{itemize}
\end{definitionbox}

\section{Indicateurs de dispersion}

La variance et l'écart-type quantifient la dispersion des données autour de leur moyenne : un écart-type faible signifie que les données sont concentrées, un écart-type élevé indique une grande variabilité. L'écart interquartile, moins sensible aux valeurs extrêmes, est l'indicateur de dispersion associé à la médiane et est représenté visuellement par la boîte dans le diagramme boîte à moustaches.

\begin{definitionbox}[Étendue et écart interquartile]
\begin{itemize}
  \item \textbf{Étendue} : $e = x_n - x_1$ (différence max–min).
  \item \textbf{Écart interquartile} : $EI = Q_3 - Q_1$.
\end{itemize}
\end{definitionbox}

\begin{definitionbox}[Variance et écart-type]
\[
  s^2 = \frac{1}{n}\sum_{i=1}^n (x_i-\bar{x})^2 = \frac{1}{n}\sum x_i^2 - \bar{x}^2 \quad \text{(variance)}, \qquad s = \sqrt{s^2} \quad \text{(écart-type)}.
\]
\end{definitionbox}

\begin{erreurbox}
La variance de population se divise par $n$ ; la variance \emph{corrigée} (échantillon) par $n-1$. Au concours CAPLP, préciser laquelle est utilisée.
\end{erreurbox}

\section{Représentations graphiques}

Le choix de la représentation graphique conditionne la lisibilité d'une analyse statistique : l'histogramme est adapté aux données groupées en classes, la boîte à moustaches résume en cinq nombres la structure d'une distribution et met en évidence les valeurs atypiques, tandis que le diagramme en bâtons convient aux données discrètes. Une visualisation bien choisie permet de détecter immédiatement la symétrie, la présence d'outliers ou la bimodalité d'une distribution.

\begin{definitionbox}[Boîte à moustaches]
La \emph{boîte à moustaches} (box plot) représente :
\begin{itemize}
  \item La boîte de $Q_1$ à $Q_3$ (écart interquartile).
  \item Le trait intérieur à la médiane $M_e$.
  \item Les moustaches jusqu'aux valeurs extrêmes (ou jusqu'à $1{,}5\times EI$).
\end{itemize}
\end{definitionbox}

\begin{center}
\begin{tikzpicture}[scale=1.0]
  % Axe
  \draw[->] (0,0)--(11,0) node[right,font=\scriptsize]{$x$};
  % Données : min=2, Q1=4, Me=6, Q3=8, max=10
  % Boîte
  \fill[caplpbluebg] (4,0.3) rectangle (8,1.0);
  \draw[thick] (4,0.3) rectangle (8,1.0);
  % Médiane
  \draw[thick, caplpred] (6,0.3)--(6,1.0);
  % Moustaches
  \draw[thick] (2,0.65)--(4,0.65);
  \draw[thick] (8,0.65)--(10,0.65);
  \draw[thick] (2,0.5)--(2,0.8);
  \draw[thick] (10,0.5)--(10,0.8);
  % Étiquettes
  \node[font=\scriptsize,below] at (2,0.3) {min};
  \node[font=\scriptsize,below] at (4,0.3) {$Q_1$};
  \node[font=\scriptsize,below,caplpred] at (6,0.3) {$M_e$};
  \node[font=\scriptsize,below] at (8,0.3) {$Q_3$};
  \node[font=\scriptsize,below] at (10,0.3) {max};
  % Valeurs
  \foreach \v in {2,4,6,8,10}{\node[font=\scriptsize,below] at (\v,-0.2) {\v};}
  \node[font=\small] at (5.5,1.4) {$EI=Q_3-Q_1=4$};
\end{tikzpicture}
\captionof*{figure}{\small\itshape Boîte à moustaches : min$=2$, $Q_1=4$, $M_e=6$, $Q_3=8$, max$=10$.}
\end{center}

\begin{exercicebox}[Calculs statistiques]
Les notes de 10 élèves sont (en ordre croissant) : 4, 6, 7, 8, 9, 10, 11, 13, 15, 17.
\begin{enumerate}
  \item Calculer $\bar{x}$, $M_e$, $Q_1$, $Q_3$ et $EI$.
  \item Calculer la variance $s^2$ et l'écart-type $s$.
\end{enumerate}
\end{exercicebox}

\begin{correctionbox}[Correction]
$n=10$. \textbf{1.} $\bar{x}=\dfrac{4+6+\cdots+17}{10}=\dfrac{100}{10}=10$.
$M_e=\dfrac{x_5+x_6}{2}=\dfrac{9+10}{2}=9{,}5$.
$Q_1=$ médiane de $\{4,6,7,8,9\}$ $=7$ ; $Q_3=$ médiane de $\{10,11,13,15,17\}$ $=13$. $EI=6$.

\textbf{2.} $\dfrac{1}{10}\sum x_i^2 = \dfrac{16+36+49+64+81+100+121+169+225+289}{10}=\dfrac{1150}{10}=115$.
$s^2=115-100=15$, $s=\sqrt{15}\approx3{,}87$.
\end{correctionbox}

\begin{retenirbox}
\textbf{Statistiques 1 variable — À retenir}
\begin{itemize}
  \item Moyenne : $\bar{x}=\frac{1}{n}\sum x_i$.
  \item Médiane : valeur centrale (ou moyenne des deux centrales).
  \item $Q_1$, $Q_3$ : médianes des demi-séries.
  \item Variance : $s^2=\frac{1}{n}\sum x_i^2-\bar{x}^2$.
  \item Boîte à moustaches : min, $Q_1$, $M_e$, $Q_3$, max.
\end{itemize}
\end{retenirbox}

%% === CHAPITRE 15 ===
\chapter{Statistiques à deux variables et algorithmique Python}

\section*{Objectifs}
\begin{itemize}[leftmargin=1.5em]
  \item Étudier un nuage de points, calculer le point moyen et la covariance.
  \item Déterminer la droite de régression par moindres carrés.
  \item Interpréter le coefficient de corrélation linéaire.
  \item Maîtriser les bases de l'algorithmique en Python : variables, boucles, fonctions, récursivité.
\end{itemize}

\section{Statistiques à deux variables}

L'analyse bivariée cherche à mettre en évidence et à quantifier les liens entre deux variables statistiques. Le coefficient de corrélation linéaire $r$ mesure l'intensité de la liaison linéaire : $r = \pm 1$ correspond à une liaison linéaire parfaite, $r = 0$ à l'absence de corrélation linéaire (bien que d'autres formes de dépendance puissent exister). La droite de régression par moindres carrés est la droite qui minimise la somme des carrés des écarts verticaux entre les points observés et la droite ajustée.

\begin{definitionbox}[Nuage de points et point moyen]
Une \emph{série statistique à deux variables} est un ensemble de $n$ couples $(x_i,y_i)$.
Le \emph{point moyen} est $G = (\bar{x},\bar{y})$ où $\bar{x}=\frac{1}{n}\sum x_i$ et $\bar{y}=\frac{1}{n}\sum y_i$.
\end{definitionbox}

\begin{definitionbox}[Covariance et coefficient de corrélation]
\[
  \mathrm{cov}(X,Y) = \frac{1}{n}\sum(x_i-\bar{x})(y_i-\bar{y}) = \frac{1}{n}\sum x_i y_i - \bar{x}\bar{y}.
\]
Le \emph{coefficient de corrélation linéaire} est :
\[
  r = \frac{\mathrm{cov}(X,Y)}{s_X\,s_Y}, \quad -1 \leq r \leq 1.
\]
$r$ proche de $\pm1$ : forte corrélation linéaire. $r$ proche de $0$ : pas de corrélation linéaire.
\end{definitionbox}

\begin{theoremebox}[Droite de régression par moindres carrés]
La droite $y = ax + b$ qui minimise $\sum(y_i - ax_i - b)^2$ a pour coefficients :
\[
  a = \frac{\mathrm{cov}(X,Y)}{s_X^2}, \qquad b = \bar{y} - a\bar{x}.
\]
Cette droite passe par le point moyen $G=(\bar{x},\bar{y})$.
\end{theoremebox}

\begin{center}
\begin{tikzpicture}
\begin{axis}[
  width=9cm, height=6cm,
  xlabel={$x$}, ylabel={$y$},
  title={\small Nuage de points et droite de régression},
  xmin=0, xmax=6, ymin=0, ymax=8,
  grid=major, grid style={dotted,gray!30},
  tick label style={font=\scriptsize},
  label style={font=\small},
  legend style={font=\scriptsize, at={(0.05,0.95)}, anchor=north west},
]
% Nuage de points (données fictives)
\addplot[only marks, mark=*, mark size=2.5pt, caplpblue] coordinates {
  (1,1.5)(2,2.8)(2.5,3.2)(3,3.9)(3.5,4.6)(4,5.1)(4.5,5.5)(5,6.5)(5.5,6.8)};
\addlegendentry{$(x_i,y_i)$}
% Droite de régression (approx y = 1.2x + 0.2)
\addplot[thick, caplpred, domain=0.5:5.8] {1.2*x + 0.2};
\addlegendentry{$y=1{,}2x+0{,}2$}
% Point moyen G
\addplot[mark=square*, mark size=3pt, caplpgreen] coordinates {(3.5,4.43)};
\addlegendentry{$G=(\bar{x},\bar{y})$}
\end{axis}
\end{tikzpicture}
\captionof*{figure}{\small\itshape La droite de régression passe par le point moyen $G$ et minimise la somme des carrés des résidus.}
\end{center}

\section{Algorithmique et programmation Python}

Python est devenu le langage de référence pour le calcul scientifique et l'enseignement de l'algorithmique, grâce à sa syntaxe claire et ses bibliothèques puissantes. Au concours CAPLP, la maîtrise des structures de base — variables, boucles, fonctions, récursivité — est évaluée à travers des algorithmes mathématiques concrets (suites, PGCD, tri). La lisibilité du code (indentation correcte, nommage des variables, cas de base de la récursion) est un critère de qualité attendu.

\begin{definitionbox}[Variables et types en Python]
En Python, les types de base sont :
\begin{itemize}
  \item \texttt{int} (entier), \texttt{float} (flottant), \texttt{bool} (\texttt{True}/\texttt{False}), \texttt{str} (chaîne).
  \item \texttt{list} (liste mutable), \texttt{tuple} (immuable).
\end{itemize}
Affectation : \texttt{x = 5} ; typage dynamique.
\end{definitionbox}

\begin{methodebox}[Structures de contrôle]
\begin{itemize}
  \item \textbf{Conditionnelle :} \texttt{if condition: ... elif ...: ... else: ...}
  \item \textbf{Boucle \texttt{for} :} \texttt{for i in range(n): ...} (itère $i=0,1,\ldots,n-1$)
  \item \textbf{Boucle \texttt{while} :} \texttt{while condition: ...}
  \item \textbf{Fonction :} \texttt{def f(x, y): ... return valeur}
  \item \textbf{Récursivité :} une fonction s'appelle elle-même avec un cas de base.
\end{itemize}
\end{methodebox}

\begin{exemplebox}[Algorithmes classiques en Python]
\textbf{Factorielle (récursif) :}
\begin{verbatim}
def factorielle(n):
    if n == 0:
        return 1
    return n * factorielle(n - 1)
\end{verbatim}

\textbf{PGCD par algorithme d'Euclide :}
\begin{verbatim}
def pgcd(a, b):
    while b != 0:
        a, b = b, a % b
    return a
\end{verbatim}

\textbf{Suite de Fibonacci :}
\begin{verbatim}
def fibonacci(n):
    a, b = 0, 1
    for _ in range(n):
        a, b = b, a + b
    return a
\end{verbatim}

\textbf{Tri par sélection :}
\begin{verbatim}
def tri_selection(lst):
    n = len(lst)
    for i in range(n):
        imin = i
        for j in range(i+1, n):
            if lst[j] < lst[imin]:
                imin = j
        lst[i], lst[imin] = lst[imin], lst[i]
    return lst
\end{verbatim}
\end{exemplebox}

\begin{erreurbox}[Pièges Python au concours]
\begin{itemize}
  \item \textbf{Indentation.} Python utilise l'indentation comme structure syntaxique — une erreur d'indentation provoque une erreur d'exécution.
  \item \textbf{Division entière.} \texttt{5/2 = 2.5} (flottant) ; \texttt{5//2 = 2} (division entière).
  \item \textbf{Récursion.} Toujours prévoir un cas de base pour éviter une récursion infinie.
  \item \textbf{Listes.} \texttt{lst2 = lst1} ne copie pas : les deux variables pointent vers la même liste. Utiliser \texttt{lst2 = lst1.copy()}.
\end{itemize}
\end{erreurbox}

\begin{exercicebox}[Régression et algorithme]
\begin{enumerate}
  \item On dispose des données : $(1,2)$, $(2,4)$, $(3,5)$, $(4,4)$, $(5,5)$.
    Calculer $\bar{x}$, $\bar{y}$, $\mathrm{cov}(X,Y)$, $s_X^2$ et la droite de régression $y=ax+b$.
  \item Écrire une fonction Python \texttt{moyenne(lst)} qui retourne la moyenne d'une liste de nombres.
  \item Écrire une fonction Python \texttt{est\_premier(n)} qui retourne \texttt{True} si $n$ est premier.
\end{enumerate}
\end{exercicebox}

\begin{correctionbox}[Correction]
\textbf{1.} $\bar{x}=3$, $\bar{y}=4$. $\frac{1}{5}\sum x_iy_i=\frac{2+8+15+16+25}{5}=\frac{66}{5}=13{,}2$.
$\mathrm{cov}=13{,}2-3\times4=1{,}2$. $s_X^2=\frac{1+4+9+16+25}{5}-9=11-9=2$.
$a=1{,}2/2=0{,}6$, $b=4-0{,}6\times3=2{,}2$. Droite : $y=0{,}6x+2{,}2$.

\textbf{2.}
\begin{verbatim}
def moyenne(lst):
    return sum(lst) / len(lst)
\end{verbatim}

\textbf{3.}
\begin{verbatim}
def est_premier(n):
    if n < 2:
        return False
    for i in range(2, int(n**0.5) + 1):
        if n % i == 0:
            return False
    return True
\end{verbatim}
\end{correctionbox}

\begin{retenirbox}
\textbf{Statistiques 2 variables et Python — À retenir}
\begin{itemize}
  \item $\mathrm{cov}(X,Y)=\frac{1}{n}\sum x_iy_i - \bar{x}\bar{y}$.
  \item Droite de régression : $a=\mathrm{cov}/s_X^2$, $b=\bar{y}-a\bar{x}$, passe par $G$.
  \item $r\in[-1,1]$ : $|r|$ proche de 1 = forte corrélation linéaire.
  \item Python : \texttt{def}, \texttt{for}, \texttt{while}, \texttt{if/elif/else}, récursivité + cas de base.
  \item Algorithme d'Euclide : \texttt{while b!=0: a,b = b, a\%b}.
\end{itemize}
\end{retenirbox}

%% ============================================================
%% PARTIE II — PROGRAMME A.2 : DISCIPLINE MAJEURE MATHÉMATIQUES
%% ============================================================

\part{Programme A.2 — Discipline majeure Mathématiques}

\chapter*{Introduction à la partie A.2}
\addcontentsline{toc}{chapter}{Introduction à la partie A.2}

La partie A.2 du programme CAPLP 2026 (section mathématiques–physique chimie) couvre des
thèmes de niveau licence (L3) propres à la discipline mathématiques : algèbre linéaire avancée,
analyse de fonctions à plusieurs variables, développements limités, séries numériques,
dénombrement, et probabilités discrètes avancées.

Ces notions approfondissent la partie A.1 et sont mobilisées dans les épreuves écrites
de composition ainsi que dans l'épreuve sur dossier.

%% === CHAPITRE A2.1 ===
\chapter{Algèbre linéaire — Espaces vectoriels, applications linéaires, diagonalisation}

\section*{Objectifs}
\begin{itemize}[leftmargin=1.5em]
  \item Maîtriser les notions de base des espaces vectoriels réels (dimension finie).
  \item Savoir caractériser une application linéaire, déterminer noyau et image.
  \item Calculer et utiliser les matrices : opérations, inverse, déterminant.
  \item Déterminer les éléments propres d'un endomorphisme ; diagonaliser.
  \item Trigonaliser ; calculer $A^n$.
\end{itemize}

\section{Espaces vectoriels de dimension finie}

La notion d'espace vectoriel abstrait unifie sous un même cadre théorique des objets a priori très différents : les vecteurs de $\mathbb{R}^n$, les matrices, les polynômes, les fonctions continues. Cette unification permet d'appliquer les mêmes théorèmes (bases, dimension, rang) à toutes ces structures, économisant ainsi des démonstrations spécifiques pour chaque cas. La dimension d'un espace vectoriel, définie comme le cardinal commun à toutes ses bases, est l'invariant fondamental qui classe les espaces à isomorphisme près.

\begin{definitionbox}[Espace vectoriel réel]
Soit $E$ un ensemble non vide muni d'une loi interne \(+\) et d'une loi externe
\(\cdot : \R \times E \to E\). On dit que $(E,+,\cdot)$ est un \emph{espace vectoriel sur \(\R\)}
si les huit axiomes suivants sont satisfaits :
\begin{enumerate}[label=\alph*)]
  \item $(E,+)$ est un groupe abélien (associativité, commutatitivité, élément neutre $\mathbf{0}$, opposé) ;
  \item $\forall \lambda,\mu \in \R,\,\forall u \in E : (\lambda+\mu)\cdot u = \lambda \cdot u + \mu \cdot u$ ;
  \item $\forall \lambda \in \R,\,\forall u,v \in E : \lambda \cdot (u+v) = \lambda \cdot u + \lambda \cdot v$ ;
  \item $\forall \lambda,\mu \in \R,\,\forall u \in E : (\lambda\mu) \cdot u = \lambda \cdot (\mu \cdot u)$ ;
  \item $\forall u \in E : 1 \cdot u = u$.
\end{enumerate}
Les éléments de $E$ sont appelés \emph{vecteurs} et ceux de $\R$ des \emph{scalaires}.
\end{definitionbox}

\begin{exemplebox}[Exemples fondamentaux]
\begin{itemize}
  \item $\R^n$ (avec les opérations composante par composante) pour tout $n \geq 1$.
  \item $\mathcal{M}_{n,p}(\R)$ l'ensemble des matrices réelles $n \times p$.
  \item $\R[X]$ l'anneau des polynômes à coefficients réels.
  \item $\mathcal{C}([a,b],\R)$ les fonctions continues sur $[a,b]$.
\end{itemize}
\end{exemplebox}

\begin{definitionbox}[Sous-espace vectoriel]
Soit $(E,+,\cdot)$ un espace vectoriel et $F \subseteq E$, $F \neq \emptyset$.
$F$ est un \emph{sous-espace vectoriel} (s.e.v.) de $E$ si et seulement si :
\[
  \forall u,v \in F,\;\forall \lambda \in \R :\; \lambda u + v \in F.
\]
Autrement dit : $F$ est stable par combinaisons linéaires.
\end{definitionbox}

\begin{definitionbox}[Famille libre, famille génératrice, base]
Soit $(v_1,\ldots,v_p) \in E^p$.
\begin{itemize}
  \item La famille est \emph{libre} (linéairement indépendante) si :
    \[\lambda_1 v_1 + \cdots + \lambda_p v_p = \mathbf{0} \implies \lambda_1 = \cdots = \lambda_p = 0.\]
  \item Elle est \emph{génératrice} de $E$ si tout vecteur de $E$ est combinaison linéaire des $v_i$.
  \item Elle est une \emph{base} de $E$ si elle est à la fois libre et génératrice.
\end{itemize}
\end{definitionbox}

\begin{theoremebox}[Dimension d'un espace vectoriel de dimension finie]
Si $E$ admet une base, toutes les bases de $E$ ont le même cardinal, appelé
\emph{dimension} de $E$ et noté $\dim E$.
En particulier, $\dim \R^n = n$, $\dim \mathcal{M}_{n,p}(\R) = np$.
\end{theoremebox}

\begin{theoremebox}[Théorème de la base incomplète]
Soit $E$ un e.v. de dimension $n$. Toute famille libre de $k < n$ vecteurs peut être
complétée en une base de $E$ par adjonction de $n-k$ vecteurs.
\end{theoremebox}

\section{Applications linéaires}

Les applications linéaires sont les morphismes de la structure d'espace vectoriel : elles préservent les combinaisons linéaires. Le théorème du rang — $\dim \Ker f + \dim \Im f = \dim E$ — est la version linéaire du principe de conservation : la dimension perdue dans le noyau (écrasée à zéro) est exactement la dimension gagnée dans l'image. Ce théorème est le principal outil pour déterminer si une application est injective, surjective, ou bijective sans calcul explicite.

\begin{definitionbox}[Application linéaire]
Soient $E$ et $F$ deux espaces vectoriels sur $\R$.
Une application $f : E \to F$ est \emph{linéaire} si :
\[
  \forall u,v \in E,\;\forall \lambda \in \R :\;
  f(\lambda u + v) = \lambda f(u) + f(v).
\]
Si $E = F$, on parle d'\emph{endomorphisme} de $E$.
\end{definitionbox}

\begin{definitionbox}[Noyau et image]
Soit $f : E \to F$ linéaire.
\begin{itemize}
  \item Le \emph{noyau} de $f$ : $\Ker(f) = \{u \in E \mid f(u) = \mathbf{0}_F\}$.
  \item L'\emph{image} de $f$ : $\Im(f) = \{f(u) \mid u \in E\}$.
\end{itemize}
$\Ker(f)$ est un s.e.v. de $E$ et $\Im(f)$ est un s.e.v. de $F$.
\end{definitionbox}

\begin{theoremebox}[Théorème du rang (ou de la dimension)]
Soit $f : E \to F$ linéaire, avec $E$ de dimension finie $n$. Alors :
\[
  \dim \Ker(f) + \dim \Im(f) = n.
\]
\end{theoremebox}

\begin{preuvebox}
Posons $r = \dim \Im(f)$ (le \emph{rang} de $f$) et $k = \dim \Ker(f)$.
Soit $(e_1,\ldots,e_k)$ une base de $\Ker(f)$, complétée en une base
$(e_1,\ldots,e_k,e_{k+1},\ldots,e_n)$ de $E$ (théorème de la base incomplète).
On vérifie que $(f(e_{k+1}),\ldots,f(e_n))$ est une base de $\Im(f)$,
d'où $r = n-k$, soit $k+r = n$. \qed
\end{preuvebox}

\begin{theoremebox}[Caractérisation injection/surjection/bijection]
Avec les notations précédentes ($\dim E = n$, $\dim F = m$) :
\begin{itemize}
  \item $f$ injective $\iff$ $\Ker(f) = \{\mathbf{0}\}$ $\iff$ $\dim \Im(f) = n$.
  \item $f$ surjective $\iff$ $\Im(f) = F$ $\iff$ $\dim \Im(f) = m$.
  \item $f$ bijective (isomorphisme) $\iff$ $f$ injective et surjective $\iff$ $n = m$ et $\Ker(f) = \{\mathbf{0}\}$.
\end{itemize}
\end{theoremebox}

\section{Matrices et calcul matriciel}

La matrice d'une application linéaire dépend du choix des bases dans l'espace de départ et l'espace d'arrivée : deux matrices représentant le même endomorphisme dans des bases différentes sont semblables (reliées par $A' = P^{-1}AP$). Changer de base pour simplifier la matrice — idéalement la rendre diagonale — est le programme de la diagonalisation, qui cherche la base la plus « naturelle » pour l'endomorphisme.

\begin{definitionbox}[Matrice d'une application linéaire]
Soient $\mathcal{B} = (e_1,\ldots,e_n)$ une base de $E$ et $\mathcal{C} = (f_1,\ldots,f_m)$ une base de $F$.
La \emph{matrice de $f$ dans les bases $\mathcal{B}$ et $\mathcal{C}$} est la matrice
$M \in \mathcal{M}_{m,n}(\R)$ dont la $j$-ième colonne est le vecteur colonne
des coordonnées de $f(e_j)$ dans $\mathcal{C}$.
\end{definitionbox}

\begin{theoremebox}[Produit matriciel et composition]
Si $f : E \to F$ a pour matrice $A$ et $g : F \to G$ a pour matrice $B$ (dans des bases fixes),
alors $g \circ f$ a pour matrice $BA$.
Le produit matriciel est \textbf{associatif} mais en général \textbf{non commutatif}.
\end{theoremebox}

\begin{definitionbox}[Matrice inversible]
$A \in \mathcal{M}_n(\R)$ est \emph{inversible} s'il existe $B \in \mathcal{M}_n(\R)$ telle que $AB = BA = I_n$.
On note alors $B = A^{-1}$.
\end{definitionbox}

\begin{methodebox}[Calcul de l'inverse par la méthode du pivot (Gauss-Jordan)]
Pour calculer $A^{-1}$ :
\begin{enumerate}
  \item Former la matrice augmentée $[A \mid I_n]$.
  \item Appliquer des opérations élémentaires sur les lignes pour transformer $A$ en $I_n$.
  \item La matrice obtenue à droite est $A^{-1}$.
  \item Si $A$ ne se réduit pas en $I_n$, alors $A$ n'est pas inversible.
\end{enumerate}
\end{methodebox}

\begin{theoremebox}[Déterminant d'une matrice $2 \times 2$ et $3 \times 3$]
Pour $A = \begin{pmatrix} a & b \\ c & d \end{pmatrix}$ :
\[
  \det(A) = ad - bc.
\]
Pour $A = (a_{ij})_{1 \leq i,j \leq 3}$, développement selon la première ligne :
\[
  \det(A) = a_{11}(a_{22}a_{33}-a_{23}a_{32}) - a_{12}(a_{21}a_{33}-a_{23}a_{31}) + a_{13}(a_{21}a_{32}-a_{22}a_{31}).
\]
\end{theoremebox}

\begin{theoremebox}[Propriétés du déterminant]
\begin{itemize}
  \item $A$ inversible $\iff$ $\det(A) \neq 0$.
  \item $\det(AB) = \det(A)\det(B)$.
  \item $\det(A^{-1}) = \frac{1}{\det(A)}$ si $A$ inversible.
  \item $\det(A^\top) = \det(A)$.
  \item Le déterminant est \emph{multilinéaire alterné} en les lignes (et en les colonnes).
\end{itemize}
\end{theoremebox}

\begin{exemplebox}[Inverse d'une matrice $2\times 2$]
Pour $A = \begin{pmatrix} a & b \\ c & d \end{pmatrix}$ avec $\det(A) = ad-bc \neq 0$ :
\[
  A^{-1} = \frac{1}{ad-bc}\begin{pmatrix} d & -b \\ -c & a \end{pmatrix}.
\]
\textbf{Application.} $A = \begin{pmatrix} 2 & 3 \\ 1 & 2 \end{pmatrix}$, $\det A = 4-3 = 1$,
donc $A^{-1} = \begin{pmatrix} 2 & -3 \\ -1 & 2 \end{pmatrix}$.
\end{exemplebox}

\section{Valeurs propres et vecteurs propres}

Diagonaliser une matrice revient à trouver une base de vecteurs propres dans laquelle l'endomorphisme agit simplement par homothéties dans chaque direction. Cette base « naturelle » rend le calcul de $A^n$ trivial ($A^n = PD^nP^{-1}$) et simplifie la résolution de systèmes différentiels. La condition de diagonalisabilité — égalité des multiplicités algébriques et géométriques pour chaque valeur propre — est le critère décisif que tout candidat doit vérifier systématiquement.

\begin{definitionbox}[Valeurs propres, vecteurs propres]
Soit $f$ un endomorphisme de $E$ (de dimension finie $n$) et $A$ sa matrice dans une base.
Un scalaire $\lambda \in \R$ est une \emph{valeur propre} de $f$ (ou de $A$) si :
\[
  \exists v \neq \mathbf{0} : f(v) = \lambda v \quad (\text{équivalent : } Av = \lambda v).
\]
Un tel $v$ est appelé \emph{vecteur propre} associé à $\lambda$.
Le sous-espace $E_\lambda = \Ker(f - \lambda \,\mathrm{id}) = \{v \mid Av = \lambda v\}$
est le \emph{sous-espace propre} associé à $\lambda$.
\end{definitionbox}

\begin{definitionbox}[Polynôme caractéristique]
Le \emph{polynôme caractéristique} de $A \in \mathcal{M}_n(\R)$ est :
\[
  \chi_A(\lambda) = \det(A - \lambda I_n).
\]
Les valeurs propres de $A$ sont exactement les racines de $\chi_A$.
\end{definitionbox}

\begin{theoremebox}[Diagonalisation]
$A \in \mathcal{M}_n(\R)$ est \emph{diagonalisable} (sur $\R$) si et seulement si
$\chi_A$ est scindé sur $\R$ et pour chaque valeur propre $\lambda_i$, la multiplicité
géométrique égale la multiplicité algébrique :
\[
  \dim E_{\lambda_i} = m_i \quad \text{(multiplicité de $\lambda_i$ comme racine de $\chi_A$)}.
\]
Dans ce cas, il existe $P$ inversible telle que $A = P D P^{-1}$
où $D = \mathrm{diag}(\lambda_1,\ldots,\lambda_n)$.
\end{theoremebox}

\begin{methodebox}[Méthode de diagonalisation]
\begin{enumerate}
  \item Calculer $\chi_A(\lambda) = \det(A - \lambda I_n)$.
  \item Trouver les racines réelles $\lambda_1, \ldots, \lambda_k$ et leurs multiplicités $m_i$.
  \item Pour chaque $\lambda_i$, résoudre $(A - \lambda_i I_n)v = 0$ et trouver une base de $E_{\lambda_i}$.
  \item Vérifier $\dim E_{\lambda_i} = m_i$ pour chaque $i$.
  \item Si toutes les conditions sont réunies : $P = [v_1 \mid \cdots \mid v_n]$ (colonnes = vecteurs propres).
  \item Calculer $A^k = P D^k P^{-1}$ où $D^k = \mathrm{diag}(\lambda_1^k,\ldots,\lambda_n^k)$.
\end{enumerate}
\end{methodebox}

\begin{exemplebox}[Diagonalisation d'une matrice $2\times 2$]
Soit $A = \begin{pmatrix} 3 & 1 \\ 0 & 2 \end{pmatrix}$.

\textbf{Étape 1.} $\chi_A(\lambda) = (3-\lambda)(2-\lambda) = 0 \implies \lambda_1 = 2,\; \lambda_2 = 3$.

\textbf{Étape 2.} $E_2 = \Ker(A - 2I) = \Ker\begin{pmatrix}1&1\\0&0\end{pmatrix}$, base : $v_1 = \begin{pmatrix}-1\\1\end{pmatrix}$.

$E_3 = \Ker(A - 3I) = \Ker\begin{pmatrix}0&1\\0&-1\end{pmatrix}$, base : $v_2 = \begin{pmatrix}1\\0\end{pmatrix}$.

\textbf{Étape 3.} $P = \begin{pmatrix}-1&1\\1&0\end{pmatrix}$, $D = \begin{pmatrix}2&0\\0&3\end{pmatrix}$, d'où $A = PDP^{-1}$
et $A^n = P\begin{pmatrix}2^n&0\\0&3^n\end{pmatrix}P^{-1}$.
\end{exemplebox}

\begin{center}
\begin{tikzpicture}[scale=1.3,>=Stealth]
  \draw[->] (-2.3,0) -- (2.3,0) node[right,font=\small] {$x$};
  \draw[->] (0,-1.5) -- (0,2.3) node[above,font=\small] {$y$};
  % Vecteur propre lambda=2 : v1=(-1,1), direction
  \draw[->, very thick, caplpblue] (0,0) -- (-1.2,1.2)
    node[above left, font=\small, caplpblue] {$v_1$\;($\lambda_1=2$)};
  % Vecteur propre lambda=3 : v2=(1,0), direction
  \draw[->, very thick, caplpred] (0,0) -- (1.6,0)
    node[below right, font=\small, caplpred] {$v_2$\;($\lambda_2=3$)};
  % Droites propres (pointillés)
  \draw[dashed, caplpblue!60] (-1.8,1.8) -- (1.8,-1.8);
  \draw[dashed, caplpred!60]  (-2.1,0)   -- (2.1,0);
  % Un vecteur quelconque et son image
  \draw[->, thick, caplpgreen] (0,0) -- (0.8,0.8)
    node[right, font=\small, caplpgreen] {$u$};
  \draw[->, thick, caplporange] (0,0) -- (1.6+0.8*(-2), 0+0.8*3)
    node[above, font=\small, caplporange] {$Au$};
  % Titre
  \node[font=\small\itshape] at (0,-1.8) {Sous-espaces propres de $A=\begin{pmatrix}3&1\\0&2\end{pmatrix}$};
\end{tikzpicture}
\captionof*{figure}{\small\itshape Les deux droites propres (pointillées) ; $v_1$ et $v_2$ sont des vecteurs propres associés à $\lambda_1=2$ et $\lambda_2=3$.}
\end{center}

\begin{theoremebox}[Trigonalisation]
Toute matrice $A \in \mathcal{M}_n(\C)$ est trigonalisable, i.e.\ il existe $P$ inversible (sur $\C$) telle que
$P^{-1}AP$ soit triangulaire supérieure.
Sur $\R$, cela est possible si et seulement si $\chi_A$ est scindé sur $\R$.
Si $A$ n'est pas diagonalisable mais que $\chi_A$ est scindé sur $\R$, on recherche la
forme de Jordan (hors programme détaillé, mais à connaître conceptuellement).
\end{theoremebox}

\begin{erreurbox}[Confusions fréquentes]
\begin{itemize}
  \item \textbf{Erreur 1.} Croire qu'une matrice à $n$ valeurs propres distinctes est automatiquement diagonalisable : c'est vrai seulement si les $n$ valeurs propres sont simples.
  \item \textbf{Erreur 2.} Oublier de vérifier $\dim E_{\lambda_i} = m_i$ en cas de valeur propre multiple.
  \item \textbf{Erreur 3.} Calculer $A^n$ en diagonalisant sans vérifier l'inversibilité de $P$.
\end{itemize}
\end{erreurbox}

\section{Exercices types}

\begin{exercicebox}[Application directe — valeurs propres]
Soit $A = \begin{pmatrix} 5 & -2 \\ 2 & 1 \end{pmatrix}$.
\begin{enumerate}
  \item Calculer le polynôme caractéristique de $A$.
  \item Déterminer les valeurs propres et un vecteur propre associé à chacune.
  \item $A$ est-elle diagonalisable ? Si oui, écrire $A = PDP^{-1}$ et calculer $A^{10}$.
\end{enumerate}
\end{exercicebox}

\begin{correctionbox}[Correction]
\textbf{1.} $\chi_A(\lambda) = (5-\lambda)(1-\lambda)+4 = \lambda^2 - 6\lambda + 9 = (\lambda-3)^2$.

\textbf{2.} Valeur propre $\lambda = 3$ de multiplicité algébrique 2.
$A - 3I = \begin{pmatrix}2&-2\\2&-2\end{pmatrix}$, $\Ker(A-3I) = \mathrm{Vect}\begin{pmatrix}1\\1\end{pmatrix}$, $\dim E_3 = 1 \neq 2$.

\textbf{3.} $A$ n'est \textbf{pas diagonalisable} (multiplicité géométrique $1 <$ multiplicité algébrique $2$). La forme de Jordan est $J = \begin{pmatrix}3&1\\0&3\end{pmatrix}$.
\end{correctionbox}

\begin{exercicebox}[Noyau, image et rang]
Soit $f : \R^3 \to \R^3$ définie par $f(x,y,z) = (x+y, y+z, x-z)$.
\begin{enumerate}
  \item Montrer que $f$ est linéaire.
  \item Déterminer une base de $\Ker(f)$ et sa dimension.
  \item En déduire la dimension de $\Im(f)$. Préciser si $f$ est injective, surjective.
\end{enumerate}
\end{exercicebox}

\begin{correctionbox}[Correction]
\textbf{1.} On vérifie $f(\lambda u + v) = \lambda f(u) + f(v)$ (calcul direct, linéarité de chaque composante).

\textbf{2.} La matrice de $f$ dans la base canonique est $A = \begin{pmatrix}1&1&0\\0&1&1\\1&0&-1\end{pmatrix}$.
On résout $Ax = 0$ par pivot :
$L_3 \leftarrow L_3 - L_1$ donne $\begin{pmatrix}1&1&0\\0&1&1\\0&-1&-1\end{pmatrix}$,
puis $L_3 + L_2$ donne $\begin{pmatrix}1&1&0\\0&1&1\\0&0&0\end{pmatrix}$.
On obtient $y = -z$, $x = z$, donc $\Ker(f) = \mathrm{Vect}(1,-1,1)$, $\dim \Ker(f) = 1$.

\textbf{3.} Par le théorème du rang : $\dim \Im(f) = 3 - 1 = 2$. $f$ n'est ni injective ($\Ker \neq \{0\}$) ni surjective ($\Im \neq \R^3$).
\end{correctionbox}

\begin{retenirbox}
\textbf{Algèbre linéaire — À retenir}
\begin{itemize}
  \item $F$ s.e.v. $\iff$ stable par combinaisons linéaires.
  \item Base : famille libre et génératrice. $\dim E$ = cardinal de toute base.
  \item Théorème du rang : $\dim \Ker f + \dim \Im f = \dim E$.
  \item $f$ injective $\iff \Ker f = \{0\}$ ; $f$ surjective $\iff \Im f = F$.
  \item Valeurs propres = racines de $\chi_A = \det(A-\lambda I)$.
  \item Diagonalisable $\iff$ multiplicités géométriques = multiplicités algébriques.
  \item $A^n = PD^nP^{-1}$ si $A$ diagonalisable.
\end{itemize}
\end{retenirbox}

%% === CHAPITRE A2.2 ===
\chapter{Fonctions de deux variables réelles}

\section*{Objectifs}
\begin{itemize}[leftmargin=1.5em]
  \item Étudier les notions de limite, continuité et différentiabilité pour des fonctions $f : \R^2 \to \R$.
  \item Calculer les dérivées partielles, le gradient et la matrice hessienne.
  \item Déterminer le plan tangent ; caractériser les points critiques.
  \item Appliquer le développement de Taylor-Young à l'ordre 2.
\end{itemize}

\section{Limites et continuité}

La notion de limite en dimension 2 est plus délicate qu'en dimension 1 : il ne suffit pas que la limite existe le long de chaque droite passant par le point, il faut qu'elle soit la même quelle que soit la façon dont on s'approche. Cette richesse topologique de $\mathbb{R}^2$ — où les chemins d'approche sont infiniment plus variés qu'en $\mathbb{R}$ — est à l'origine de nombreux exemples de fonctions admettant des limites directionnelles toutes égales sans avoir de limite globale.

\begin{definitionbox}[Limite en un point]
Soit $f : U \subseteq \R^2 \to \R$ et $(a,b) \in \overline{U}$.
On dit que $\lim_{(x,y)\to(a,b)} f(x,y) = L$ si :
\[
  \forall \varepsilon > 0,\; \exists \delta > 0 : \|(x,y)-(a,b)\| < \delta \implies |f(x,y) - L| < \varepsilon.
\]
$f$ est \emph{continue} en $(a,b)$ si $f(a,b) = \lim_{(x,y)\to(a,b)} f(x,y)$.
\end{definitionbox}

\begin{erreurbox}[Piège : limites suivant des chemins]
L'existence de la limite de $f$ le long de tous les chemins passant par $(a,b)$ ne suffit
\textbf{pas} à garantir l'existence de la limite globale. Il faut utiliser $\varepsilon$-$\delta$ ou des majorations.
\end{erreurbox}

\begin{exemplebox}[Non-existence d'une limite]
Soit $f(x,y) = \dfrac{xy}{x^2+y^2}$ (définie hors $(0,0)$).
\begin{itemize}
  \item Chemin $y = 0$ : $f(x,0) = 0 \to 0$.
  \item Chemin $y = x$ : $f(x,x) = \frac{x^2}{2x^2} = \frac{1}{2} \to \frac{1}{2}$.
\end{itemize}
Les limites diffèrent selon le chemin : $f$ n'a \textbf{pas} de limite en $(0,0)$.
\end{exemplebox}

\section{Dérivées partielles et différentielle}

Les dérivées partielles mesurent le taux de variation de $f$ dans chacune des directions des axes coordonnés, en « gelant » les autres variables. Mais l'existence de toutes les dérivées partielles ne garantit pas la différentiabilité : c'est la différentielle $df(a,b)$ — forme linéaire approchant $f(a+h,b+k) - f(a,b)$ avec un reste $o(\|(h,k)\|)$ — qui constitue la vraie généralisation du concept de dérivée à plusieurs variables.

\begin{definitionbox}[Dérivées partielles]
Soit $f : U \subseteq \R^2 \to \R$ et $(a,b) \in U$.
\begin{align*}
  \frac{\partial f}{\partial x}(a,b) &= \lim_{h \to 0} \frac{f(a+h,b) - f(a,b)}{h}, \\
  \frac{\partial f}{\partial y}(a,b) &= \lim_{h \to 0} \frac{f(a,b+h) - f(a,b)}{h}.
\end{align*}
\end{definitionbox}

\begin{definitionbox}[Différentiabilité et gradient]
$f$ est \emph{différentiable} en $(a,b)$ s'il existe $(l_1,l_2) \in \R^2$ tel que :
\[
  f(a+h,b+k) = f(a,b) + l_1 h + l_2 k + o\!\left(\|(h,k)\|\right) \quad \text{quand } (h,k) \to (0,0).
\]
On a alors $l_1 = \frac{\partial f}{\partial x}(a,b)$ et $l_2 = \frac{\partial f}{\partial y}(a,b)$.
Le \emph{gradient} de $f$ en $(a,b)$ est :
\[
  \nabla f(a,b) = \left(\frac{\partial f}{\partial x}(a,b),\;\frac{\partial f}{\partial y}(a,b)\right).
\]
\end{definitionbox}

\begin{theoremebox}[Condition suffisante de différentiabilité]
Si $\dfrac{\partial f}{\partial x}$ et $\dfrac{\partial f}{\partial y}$ existent et sont
\emph{continues} au voisinage de $(a,b)$, alors $f$ est différentiable en $(a,b)$
(on dit que $f$ est de classe $\mathcal{C}^1$).
\end{theoremebox}

\begin{definitionbox}[Plan tangent]
Si $f$ est différentiable en $(a,b)$, le \emph{plan tangent} au graphe de $f$
au point $(a,b,f(a,b))$ a pour équation :
\[
  z = f(a,b) + \frac{\partial f}{\partial x}(a,b)\,(x-a) + \frac{\partial f}{\partial y}(a,b)\,(y-b).
\]
\end{definitionbox}

\section{Dérivées partielles d'ordre 2 et hessien}

La matrice hessienne encode la courbure d'une surface dans toutes les directions : ses valeurs propres positives (resp. négatives) indiquent que la surface est concave vers le haut (resp. vers le bas) dans les directions correspondantes. Le signe de son déterminant — négatif pour un point selle, positif pour un extremum — est l'outil de classification des points critiques en analyse multivariable, analogue au signe de $f''(a)$ en dimension 1.

\begin{definitionbox}[Matrice hessienne]
Si $f$ est de classe $\mathcal{C}^2$, la \emph{matrice hessienne} de $f$ en $(a,b)$ est :
\[
  H_f(a,b) = \begin{pmatrix}
    \dfrac{\partial^2 f}{\partial x^2} & \dfrac{\partial^2 f}{\partial x \partial y} \\[8pt]
    \dfrac{\partial^2 f}{\partial y \partial x} & \dfrac{\partial^2 f}{\partial y^2}
  \end{pmatrix}_{(a,b)}.
\]
D'après le théorème de Schwarz ($f \in \mathcal{C}^2$) :
$\dfrac{\partial^2 f}{\partial x \partial y} = \dfrac{\partial^2 f}{\partial y \partial x}$,
donc $H_f$ est symétrique.
\end{definitionbox}

\begin{theoremebox}[Développement de Taylor-Young à l'ordre 2]
Si $f$ est de classe $\mathcal{C}^2$ au voisinage de $(a,b)$, alors pour $(h,k) \to (0,0)$ :
\[
  f(a+h,b+k) = f(a,b) + \nabla f(a,b) \cdot \begin{pmatrix}h\\k\end{pmatrix}
  + \frac{1}{2}\begin{pmatrix}h&k\end{pmatrix} H_f(a,b) \begin{pmatrix}h\\k\end{pmatrix}
  + o\!\left(h^2+k^2\right).
\]
\end{theoremebox}

\section{Points critiques}

Les extrema locaux d'une fonction différentiable se trouvent nécessairement parmi ses points critiques (là où le gradient est nul), mais tous les points critiques ne sont pas des extrema : les points selles, où la fonction croît dans certaines directions et décroît dans d'autres, échappent à cette classification. La stratégie d'étude — calculer $\nabla f$, résoudre $\nabla f = 0$, calculer $H_f$ et examiner le signe de $\det H_f$ — est le protocole systématique à appliquer en concours.

\begin{definitionbox}[Point critique]
$(a,b)$ est un \emph{point critique} de $f$ si $\nabla f(a,b) = (0,0)$,
c'est-à-dire $\dfrac{\partial f}{\partial x}(a,b) = \dfrac{\partial f}{\partial y}(a,b) = 0$.
\end{definitionbox}

\begin{theoremebox}[Nature d'un point critique via le hessien]
Soit $(a,b)$ un point critique de $f \in \mathcal{C}^2$.
Posons $\Delta = \det H_f(a,b)$ et $r = \dfrac{\partial^2 f}{\partial x^2}(a,b)$.
\begin{itemize}
  \item $\Delta > 0$ et $r > 0$ : \emph{minimum local}.
  \item $\Delta > 0$ et $r < 0$ : \emph{maximum local}.
  \item $\Delta < 0$ : \emph{point selle} (ni maximum ni minimum).
  \item $\Delta = 0$ : \emph{cas douteux} (le théorème ne conclut pas).
\end{itemize}
\end{theoremebox}

\begin{exemplebox}[Étude d'une fonction de deux variables]
Soit $f(x,y) = x^3 - 3x + y^2 - 2y$.

\textbf{Points critiques.}
$\dfrac{\partial f}{\partial x} = 3x^2 - 3 = 0 \implies x = \pm 1$,\quad
$\dfrac{\partial f}{\partial y} = 2y - 2 = 0 \implies y = 1$.
Points critiques : $(1,1)$ et $(-1,1)$.

\textbf{Hessien.} $\dfrac{\partial^2 f}{\partial x^2} = 6x$, $\dfrac{\partial^2 f}{\partial y^2} = 2$, $\dfrac{\partial^2 f}{\partial x\partial y} = 0$.
$H_f = \begin{pmatrix}6x&0\\0&2\end{pmatrix}$.

En $(1,1)$ : $\Delta = 12 > 0$, $r = 6 > 0$ $\to$ minimum local.
En $(-1,1)$ : $\Delta = -12 < 0$ $\to$ point selle.
\end{exemplebox}

\begin{retenirbox}
\textbf{Fonctions de deux variables — À retenir}
\begin{itemize}
  \item La limite le long d'un chemin ne suffit pas : des chemins différents peuvent donner des résultats différents.
  \item $f \in \mathcal{C}^1 \implies f$ différentiable.
  \item Gradient $\nabla f$ = vecteur des dérivées partielles.
  \item Plan tangent : $z = f(a,b) + \nabla f(a,b) \cdot (x-a, y-b)$.
  \item Point critique $\iff \nabla f = 0$ ; nature via $\Delta = \det H_f$.
\end{itemize}
\end{retenirbox}

\begin{center}
\begin{tikzpicture}
\begin{axis}[
  width=7cm, height=6cm,
  xlabel={$x$}, ylabel={$y$},
  title={\small Courbes de niveau de $f(x,y)=x^2+2y^2$},
  view={0}{90},
  colormap/cool,
  xtick={-2,-1,0,1,2}, ytick={-2,-1,0,1,2},
  tick label style={font=\scriptsize},
  label style={font=\small},
]
\addplot3[contour gnuplot={levels={0.5,1,2,3,4,6}},
  thick, samples=60, domain=-2.2:2.2] {x^2 + 2*y^2};
\addplot[mark=*,mark size=2pt, caplpred] coordinates {(0,0)};
\node[font=\scriptsize, caplpred] at (axis cs:0.15,0.15) {min};
\end{axis}
\end{tikzpicture}
\quad
\begin{tikzpicture}
\begin{axis}[
  width=7cm, height=6cm,
  xlabel={$x$}, ylabel={$y$},
  zlabel={$z$},
  title={\small Surface $z = x^2 + 2y^2$},
  colormap/cool,
  tick label style={font=\scriptsize},
  label style={font=\small},
  view={45}{35},
  samples=20, domain=-2:2,
]
\addplot3[surf, opacity=0.85] {x^2 + 2*y^2};
\end{axis}
\end{tikzpicture}
\captionof*{figure}{\small\itshape Courbes de niveau (gauche) et surface 3D (droite) de $f(x,y)=x^2+2y^2$ — minimum en $(0,0)$.}
\end{center}

%% === CHAPITRE A2.3 ===
\chapter{Développements limités}

\section*{Objectifs}
\begin{itemize}[leftmargin=1.5em]
  \item Définir et calculer un développement limité (DL) en 0 et en un point $a$.
  \item Utiliser les DL usuels (liste à connaître).
  \item Appliquer les DL : calcul de limites, étude locale de fonctions, tangente, branche parabolique.
\end{itemize}

\section{Définition et existence}

Un développement limité est une approximation polynomiale d'une fonction au voisinage d'un point : la partie régulière $\sum_{k=0}^n a_k x^k$ approche $f(x)$ avec un reste $o(x^n)$ qui devient négligeable devant $x^n$ quand $x \to 0$. Cette approximation est d'autant plus précise que l'ordre $n$ est élevé, mais elle n'est valable que dans un petit voisinage du point de développement — contrairement à la convergence d'une série entière qui peut s'étendre sur un intervalle fini.

\begin{definitionbox}[Développement limité à l'ordre $n$ en $0$]
On dit que $f$ admet un \emph{DL à l'ordre $n$ en $0$} s'il existe des réels
$a_0, a_1, \ldots, a_n$ tels que :
\[
  f(x) = a_0 + a_1 x + a_2 x^2 + \cdots + a_n x^n + o(x^n) \quad (x \to 0).
\]
La somme $a_0 + a_1 x + \cdots + a_n x^n$ est la \emph{partie régulière} du DL.
\end{definitionbox}

\begin{theoremebox}[Unicité du DL]
Si $f$ admet un DL à l'ordre $n$ en $0$, il est \emph{unique}.
\end{theoremebox}

\begin{theoremebox}[DL des fonctions $\mathcal{C}^n$ — Formule de Taylor-Young]
Si $f$ est de classe $\mathcal{C}^n$ en $0$, alors $f$ admet un DL à l'ordre $n$ en $0$ et :
\[
  a_k = \frac{f^{(k)}(0)}{k!} \quad (0 \leq k \leq n).
\]
\end{theoremebox}

\section{DL usuels en $0$}

Les développements limités des huit fonctions usuelles — exponentielle, sinus, cosinus, logarithme, puissance générale, $1/(1-x)$, tangente, arctangente — constituent la boîte à outils fondamentale. Ces formules s'obtiennent par la formule de Taylor-Young à partir des dérivées successives, et toutes les autres peuvent s'en déduire par composition, produit ou intégration. Leur maîtrise est un prérequis absolu pour aborder efficacement les calculs de limites et les études locales de fonctions.

\begin{retenirbox}
Les DL suivants sont à connaître par cœur (pour $x \to 0$) :
\begin{align*}
  e^x &= 1 + x + \frac{x^2}{2!} + \cdots + \frac{x^n}{n!} + o(x^n) \\
  \cos x &= 1 - \frac{x^2}{2!} + \frac{x^4}{4!} - \cdots + \frac{(-1)^n x^{2n}}{(2n)!} + o(x^{2n+1}) \\
  \sin x &= x - \frac{x^3}{3!} + \frac{x^5}{5!} - \cdots + \frac{(-1)^n x^{2n+1}}{(2n+1)!} + o(x^{2n+2}) \\
  \ln(1+x) &= x - \frac{x^2}{2} + \frac{x^3}{3} - \cdots + \frac{(-1)^{n-1}x^n}{n} + o(x^n) \\
  (1+x)^\alpha &= 1 + \alpha x + \frac{\alpha(\alpha-1)}{2!}x^2 + \cdots + \binom{\alpha}{n}x^n + o(x^n) \\
  \frac{1}{1-x} &= 1 + x + x^2 + \cdots + x^n + o(x^n) \quad (|x|<1) \\
  \tan x &= x + \frac{x^3}{3} + \frac{2x^5}{15} + o(x^6) \\
  \arctan x &= x - \frac{x^3}{3} + \frac{x^5}{5} - \cdots + \frac{(-1)^n x^{2n+1}}{2n+1} + o(x^{2n+2})
\end{align*}
\end{retenirbox}

\section{Opérations sur les DL}

La puissance des développements limités tient à la facilité avec laquelle on peut les combiner : somme, produit (par troncature à l'ordre voulu), composition (en substituant un DL dans un autre) et intégration terme à terme. Ces opérations permettent de calculer des DL de fonctions compliquées sans jamais calculer de dérivées successives, pourvu que l'on parte des DL usuels et que l'on gère soigneusement les ordres de troncature.

\begin{theoremebox}[Règles opératoires]
\begin{itemize}
  \item \textbf{Somme/différence :} $DL_n(f) \pm DL_n(g) = DL_n(f\pm g)$.
  \item \textbf{Produit :} multiplier les parties régulières et tronquer à l'ordre $n$.
  \item \textbf{Composition :} substituer $u = g(x)$ dans $DL_n(f(u))$ si $g(0)=0$.
  \item \textbf{Intégration :} intégrer terme à terme la partie régulière.
  \item \textbf{Dérivation :} dériver terme à terme (attention : l'ordre peut diminuer).
\end{itemize}
\end{theoremebox}

\begin{erreurbox}[Pièges classiques]
\begin{itemize}
  \item \textbf{Erreur 1.} Tronquer trop tôt lors d'un produit.
    Exemple : $e^x \cdot \sin x$ à l'ordre 3 $\neq$ $(1+x+\frac{x^2}{2})(x)$.
    Il faut développer chaque facteur à l'ordre 3 \emph{avant} de multiplier.
  \item \textbf{Erreur 2.} Composer sans vérifier que $g(0) = 0$.
\end{itemize}
\end{erreurbox}

\section{Applications des DL}

Les développements limités ont trois applications directes en analyse : le calcul de limites indéterminées (formes $0/0$ ou $\infty/\infty$) par identification du terme dominant, l'étude locale d'une courbe au voisinage d'un point (nature d'un extremum, position par rapport à la tangente, point d'inflexion), et la détermination d'équivalents de fonctions. Pour les formes $0/0$, la méthode par DL est souvent plus systématique et moins sujette à erreur que la règle de l'Hôpital.

\begin{methodebox}[Calcul de limite par DL]
Pour calculer $\lim_{x\to 0} \dfrac{f(x)}{g(x)}$ :
\begin{enumerate}
  \item Écrire $f(x) = a_k x^k + o(x^k)$ et $g(x) = b_j x^j + o(x^j)$ avec $a_k \neq 0$, $b_j \neq 0$.
  \item Si $k = j$ : limite $= \dfrac{a_k}{b_j}$.
  \item Si $k > j$ : limite $= 0$.
  \item Si $k < j$ : limite $= \pm\infty$ (selon les signes).
\end{enumerate}
\end{methodebox}

\begin{exemplebox}[Limite par DL]
$\displaystyle\lim_{x\to 0} \frac{\tan x - \sin x}{x^3}$.

$\tan x = x + \dfrac{x^3}{3} + o(x^3)$, $\sin x = x - \dfrac{x^3}{6} + o(x^3)$.

$\tan x - \sin x = \dfrac{x^3}{3} + \dfrac{x^3}{6} + o(x^3) = \dfrac{x^3}{2} + o(x^3)$.

Donc la limite vaut $\dfrac{1/2}{1} = \dfrac{1}{2}$.
\end{exemplebox}

\begin{methodebox}[Position d'une courbe par rapport à sa tangente]
En un point $x_0$, si $f(x) = f(x_0) + f'(x_0)(x-x_0) + a_k(x-x_0)^k + \cdots$
avec $a_k$ premier terme non nul d'ordre $\geq 2$ :
\begin{itemize}
  \item $k$ impair : la courbe traverse la tangente en $x_0$ (point d'inflexion).
  \item $k$ pair et $a_k > 0$ : courbe au-dessus de la tangente (minimum local).
  \item $k$ pair et $a_k < 0$ : courbe en dessous (maximum local).
\end{itemize}
\end{methodebox}

\begin{center}
\begin{tikzpicture}
\begin{axis}[
  width=10cm, height=6cm,
  xlabel={$x$}, ylabel={$y$},
  title={\small $\sin x$ et ses approximations par DL en $0$},
  xmin=-3.5, xmax=3.5, ymin=-1.8, ymax=1.8,
  domain=-3.3:3.3, samples=120,
  tick label style={font=\scriptsize},
  label style={font=\small},
  legend style={font=\scriptsize, at={(1.02,1)}, anchor=north west},
  grid=major, grid style={dotted,gray!40},
]
\addplot[thick, black] {sin(deg(x))};
\addlegendentry{$\sin x$}
\addplot[thick, caplpblue, dashed] {x};
\addlegendentry{DL ord.\ 1 : $x$}
\addplot[thick, caplpgreen, dash dot] {x - x^3/6};
\addlegendentry{DL ord.\ 3 : $x-\frac{x^3}{6}$}
\addplot[thick, caplpred, dotted] {x - x^3/6 + x^5/120};
\addlegendentry{DL ord.\ 5 : $x-\frac{x^3}{6}+\frac{x^5}{120}$}
\end{axis}
\end{tikzpicture}
\captionof*{figure}{\small\itshape Plus l'ordre augmente, plus l'approximation est précise, mais sur un intervalle de plus en plus large autour de $0$.}
\end{center}

\begin{exercicebox}[DL et limite]
Calculer $\displaystyle\lim_{x\to 0}\frac{e^x - 1 - x - \frac{x^2}{2}}{x^3}$.
\end{exercicebox}
\begin{correctionbox}[Correction]
$e^x = 1 + x + \dfrac{x^2}{2} + \dfrac{x^3}{6} + o(x^3)$.

Donc $e^x - 1 - x - \dfrac{x^2}{2} = \dfrac{x^3}{6} + o(x^3)$.

La limite vaut $\dfrac{1/6}{1} = \dfrac{1}{6}$.
\end{correctionbox}

\begin{retenirbox}
\textbf{DL — À retenir}
\begin{itemize}
  \item DL à l'ordre $n$ : $f(x) = \sum_{k=0}^n \frac{f^{(k)}(0)}{k!}x^k + o(x^n)$.
  \item Connaître les 8 DL usuels par cœur.
  \item Produit : multiplier et tronquer. Composition : substituer si $g(0)=0$.
  \item Applications : limites ($0/0$), nature des points critiques, asymptotes.
\end{itemize}
\end{retenirbox}

%% === CHAPITRE A2.4 ===
\chapter{Séries numériques}

\section*{Objectifs}
\begin{itemize}[leftmargin=1.5em]
  \item Définir la convergence d'une série numérique.
  \item Connaître et appliquer les critères de convergence (termes positifs, comparaison, Riemann, règle de d'Alembert).
  \item Reconnaître les séries de référence.
\end{itemize}

\section{Définitions fondamentales}

Une série numérique est une somme infinie dont la convergence n'est pas automatique : elle est définie comme la limite de la suite de ses sommes partielles, et cette limite peut ne pas exister. La difficulté principale est que la condition nécessaire $u_n \to 0$ est loin d'être suffisante — la série harmonique $\sum 1/n$ en est le contre-exemple classique. L'étude de la convergence des séries requiert donc des critères spécifiques qui vont au-delà du comportement du terme général.

\begin{definitionbox}[Série numérique]
Soit $(u_n)_{n \geq 0}$ une suite réelle.
On appelle \emph{série de terme général $u_n$} et on note $\sum u_n$ la suite des
\emph{sommes partielles} $S_N = \sum_{n=0}^N u_n$.

La série \emph{converge} si la suite $(S_N)$ converge, et l'on note $\sum_{n=0}^{+\infty} u_n$
sa limite (la \emph{somme} de la série).

Si $(S_N)$ diverge, la série \emph{diverge}.
\end{definitionbox}

\begin{theoremebox}[Condition nécessaire de convergence]
Si $\sum u_n$ converge, alors $\lim_{n\to+\infty} u_n = 0$.
\end{theoremebox}

\begin{preuvebox}
$u_n = S_n - S_{n-1} \to S - S = 0$ si $S_n \to S$. \qed
\end{preuvebox}

\begin{erreurbox}
La réciproque est \textbf{fausse} : $u_n \to 0$ ne suffit pas pour conclure à la convergence.
Contre-exemple : la série harmonique $\sum \dfrac{1}{n}$ diverge bien que $\dfrac{1}{n} \to 0$.
\end{erreurbox}

\section{Séries de référence}

La série géométrique et la série de Riemann sont les deux pierres de touche de l'étude de convergence des séries : elles servent de termes de comparaison pour une immense variété de séries. La série géométrique, dont la somme est explicite pour $|q| < 1$, est aussi la série la plus utilisée en pratique pour les calculs ; la série de Riemann, par son comportement au seuil $\alpha = 1$, illustre la subtilité de la frontière entre convergence et divergence.

\begin{theoremebox}[Série géométrique]
Pour $q \in \R$ : $\displaystyle\sum_{n=0}^{N} q^n = \frac{1-q^{N+1}}{1-q}$ si $q \neq 1$.
La série $\sum_{n=0}^{+\infty} q^n$ converge \emph{si et seulement si} $|q| < 1$,
et sa somme est $\dfrac{1}{1-q}$.
\end{theoremebox}

\begin{theoremebox}[Série de Riemann]
La série $\sum_{n=1}^{+\infty} \dfrac{1}{n^\alpha}$ converge si et seulement si $\alpha > 1$.
\end{theoremebox}

\section{Critères de convergence (séries à termes positifs)}

Pour les séries à termes positifs, la stratégie est de comparer le terme général à celui d'une série de référence (Riemann, géométrique) dont on connaît la nature. Le critère des équivalents est le plus commode : il suffit de déterminer l'équivalent dominant de $u_n$ quand $n \to +\infty$. La règle de d'Alembert est préférable lorsque le terme général fait intervenir des factorielles ou des exponentielles.

\begin{theoremebox}[Critère de comparaison]
Soient $u_n \geq 0$ et $v_n \geq 0$ avec $u_n \leq v_n$ pour tout $n$ assez grand.
\begin{itemize}
  \item Si $\sum v_n$ converge, alors $\sum u_n$ converge.
  \item Si $\sum u_n$ diverge, alors $\sum v_n$ diverge.
\end{itemize}
\end{theoremebox}

\begin{theoremebox}[Critère des équivalents]
Si $u_n \sim v_n$ (i.e.\ $\dfrac{u_n}{v_n} \to 1$) avec $u_n, v_n > 0$,
alors $\sum u_n$ et $\sum v_n$ ont même nature.
\end{theoremebox}

\begin{theoremebox}[Règle de d'Alembert]
Soit $u_n > 0$ et $\ell = \lim_{n\to+\infty} \dfrac{u_{n+1}}{u_n}$.
\begin{itemize}
  \item Si $\ell < 1$ : $\sum u_n$ converge.
  \item Si $\ell > 1$ : $\sum u_n$ diverge.
  \item Si $\ell = 1$ : la règle est \textbf{inconclusive}.
\end{itemize}
\end{theoremebox}

\begin{theoremebox}[Règle de Cauchy (racine)]
Soit $u_n \geq 0$ et $\ell = \limsup_{n\to+\infty} u_n^{1/n}$.
\begin{itemize}
  \item Si $\ell < 1$ : $\sum u_n$ converge absolument.
  \item Si $\ell > 1$ : $\sum u_n$ diverge.
  \item Si $\ell = 1$ : inconclusive.
\end{itemize}
\end{theoremebox}

\section{Convergence absolue}

La convergence absolue est une forme de convergence plus forte que la simple convergence : si $\sum |u_n|$ converge, alors $\sum u_n$ converge aussi, et ce quels que soient les signes des termes. La convergence conditionnelle — convergence de $\sum u_n$ sans convergence de $\sum |u_n|$ — est plus fragile : la somme d'une telle série dépend de l'ordre dans lequel on additionne les termes (théorème de réarrangement de Riemann).

\begin{definitionbox}[Convergence absolue]
$\sum u_n$ \emph{converge absolument} si $\sum |u_n|$ converge.
\end{definitionbox}

\begin{theoremebox}
Convergence absolue $\implies$ convergence. La réciproque est fausse en général.
\end{theoremebox}

\begin{exemplebox}[Exemples fondamentaux]
\begin{itemize}
  \item $\sum \dfrac{1}{n^2}$ : série de Riemann avec $\alpha=2>1$ $\to$ \textbf{converge}. Sa somme est $\dfrac{\pi^2}{6}$ (hors programme).
  \item $\sum \dfrac{(-1)^n}{n}$ : converge (critère d'Abel), mais pas absolument (série harmonique diverge).
  \item $\sum \dfrac{n!}{3^n}$ : $\dfrac{u_{n+1}}{u_n} = \dfrac{n+1}{3} \to +\infty > 1$ $\to$ \textbf{diverge}.
  \item $\sum \dfrac{1}{n^{3/2}}$ : $\alpha = \dfrac{3}{2} > 1$ $\to$ \textbf{converge}.
\end{itemize}
\end{exemplebox}

\begin{exercicebox}[Nature d'une série]
Étudier la nature des séries : (1) $\displaystyle\sum_{n\geq 1} \frac{n+1}{n^3+2}$ ;
\quad (2) $\displaystyle\sum_{n\geq 0} \frac{2^n}{n!}$.
\end{exercicebox}

\begin{correctionbox}[Correction]
\textbf{(1).} Pour $n$ grand, $\dfrac{n+1}{n^3+2} \sim \dfrac{n}{n^3} = \dfrac{1}{n^2}$.
Comme $\sum \dfrac{1}{n^2}$ converge (Riemann, $\alpha=2>1$), par critère des équivalents, $\sum \dfrac{n+1}{n^3+2}$ \textbf{converge}.

\textbf{(2).} $u_n = \dfrac{2^n}{n!} > 0$. Règle de d'Alembert :
$\dfrac{u_{n+1}}{u_n} = \dfrac{2}{n+1} \to 0 < 1$.
Donc $\sum \dfrac{2^n}{n!}$ \textbf{converge} (et sa somme est $e^2 - 1$ d'après le DL de $e^x$).
\end{correctionbox}

\begin{center}
\begin{tikzpicture}
\begin{axis}[
  width=10cm, height=5.5cm,
  xlabel={$N$ (rang de la somme partielle)},
  ylabel={$S_N$},
  title={\small Sommes partielles de $\displaystyle\sum_{n=0}^N q^n$ pour différentes valeurs de $q$},
  xmin=0, xmax=14, ymin=0, ymax=5.5,
  xtick={0,2,...,14}, ytick={0,1,2,3,4,5},
  tick label style={font=\scriptsize},
  label style={font=\small},
  legend style={font=\scriptsize, at={(1.02,1)}, anchor=north west},
  grid=major, grid style={dotted, gray!40},
]
% q=0.5, somme = 1/(1-0.5)=2
\addplot+[mark=*, mark size=1.5pt, thick, caplpblue]
  coordinates {(0,1)(1,1.5)(2,1.75)(3,1.875)(4,1.9375)(5,1.969)(6,1.984)(7,1.992)(8,1.996)(9,1.998)(10,1.999)(11,2.0)(12,2.0)(13,2.0)(14,2.0)};
\addlegendentry{$q=0{,}5$ (CV vers 2)}
% q=0.8, somme = 1/(1-0.8)=5
\addplot+[mark=square*, mark size=1.5pt, thick, caplpgreen]
  coordinates {(0,1)(1,1.8)(2,2.44)(3,2.952)(4,3.362)(5,3.689)(6,3.951)(7,4.161)(8,4.329)(9,4.463)(10,4.571)(11,4.657)(12,4.725)(13,4.780)(14,4.824)};
\addlegendentry{$q=0{,}8$ (CV vers 5)}
% Ligne asymptote q=0.5
\addplot[dashed, caplpblue!50, domain=0:14] {2};
\addplot[dashed, caplpgreen!50, domain=0:14] {5};
\end{axis}
\end{tikzpicture}
\captionof*{figure}{\small\itshape Les sommes partielles $S_N$ convergent vers $\frac{1}{1-q}$ quand $|q|<1$. Plus $q$ est proche de 1, plus la convergence est lente.}
\end{center}

\begin{retenirbox}
\textbf{Séries numériques — À retenir}
\begin{itemize}
  \item Condition nécessaire : $u_n \to 0$. Insuffisant ! (série harmonique).
  \item Série géométrique : converge $\iff |q|<1$.
  \item Série de Riemann $\sum 1/n^\alpha$ : converge $\iff \alpha>1$.
  \item D'Alembert : $\ell < 1 \to$ CV ; $\ell > 1 \to$ DV ; $\ell=1 \to$ nul.
  \item CV absolue $\implies$ CV (réciproque fausse).
\end{itemize}
\end{retenirbox}

%% === CHAPITRE A2.5 ===
\chapter{Dénombrement}

\section*{Objectifs}
\begin{itemize}[leftmargin=1.5em]
  \item Maîtriser les règles du dénombrement : cardinal d'une réunion, produit cartésien.
  \item Compter les arrangements, permutations et combinaisons.
  \item Connaître et appliquer la formule du binôme de Newton.
\end{itemize}

\section{Cardinal d'un ensemble fini}

La formule d'inclusion-exclusion est le résultat fondateur du dénombrement d'unions : elle corrige le double comptage des éléments appartenant à plusieurs ensembles à la fois. Pour deux ensembles, elle se réduit à $|A \cup B| = |A| + |B| - |A \cap B|$ ; pour $n$ ensembles, la formule de Poincaré alterne les contributions des intersections de toutes les tailles, en additionnant les singletons, soustrayant les paires, etc.

\begin{definitionbox}[Cardinal]
Soit $E$ un ensemble fini. Son \emph{cardinal} (ou nombre d'éléments) est noté $\Card(E)$ ou $|E|$.
\end{definitionbox}

\begin{theoremebox}[Formule d'inclusion-exclusion]
Pour deux ensembles finis $A$ et $B$ :
\[
  |A \cup B| = |A| + |B| - |A \cap B|.
\]
Plus généralement, pour $n$ ensembles $A_1, \ldots, A_n$ (formule de Poincaré) :
\[
  \left|\bigcup_{i=1}^n A_i\right| = \sum_i |A_i| - \sum_{i<j}|A_i \cap A_j| + \cdots + (-1)^{n-1}|A_1 \cap \cdots \cap A_n|.
\]
\end{theoremebox}

\section{Arrangements et permutations}

Les arrangements comptent des sélections ordonnées : choisir $k$ éléments parmi $n$ en tenant compte de l'ordre dans lequel ils sont placés. Le passage de $A_n^k = n!/(n-k)!$ à $\binom{n}{k} = A_n^k/k!$ reflète exactement la suppression de l'ordre : diviser par $k!$ revient à regrouper en une seule combinaison toutes les permutations d'un même sous-ensemble de $k$ éléments.

\begin{definitionbox}[Arrangements sans répétition]
Un \emph{arrangement de $k$ éléments parmi $n$} est une liste ordonnée de $k$ éléments
distincts d'un ensemble à $n$ éléments ($0 \leq k \leq n$). Le nombre de tels arrangements est :
\[
  A_n^k = \frac{n!}{(n-k)!} = n(n-1)\cdots(n-k+1).
\]
\end{definitionbox}

\begin{definitionbox}[Permutations]
Une \emph{permutation} de $n$ éléments est un arrangement de $n$ éléments parmi $n$.
Le nombre de permutations est $n! = n \times (n-1) \times \cdots \times 1$.
Conventions : $0! = 1$.
\end{definitionbox}

\section{Combinaisons}

Les combinaisons comptent des sélections non ordonnées : $\binom{n}{k}$ est le nombre de sous-ensembles de taille $k$ d'un ensemble à $n$ éléments. Le triangle de Pascal, construit par la relation $\binom{n+1}{k} = \binom{n}{k-1} + \binom{n}{k}$, permet de calculer rapidement tous les coefficients binomiaux et de visualiser leur symétrie et leurs propriétés de sommation.

\begin{definitionbox}[Combinaisons]
Une \emph{combinaison de $k$ éléments parmi $n$} est un sous-ensemble (sans ordre) de $k$
éléments d'un ensemble à $n$ éléments. Le nombre de combinaisons est :
\[
  \binom{n}{k} = \frac{n!}{k!(n-k)!} \quad (0 \leq k \leq n).
\]
\end{definitionbox}

\begin{theoremebox}[Propriétés des coefficients binomiaux]
\begin{itemize}
  \item $\dbinom{n}{0} = \dbinom{n}{n} = 1$ ; $\dbinom{n}{1} = \dbinom{n}{n-1} = n$.
  \item Symétrie : $\dbinom{n}{k} = \dbinom{n}{n-k}$.
  \item Formule de Pascal : $\dbinom{n+1}{k} = \dbinom{n}{k-1} + \dbinom{n}{k}$.
  \item $\displaystyle\sum_{k=0}^n \binom{n}{k} = 2^n$.
\end{itemize}
\end{theoremebox}

\begin{theoremebox}[Formule du binôme de Newton]
Pour tous $a, b \in \R$ et $n \in \N$ :
\[
  (a+b)^n = \sum_{k=0}^n \binom{n}{k} a^{n-k} b^k.
\]
\end{theoremebox}

\begin{preuvebox}
Par récurrence sur $n$. Initialisation : $(a+b)^0 = 1 = \dbinom{0}{0}$. \\
Hérédité : $(a+b)^{n+1} = (a+b)(a+b)^n = (a+b)\sum_k \dbinom{n}{k} a^{n-k}b^k$.
En développant et en utilisant la relation de Pascal $\dbinom{n}{k-1}+\dbinom{n}{k}=\dbinom{n+1}{k}$,
on obtient $(a+b)^{n+1} = \sum_{k=0}^{n+1}\dbinom{n+1}{k}a^{n+1-k}b^k$. \qed
\end{preuvebox}

\begin{exemplebox}[Application du binôme]
Développer $(2x-1)^4$.
\[
  (2x-1)^4 = \sum_{k=0}^4 \binom{4}{k}(2x)^{4-k}(-1)^k
  = 16x^4 - 32x^3 + 24x^2 - 8x + 1.
\]
\end{exemplebox}

\begin{exercicebox}[Problème de dénombrement]
Un comité de 4 personnes est choisi parmi 6 femmes et 4 hommes.
\begin{enumerate}
  \item Combien de comités possibles ?
  \item Combien contiennent au moins 3 femmes ?
  \item Combien contiennent exactement 1 homme ?
\end{enumerate}
\end{exercicebox}

\begin{correctionbox}[Correction]
\textbf{1.} $\dbinom{10}{4} = 210$ comités possibles.

\textbf{2.} Au moins 3 femmes : exactement 3 femmes ou 4 femmes.
\begin{itemize}
  \item 3 femmes, 1 homme : $\dbinom{6}{3}\dbinom{4}{1} = 20 \times 4 = 80$.
  \item 4 femmes, 0 homme : $\dbinom{6}{4}\dbinom{4}{0} = 15$.
\end{itemize}
Total : $80 + 15 = 95$.

\textbf{3.} Exactement 1 homme $=$ 3 femmes et 1 homme $= 80$ (calculé ci-dessus).
\end{correctionbox}

\begin{retenirbox}
\textbf{Dénombrement — À retenir}
\begin{itemize}
  \item Arrangements ordonnés : $A_n^k = \dfrac{n!}{(n-k)!}$.
  \item Combinaisons (non ordonnées) : $\dbinom{n}{k} = \dfrac{n!}{k!(n-k)!}$.
  \item Formule de Pascal : $\dbinom{n+1}{k} = \dbinom{n}{k-1}+\dbinom{n}{k}$.
  \item Binôme de Newton : $(a+b)^n = \sum_{k=0}^n \dbinom{n}{k}a^{n-k}b^k$.
  \item Inclusion-exclusion : $|A \cup B| = |A|+|B|-|A\cap B|$.
\end{itemize}
\end{retenirbox}

%% === CHAPITRE A2.6 ===
\chapter{Probabilités discrètes avancées}

\section*{Objectifs}
\begin{itemize}[leftmargin=1.5em]
  \item Maîtriser les lois discrètes : uniforme, Bernoulli, binomiale, géométrique, Poisson.
  \item Connaître la loi exponentielle (loi continue) et ses propriétés.
  \item Calculer espérance, variance, écart-type pour ces lois.
  \item Appliquer la convergence de la loi binomiale vers la loi de Poisson.
\end{itemize}

\section{Rappels : variables aléatoires discrètes}

Ce chapitre approfondit l'étude des variables aléatoires discrètes introduites au chapitre 13 en introduisant de nouvelles lois adaptées à des situations de modélisation plus précises. Les propriétés d'espérance, variance et indépendance restent les mêmes outils de calcul, mais les lois géométrique, de Poisson et exponentielle permettent de modéliser des phénomènes que ni la loi uniforme ni la loi binomiale ne capturent correctement.

\begin{definitionbox}[Variable aléatoire discrète]
Une \emph{variable aléatoire discrète} $X$ sur $(\Omega, \mathcal{A}, \mathbb{P})$ est une fonction
$X : \Omega \to \R$ à valeurs dans un ensemble fini ou dénombrable.
Sa loi est définie par les probabilités $p_k = \mathbb{P}(X = x_k)$.
\end{definitionbox}

\begin{definitionbox}[Espérance, variance, écart-type]
\begin{align*}
  \mathbb{E}(X) &= \sum_k x_k \, p_k \quad \text{(si cette somme est absolument convergente)}.\\
  \mathbb{V}(X) &= \mathbb{E}((X - \mathbb{E}(X))^2) = \mathbb{E}(X^2) - (\mathbb{E}(X))^2.\\
  \sigma(X) &= \sqrt{\mathbb{V}(X)}.
\end{align*}
\end{definitionbox}

\section{Loi géométrique}

La loi géométrique modélise le temps d'attente — le rang du premier succès — dans une suite d'épreuves de Bernoulli indépendantes et identiquement distribuées. Sa propriété caractéristique est l'absence de mémoire : savoir que les $m$ premiers essais ont échoué ne modifie pas la distribution du nombre d'essais restants nécessaires. C'est la seule loi discrète à valeurs dans $\mathbb{N}^*$ qui possède cette propriété.

\begin{definitionbox}[Loi géométrique]
Soit $p \in {]0,1]}$. On dit que $X$ suit la \emph{loi géométrique} de paramètre $p$
(notée $\mathcal{G}(p)$) si $X$ prend ses valeurs dans $\{1, 2, 3, \ldots\}$ et :
\[
  \mathbb{P}(X = k) = (1-p)^{k-1}\, p \quad (k \geq 1).
\]
$X$ modélise le rang du premier succès dans une suite d'épreuves de Bernoulli indépendantes
de paramètre $p$.
\end{definitionbox}

\begin{theoremebox}[Espérance et variance de $\mathcal{G}(p)$]
Si $X \sim \mathcal{G}(p)$ :
\[
  \mathbb{E}(X) = \frac{1}{p}, \qquad \mathbb{V}(X) = \frac{1-p}{p^2}.
\]
\end{theoremebox}

\begin{theoremebox}[Propriété d'absence de mémoire]
Soit $X \sim \mathcal{G}(p)$. Pour tous entiers $m, n \geq 1$ :
\[
  \mathbb{P}(X > m+n \mid X > m) = \mathbb{P}(X > n).
\]
C'est la \emph{propriété d'absence de mémoire}, qui caractérise la loi géométrique
parmi les lois à valeurs dans $\mathbb{N}^*$.
\end{theoremebox}

\section{Loi de Poisson}

La loi de Poisson est la loi des événements rares et nombreux : elle apparaît comme limite de la loi binomiale $\mathcal{B}(n, \lambda/n)$ quand $n \to +\infty$, et modélise le nombre d'occurrences d'un phénomène dans un intervalle de temps ou d'espace (arrivées de clients, pannes, erreurs de transmission). Sa propriété remarquable est que l'espérance et la variance sont toutes deux égales au paramètre $\lambda$.

\begin{definitionbox}[Loi de Poisson]
Soit $\lambda > 0$. On dit que $X$ suit la \emph{loi de Poisson} de paramètre $\lambda$
(notée $\mathcal{P}(\lambda)$) si $X$ prend ses valeurs dans $\N$ et :
\[
  \mathbb{P}(X = k) = e^{-\lambda} \frac{\lambda^k}{k!} \quad (k \in \N).
\]
\end{definitionbox}

\begin{preuvebox}
Vérifions que c'est bien une loi : $\displaystyle\sum_{k=0}^{+\infty} e^{-\lambda}\frac{\lambda^k}{k!}
= e^{-\lambda} \sum_{k=0}^{+\infty} \frac{\lambda^k}{k!} = e^{-\lambda} e^{\lambda} = 1$. \qed
\end{preuvebox}

\begin{theoremebox}[Espérance et variance de $\mathcal{P}(\lambda)$]
Si $X \sim \mathcal{P}(\lambda)$ :
\[
  \mathbb{E}(X) = \lambda, \qquad \mathbb{V}(X) = \lambda.
\]
\end{theoremebox}

\begin{preuvebox}
$\mathbb{E}(X) = \sum_{k=1}^{+\infty} k \, e^{-\lambda}\dfrac{\lambda^k}{k!}
= e^{-\lambda}\lambda \sum_{k=1}^{+\infty} \dfrac{\lambda^{k-1}}{(k-1)!}
= e^{-\lambda}\lambda \, e^{\lambda} = \lambda$.

De même, $\mathbb{E}(X(X-1)) = \lambda^2$, d'où $\mathbb{E}(X^2) = \lambda^2 + \lambda$
et $\mathbb{V}(X) = \lambda^2 + \lambda - \lambda^2 = \lambda$. \qed
\end{preuvebox}

\begin{theoremebox}[Convergence de la binomiale vers Poisson]
Soit $X_n \sim \mathcal{B}(n,p_n)$ avec $n \to +\infty$, $p_n \to 0$ et $np_n \to \lambda > 0$.
Alors pour tout $k \in \N$ :
\[
  \mathbb{P}(X_n = k) \underset{n \to +\infty}{\longrightarrow} e^{-\lambda}\frac{\lambda^k}{k!}.
\]
\emph{Interprétation :} la loi de Poisson est la limite de la loi binomiale pour des événements
rares ($p$ petit) et nombreux ($n$ grand).
\end{theoremebox}

\section{Loi exponentielle}

La loi exponentielle est l'analogue continu de la loi géométrique : elle modélise la durée de vie d'un composant sans vieillissement (durée avant la première panne dans un processus de Poisson), ou plus généralement tout temps d'attente continu sans mémoire. La propriété d'absence de mémoire $\mathbb{P}(X > s+t \mid X > s) = \mathbb{P}(X > t)$ signifie que la loi conditionnelle de $X - s$ sachant $X > s$ est la même loi exponentielle : le composant « ne vieillit pas ».

\begin{definitionbox}[Loi exponentielle]
Soit $\lambda > 0$. Une variable aléatoire $X$ suit la \emph{loi exponentielle} de paramètre $\lambda$
(notée $\mathcal{E}(\lambda)$) si elle est à valeurs dans $[0,+\infty[$ avec :
\[
  f_X(t) = \lambda e^{-\lambda t}\, \mathbf{1}_{t \geq 0} \quad \text{(densité de probabilité)}.
\]
La fonction de répartition est $F_X(t) = 1 - e^{-\lambda t}$ pour $t \geq 0$.
\end{definitionbox}

\begin{theoremebox}[Espérance et variance de $\mathcal{E}(\lambda)$]
Si $X \sim \mathcal{E}(\lambda)$ :
\[
  \mathbb{E}(X) = \frac{1}{\lambda}, \qquad \mathbb{V}(X) = \frac{1}{\lambda^2}.
\]
\end{theoremebox}

\begin{theoremebox}[Absence de mémoire continue]
Si $X \sim \mathcal{E}(\lambda)$, alors pour tous $s, t \geq 0$ :
\[
  \mathbb{P}(X > s+t \mid X > s) = \mathbb{P}(X > t) = e^{-\lambda t}.
\]
La loi exponentielle est la seule loi continue vérifiant la propriété d'absence de mémoire.
\end{theoremebox}

\begin{preuvebox}
$\mathbb{P}(X > s+t \mid X > s) = \dfrac{\mathbb{P}(X > s+t)}{\mathbb{P}(X > s)}
= \dfrac{e^{-\lambda(s+t)}}{e^{-\lambda s}} = e^{-\lambda t} = \mathbb{P}(X > t)$. \qed
\end{preuvebox}

\begin{exemplebox}[Modélisation par une loi de Poisson]
Un serveur web reçoit en moyenne 5 requêtes par seconde.
On modélise le nombre $X$ de requêtes par seconde par $X \sim \mathcal{P}(5)$.
\begin{itemize}
  \item $\mathbb{P}(X = 3) = e^{-5}\dfrac{5^3}{3!} = e^{-5}\dfrac{125}{6} \approx 0{,}1404$.
  \item $\mathbb{P}(X \geq 2) = 1 - \mathbb{P}(X=0) - \mathbb{P}(X=1)
    = 1 - e^{-5} - 5e^{-5} = 1 - 6e^{-5} \approx 0{,}9596$.
\end{itemize}
\end{exemplebox}

\begin{center}
\begin{tikzpicture}
\begin{axis}[
  width=10cm, height=5.5cm,
  xlabel={$k$},
  ylabel={$\mathbb{P}(X=k)$},
  title={\small Loi de Poisson $\mathcal{P}(\lambda)$ pour $\lambda \in \{1, 3, 5\}$},
  ybar, bar width=4pt,
  xmin=-0.5, xmax=14.5,
  ymin=0, ymax=0.42,
  xtick={0,1,...,14},
  tick label style={font=\scriptsize},
  label style={font=\small},
  legend style={font=\scriptsize, at={(0.98,0.98)}, anchor=north east},
  grid=major, grid style={dotted, gray!30},
]
% lambda=1 : P(X=k) = e^{-1}/k!
\addplot[fill=caplpbluebg, draw=caplpblue] coordinates {
  (0,0.3679)(1,0.3679)(2,0.1839)(3,0.0613)(4,0.0153)(5,0.0031)(6,0.0005)};
\addlegendentry{$\lambda=1$}
% lambda=3 : P(X=k) = e^{-3} * 3^k / k!
\addplot[fill=caplpgreenbg, draw=caplpgreen] coordinates {
  (0,0.0498)(1,0.1494)(2,0.2240)(3,0.2240)(4,0.1680)(5,0.1008)(6,0.0504)(7,0.0216)(8,0.0081)(9,0.0027)};
\addlegendentry{$\lambda=3$}
% lambda=5 : P(X=k) = e^{-5} * 5^k / k!
\addplot[fill=caplpredbg, draw=caplpred] coordinates {
  (0,0.0067)(1,0.0337)(2,0.0842)(3,0.1404)(4,0.1755)(5,0.1755)(6,0.1462)(7,0.1044)(8,0.0653)(9,0.0363)(10,0.0181)(11,0.0082)(12,0.0034)(13,0.0013)(14,0.0005)};
\addlegendentry{$\lambda=5$}
\end{axis}
\end{tikzpicture}
\captionof*{figure}{\small\itshape Distribution de Poisson $\mathcal{P}(\lambda)$ : l'espérance et la variance sont toutes deux égales à $\lambda$.}
\end{center}

\begin{exercicebox}[Loi géométrique et exponentielle]
Un contrôle qualité détecte une pièce défectueuse avec probabilité $p = 0{,}04$.
On note $X$ le rang de la première pièce défectueuse détectée.
\begin{enumerate}
  \item Quelle est la loi de $X$ ? Calculer $\mathbb{E}(X)$ et $\sigma(X)$.
  \item Calculer $\mathbb{P}(X > 20)$.
  \item La durée de vie (en heures) d'un composant est modélisée par $T \sim \mathcal{E}(0{,}01)$.
    Calculer $\mathbb{E}(T)$, $\mathbb{P}(T > 100)$ et $\mathbb{P}(T > 200 \mid T > 100)$.
\end{enumerate}
\end{exercicebox}

\begin{correctionbox}[Correction]
\textbf{1.} $X \sim \mathcal{G}(0{,}04)$. $\mathbb{E}(X) = \dfrac{1}{0{,}04} = 25$.
$\mathbb{V}(X) = \dfrac{0{,}96}{0{,}04^2} = 600$, $\sigma(X) = \sqrt{600} \approx 24{,}49$.

\textbf{2.} $\mathbb{P}(X > 20) = \sum_{k=21}^{+\infty}(0{,}96)^{k-1}(0{,}04)
= (0{,}96)^{20} \approx 0{,}4420$.

\textbf{3.} $\mathbb{E}(T) = \dfrac{1}{0{,}01} = 100$ h.
$\mathbb{P}(T > 100) = e^{-0{,}01 \times 100} = e^{-1} \approx 0{,}3679$.
Par absence de mémoire : $\mathbb{P}(T > 200 \mid T > 100) = \mathbb{P}(T > 100) = e^{-1} \approx 0{,}3679$.
\end{correctionbox}

\begin{retenirbox}
\textbf{Probabilités discrètes avancées — À retenir}

\begin{tabular}{lllll}
\toprule
Loi & Paramètre & $\mathbb{P}(X=k)$ & $\mathbb{E}(X)$ & $\mathbb{V}(X)$ \\
\midrule
$\mathcal{B}(n,p)$ & $n \in \N^*, p\in(0,1)$ & $\binom{n}{k}p^k(1-p)^{n-k}$ & $np$ & $np(1-p)$ \\
$\mathcal{G}(p)$ & $p\in(0,1]$ & $(1-p)^{k-1}p$, $k\geq 1$ & $1/p$ & $(1-p)/p^2$ \\
$\mathcal{P}(\lambda)$ & $\lambda>0$ & $e^{-\lambda}\lambda^k/k!$ & $\lambda$ & $\lambda$ \\
$\mathcal{E}(\lambda)$ & $\lambda>0$ & (densité) $\lambda e^{-\lambda t}$ & $1/\lambda$ & $1/\lambda^2$ \\
\bottomrule
\end{tabular}
\vspace{0.5em}
\begin{itemize}
  \item Loi géométrique et exponentielle vérifient la propriété d'absence de mémoire.
  \item $\mathcal{B}(n,p) \to \mathcal{P}(\lambda)$ quand $n \to \infty$, $p \to 0$, $np \to \lambda$.
\end{itemize}
\end{retenirbox}

\end{document}
